\documentclass[11pt]{amsart}
\usepackage{lmodern,amsmath,amssymb,amsthm,mathtools,booktabs,array,longtable}
\usepackage[UTF8,fontset=windows,scheme=plain]{ctex}
\usepackage{mathrsfs}
\usepackage{needspace}
\usepackage[a4paper,margin=28mm]{geometry}
\usepackage{xurl}
% Keep file and JSON path breaks at separators rather than inside words.
\newcommand{\filepath}[1]{\begingroup
  \def\UrlBreaks{\do\/\do\_}\def\UrlBigBreaks{}\Urlmuskip=0mu
  \path{#1}\endgroup}
\usepackage[colorlinks=true,linkcolor=blue,citecolor=blue,urlcolor=blue]{hyperref}
\hypersetup{pdftitle={Theorem (1+1.75) on the Goldbach Conjecture},pdfauthor={Slava Chan}}
\newtheorem{theorem}{定理}[section]
\newtheorem{proposition}[theorem]{命题}
\newtheorem{lemma}[theorem]{引理}
\newtheorem{corollary}[theorem]{推论}
\theoremstyle{definition}
\newtheorem{definition}[theorem]{定义}
\newtheorem{remark}[theorem]{注}
\renewcommand{\proofname}{证明}
\renewcommand{\abstractname}{摘要}
\renewcommand{\refname}{参考文献}
\renewcommand{\tablename}{表}
\renewcommand{\appendixname}{附录}
\DeclareMathOperator{\supp}{supp}
\DeclareMathOperator{\li}{li}
\newcommand{\eps}{\varepsilon}
\newcommand{\one}{\mathbf 1}
\allowdisplaybreaks[1]
\setlength{\emergencystretch}{3em}
\title{哥德巴赫猜想的 $(1+1.75)$ 定理}
\author{Slava Chan}
\address{独立研究者}
\email{ch4nslava@gmail.com}
\date{2026年11月6日}
\subjclass[2020]{11P32, 11N36, 11N13}
\keywords{陈氏定理，双重筛，半素数，等差数列中的素数}
\begin{document}
\begin{abstract}
对固定的 $1\le a\le2$，命题 $(1+a)$ 断言，每个充分大的偶数 $N$ 均可表示为
\[
 N=p+rq,\qquad r\le q^{a-1},
\]
其中 $p,q$ 为素数，$r$ 为 $1$ 或素数。当 $a=1$ 时，这就是针对充分大偶数的二元哥德巴赫猜想；当 $a=2$ 时，这就是陈氏定理。我们无条件地证明，每个充分大的偶数都具有 $a=1.75$ 时的上述表示，从而建立命题 $(1+1.75)$，加强了 Li 和 Liu 此前的 $(1+1.9)$ 定理~\cite{LiLiu}。证明从陈氏表示数的下界中减去被排除半素数的非负上界筛估计，而该下界通过对十一项陈氏权中的配对项作限制并应用双重筛得到。$N$ 的阈值是非有效的，即本证明不给出可计算的显式阈值。
\end{abstract}
\maketitle
\section{引言}
\label{sec:current-statement}

对固定的 $1\le a\le2$，若每个充分大的偶数 $N$ 都有表示
\[
 N=p+rq,\qquad r\le q^{a-1},\qquad
 p,q\text{ 为素数},\qquad r=1\text{ 或为素数},
\]
则称\emph{命题 $(1+a)$} 成立。当 $a=1$ 时，条件迫使 $r=1$，得到针对充分大偶数的二元哥德巴赫猜想；当 $a=2$ 时，将两个因子排序后便得到陈氏定理。中间的参数则要求合数余项的一个素因子比另一个小到何种程度。

我们证明以下定理。

\begin{theorem}
\label{cor:headline}
命题 $(1+1.75)$ 无条件成立。因此，每个充分大的偶数均可表示为 $N=p+rq$，其中 $p,q$ 为素数，$r$ 为 $1$ 或素数，且 $r\le q^{3/4}$。
\end{theorem}

定理~\ref{cor:headline}--\ref{thm:current-main} 及推论~\ref{cor:three-quarters}--\ref{cor:current-explicit} 中关于 $N$ 的阈值均为非有效的。

参数 $1.75$ 可以减小为 $1.74816964$，并进一步减小为 $1.74816865$，两者对应不同的表示数下界，见推论~\ref{cor:current-explicit}。

Li 和 Liu~\cite[定理~1.1]{LiLiu} 证明了命题 $(1+1.9)$，并指出其参数可降至 $1.894$。定理~\ref{cor:headline} 加强了这一素因子限制。两项结论均针对每个充分大的偶数。

在定量论证中，令 $b=a-1$，在 $r\le q^b$ 的限制下计算素数加项。素数余项情形 $r=1$ 始终保留，因为 $1\le q^b$。下述定理对每个固定的 $1/2<b<1$ 给出下界。

全文中，$N>2$ 沿偶数趋于无穷，无下标的对数均为自然对数，$\log_N x=\log x/\log N$，$\Omega(n)$ 表示按重数计算的素因子个数。令
\begin{equation}
 C_2=\prod_{\substack{\ell>2\\\ell\text{ 为素数}}}
 \left(1-\frac1{(\ell-1)^2}\right),\qquad
 C_N=C_2\prod_{\substack{\ell\mid N,\ \ell>2\\\ell\text{ 为素数}}}
 \frac{\ell-1}{\ell-2},\qquad
 H_N=C_N\frac{N}{\log^2 N}.
 \label{eq:current-normalization}
\end{equation}
对任意正偶数 $M$，$C_M$ 表示将同一乘积中的 $N$ 换成 $M$ 所得的量。
我们按不同的素数加项计数，定义
\begin{align}
 D_{1,2}(N)&=\#\{p<N:p\text{ 为素数},\ N-p>1,
                                      \ \Omega(N-p)\le2\},
                                      \label{eq:current-d12}\\
 D_b(N)&=\#\{p<N:p\text{ 为素数},\ N-p=rq,
       \ q\text{ 为素数},\ r=1\text{ 或为素数},\ r\le q^b\},
                                      \label{eq:current-db}\\
 B_b(N)&=\#\{p<N:p\text{ 为素数},\ N-p=rq,
                \ r,q\text{ 为素数},\ q^b<r\le q\}.
                                      \label{eq:current-bb}
\end{align}
有序分解 $r\le q$ 是唯一的。因此，每个合数余项只计一次；平方余项计入 $B_b$，素数余项则计入 $D_b$。

\Needspace{10\baselineskip}
\begin{theorem}
\label{thm:current-main}
令
\begin{equation}
 K=\frac{8452171}{1250000}=6.7617368,
 \qquad c=\frac{49044063}{25000000}=1.96176252.
 \label{eq:current-safe-constants}
\end{equation}
则对每个固定的 $1/2<b<1$，有
\begin{equation}
 \liminf_{\substack{N\to\infty\\2\mid N}}
      \frac{D_b(N)}{H_N}\ge c-K\log(1/b).
 \label{eq:current-general-theorem}
\end{equation}
\end{theorem}

在 $1/2<b<1$ 内，该下界为正当且仅当 $b>\exp(-c/K)$；见第~\ref{sec:current-transfer} 节。

\begin{corollary}
\label{cor:three-quarters}
对每个充分大的偶数 $N$，有
\[
                  D_{3/4}(N)\ge0.0165H_N.
\]
\end{corollary}

\begin{corollary}
\label{cor:current-explicit}
每个充分大的偶数 $N$ 均满足
\begin{align}
 D_{0.74816964}(N)&\ge10^{-5}H_N,
                                  \label{eq:current-main-count}\\
 D_{0.74816865}(N)&\ge10^{-6}H_N,
                                  \label{eq:current-secondary-count}\\
 D_{0.9}(N)&\ge1.24889H_N.
                                  \label{eq:current-common-exponent}
\end{align}
\end{corollary}

前两个不等式分别建立命题 $(1+1.74816964)$ 与 $(1+1.74816865)$，在归一化~\eqref{eq:current-normalization} 下，对应的计数系数分别为 $10^{-5}$ 与 $10^{-6}$。$b=0.9$ 的估计便于在相同的素因子限制下比较。

\begin{proof}[定理~\ref{cor:headline} 的证明]
由推论~\ref{cor:three-quarters}，对每个充分大的偶数 $N$，均有 $D_{3/4}(N)\ge0.0165H_N>0$，故命题 $(1+1.75)$ 成立。
\end{proof}

\subsection{相关工作与方法}
在归一化 \eqref{eq:current-normalization} 下，Li 和 Liu~\cite[定理~1.1]{LiLiu} 在 $b=0.9$ 时的计数系数为 $0.0004$，其所述改进对应 $b=0.894$。在 $b=0.9$ 时，本文的 $D_b(N)$ 与 Li 和 Liu 的 $D_{1,1.9}(N)$ 是同一个计数函数：两者都包含素数余项情形 $r=1$，对每个可容许的素数加项只计一次，并使用相同的奇异级数归一化。因此，这两个系数指同一计数函数。本文的系数由 $c-K\log(1/b)$ 给出，其中 $c$ 来自陈氏总计数的下界构造，$K$ 来自非负上界控制函数。推论~\ref{cor:current-explicit} 中较小的指数仍保留显式正计数。他们对下界项及带标记筛项的讨论~\cite[第~3 节及第~8.2 节]{LiLiu} 也指出了本文的一项关键限制：更高的上界筛分布水平必须与每个使用处的系数类别匹配。

双重筛在上、下界筛余项之间传递信息。我们采用 Wu~\cite{Wu} 的表述，并对两个余项函数作有限次非负转移。Li~\cite[定理~1.1]{RunboLi} 在同一归一化下给出陈氏总计数的系数 $1.9728$。本文使用的系数 $c=1.96176252$，是在因子受限论证所需的变动剩余类系数类别和非负带标记族内证明的。每次下界转移均保留实际剩余筛级及全部前缀条件。

分布估计必须对随 $N$ 变化的剩余类一致。Maynard 的三重良可分解估计~\cite[定义~2 及定理~1.1]{MaynardII} 中的常数依赖于固定剩余类。Lichtman 的大剩余类结论~\cite[定理~1.7]{Lichtman} 在 $\vartheta=7/32$ 时具有水平 $(5-\vartheta)/8-\epsilon=153/256-\epsilon$。这里的剩余类与完整局部乘积同阶，因此保留 \eqref{eq:current-retained-phase} 中的互反相位；在此范围内，以 $1$ 替换该相位所产生的误差未必趋于零。附录~\ref{app:medium}、\ref{app:classical} 结合附录~\ref{app:coefficients} 的分解，为本文所用的有限上、下界系数族在支撑指数 $0.599$ 处建立受限分布估计。每次应用均保留选定分解、完整与互补平方条件及带标记族的非负性。

证明从全部陈氏表示数中减去被排除的半素数。上界部分构造非负除数控制函数，计算其密度，并针对同一系数证明分布估计。下界部分对十一项陈氏权的二素数项和四素数项作配对限制，以单射保留逐点权不等式，再利用受限纤维上 Buchstab 函数的下界改善系数。

非负上界控制函数及其系数实现在本文中构造。下界部分采用 Wu 的双重筛框架，并加入本文建立的配对删除及系数实现。这些构造与有限转移、定向包围共同确定常数 $K,c$。这些量的数学定义及其证书对应见附录~\ref{app:numeric}。

\section{归约与证明概要}
\label{sec:outline}

\paragraph{证明概要与核验。}
首读时可暂以引理~\ref{lem:numerical-endpoints} 为有限数值输入，依次理解上界构造、十一项下界构造与配对删除，再读第~\ref{sec:current-transfer} 节的最终相减。完整核验则还需回到各附录，检查分布估计、系数实现、有限转移与定向包围，并与附录~\ref{app:numeric} 所述数值包逐项对应。

由唯一分解，有精确恒等式
\begin{equation}
                 D_b(N)=D_{1,2}(N)-B_b(N).
                 \label{eq:reduction}
\end{equation}
因此只需证明
\begin{equation}
 \liminf_{\substack{N\to\infty\\2\mid N}}
       \frac{D_{1,2}(N)}{H_N}\ge c,
 \qquad
 \limsup_{\substack{N\to\infty\\2\mid N}}
       \frac{B_b(N)}{H_N}\le K\log(1/b).
 \label{eq:two-estimates}
\end{equation}
因子 $\log(1/b)$ 来自被排除因子的区域。令 $\beta=b/(1+b)$，不满足限制的有序半素数在对数坐标中的极限区域为 $\beta<t\le1/2$，且
\[
 \int_\beta^{1/2}\frac{dt}{t(1-t)}=\log(1/b).
\]
要将这一质量转化为补数为素数的计数，我们使用非负上界筛，其实际系数的分布估计及密度均已建立。

第~\ref{sec:current-distribution} 节建立实际所用上、下界系数的分布估计；第~\ref{sec:current-upper} 节构造非负上界筛并得到 $K$；第~\ref{sec:current-chen-weight}、\ref{sec:current-lower} 节给出下界构造，其中换位后的代价使用前面已证明的上界；第~\ref{sec:current-numeric} 节列出有限数值不等式及误差预留量；最后，第~\ref{sec:current-transfer} 节证明 \eqref{eq:two-estimates} 的第二个估计并完成主定理。特别地，$K$ 的构造不依赖于下界双重筛表，也不依赖于由这些表得到的陈氏计数。

\paragraph{参数选择。}
在每个固定精度下，有限规则与带标记查询的集合、迭代次数、确定固定紧参数域的截断参数、逼近容差和支撑预留量都在 $N$ 趋于无穷之前选定。素数区间、权和算术系数仍可通过这些固定规则依赖于 $N$。共同选择与误差传播见第~\ref{sec:lower-limits} 节及 \eqref{coef:parameter-order}；第~\ref{sec:current-numeric} 节的严格标量余量允许为最终计数选定一个固定精度。

\paragraph{简短符号索引。}
以下符号在正文与附录中采用一致含义。
\begin{center}
\small
\begin{tabular}{@{}>{\raggedright\arraybackslash}p{.23\textwidth}>{\raggedright\arraybackslash}p{.49\textwidth}>{\raggedright\arraybackslash}p{.22\textwidth}@{}}
\toprule 符号 & 含义 & 位置\\
\midrule
$E_{\rm up}$；$R_{\rm low}$ & 上界与下界的标量端点 & \eqref{eq:current-E}；\eqref{eq:current-R}\\[2pt]
$\mathscr E_0,\mathscr E_1,\mathscr E$ & 配对矩形族及其并 & 第~\ref{sec:current-chen-weight} 节\\[2pt]
$a_{\rm s},h_{\rm s},d_{\rm s}$ & 固定筛截点 & 第~\ref{sec:current-chen-weight} 节\\[2pt]
$H(s),h(s)$；$H_m,h_m$ & 上、下筛改进量；第 $m$ 代包络 & 第~\ref{sec:current-lower} 节\\[2pt]
$B$；$\mathscr B$ & 算术素数前缀；几何矩形或盒 & 第~\ref{sec:current-distribution}、\ref{sec:current-lower} 节\\[2pt]
$H_N$；$B_{\rm low}$ & 计数归一化；每项下界误差预留 & \eqref{eq:current-normalization}；\eqref{eq:current-B}\\[2pt]
$A_+,A_-$ & 归一化线性筛分子 $sF/(2e^\gamma),sf/(2e^\gamma)$ & 第~\ref{sec:current-lower} 节\\[2pt]
$\Omega$；$\omega_{\rm env}$ & 按重数计的素因子个数；Buchstab 上包络 & 第~\ref{sec:current-statement} 节；\eqref{sv:omega-envelope}\\[2pt]
$\rho_N^{\rm up},\rho_N^{\rm low}$ & 实际上、下筛除数系数；$\rho_N$ 泛指任一可用系数 & 第~\ref{sec:current-upper}、\ref{sec:current-lower} 节\\[2pt]
\bottomrule
\end{tabular}
\end{center}

我们采用标准线性筛函数 $F,f$，其归一化为
\[
 sF(s)=2e^\gamma\quad(1\le s\le3),\qquad
 sf(s)=0\quad(1\le s\le2),
\]
在后续各自适用的区间上，分别由 $(sF(s))'=f(s-1)$ 和 $(sf(s))'=F(s-1)$ 定义。Buchstab 函数满足
\[
 \omega(s)=1/s\quad(1\le s\le2),\qquad
 (s\omega(s))'=\omega(s-1)\quad(s>2).
\]
以 $P^-(n)$ 表示 $n\ge2$ 的最小素因子，并约定 $P^-(1)=+\infty$。这些函数所用的有限区间数据见附录~\ref{app:numeric}。

\subsection{一个非负筛系数定理}
上界估计使用下面的单一系数定理，其中三个结论均针对同一组除数系数。这里的“非负”（positive）指所得除数和 $R_N(n)$ 是非负上界控制函数；各个除数系数 $\rho_N^{\rm up}(q)$ 可以带符号。

\begin{theorem}
\label{thm:positive-coefficient}
固定 $a_*>0$、带标签素因子个数的一个上界，以及有限组光滑权 $W_N(p_1,\ldots,p_k)$。这些权的支撑满足
\[
 p_i\ge c_*N^{a_*},\qquad p_1\cdots p_k\le C_*N,
\]
其中 $c_*,C_*>0$ 为固定常数；权的大小一致有界，其混合 Euler 导数满足固定的多对数界。除这些光滑几何条件外，不再附加其他算术筛选条件。令 $\rho_0=19/96$、$\theta_*=599/1000$，并记
\[
 P_N(y)=\prod_{\substack{p<y\\p\nmid N}}p,\qquad
 X_W^{(q)}=\sum_{(p_1\cdots p_k,q)=1}W_N(p_1,\ldots,p_k).
\]
则存在实系数 $\rho_N^{\rm up}(q)$ 以及不依赖于 $N$ 的常数 $M$，其支撑为满足 $q\le N^{\theta_*}$、$(q,N)=1$ 的无平方因子整数，且 $|\rho_N^{\rm up}(q)|\le M$，使得
\begin{equation}
 R_N(n):=\sum_{q\mid n}\rho_N^{\rm up}(q)
   \ge\one_{(n,P_N(N^{\rho_0}))=1}\ge0
                           \qquad(n\in\mathbb Z),\label{eq:core-positivity}
\end{equation}
\begin{equation}
 \sum_q\frac{\rho_N^{\rm up}(q)}{\varphi(q)}
   \le(K+o(1))\frac{C_N}{\log N},\label{eq:core-density}
\end{equation}
\begin{multline}
 \sum_q\rho_N^{\rm up}(q)\left\{
 \sum_{p_1\cdots p_k\equiv N\ (q)}W_N(p_1,\ldots,p_k)
       -\frac{X_W^{(q)}}{\varphi(q)}\right\}\\
   =O_{A,W}(N(\log N)^{-A})\qquad(A>0).
                           \label{eq:core-distribution}
\end{multline}
该有限和也定义了 $R_N(0)$。有限前缀决策规则、逼近参数和支撑余量均在 $N$ 沿所有偶数趋于无穷之前固定。常数可以依赖于上述固定的有限组权及其导数界，但对 $N$ 一致。
\end{theorem}

这组系数在第~\ref{sec:current-upper} 节构造。我们先在第~\ref{sec:current-distribution} 节针对该构造所需的精确系数表示证明分布结论。

\section{筛系数的分布估计}
\label{sec:current-distribution}

\begin{definition}
有限前缀决策规则在有序素数前缀的坐标 $t_i=\log_Np_i$ 中指定有限多个区域，并为每个区域指定一个筛级。每个前缀处的决策在细分其子节点之前固定。当选定分支在 \eqref{coef:full-squares}--\eqref{coef:selected-squares} 中的完整平方条件和选定补集条件都具有所需的严格余量时，称该分支可用；继承分支则保留其父节点的决策。带标记族是具有共同总筛级 $Q$ 的前缀 $B$ 的非负有限加权族，其剩余筛级为 $Q/B$，并保留原有的素数上限。带标记查询是对这样的一个族作筛计数。下文所谓可用系数，均指由这些规则及有限次带标记转移得到、且具有附录~\ref{app:coefficients} 所给三因子表示的系数。
\end{definition}

\subsection{变动剩余类的局部估计}

分布问题以完整乘积为局部尺度。用 $U,V$ 表示卷积长度，记 $x=UV$，并以 $\mathfrak a$ 表示剩余类；$N$ 仍表示全局的目标偶数。附录~\ref{app:medium} 证明以下估计。

\begin{proposition}[受限中等长度卷积]
\label{prop:current-medium}
固定 $\vartheta=7/32$、$0<\sigma<1/100$、$c_0>0$、除数函数的阶数 $B$ 以及 $0\le\kappa\le\sigma/100$。设 $U,V,Q_1,Q_2,Q_3\ge1$，并假设
\begin{align}
 c_0x&\le |\mathfrak a|\le x^{1+\kappa},&
 Q_1&<Ux^{-\sigma}, \notag\\
 U^2Q_2Q_3^2&<x^{1-8\sigma},&
 |\mathfrak a|^{\vartheta}UQ_1Q_2^5Q_3^2&<x^{2-15\sigma},
       \label{eq:current-medium-ranges}\\
 U^{1+2\vartheta}Q_1Q_2^{5-2\vartheta}Q_3^{2+2\vartheta}
      &<x^{2-15\sigma}.&&\notag
\end{align}
设 $\alpha$ 的绝对值不超过一，支撑在 $[U,2U]$ 中的某个区间上，并满足 $P^-(n)>x^{1/(\log\log x)^3}$。假定它满足带辅助互素条件的 Siegel--Walfisz 估计：对任意 $A>0$、任意 $k,q\ge1$ 以及 $(b_0,q)=1$，有
\begin{equation}
 \left|\sum_{\substack{n\equiv b_0\pmod q\\(n,k)=1}}\alpha_n
 -\frac1{\varphi(q)}\sum_{(n,kq)=1}\alpha_n\right|
 \le C_AU\tau(k)^{C_0}(\log x)^{-A},
 \label{eq:current-aux-sw}
\end{equation}
其中 $C_0$ 固定，族中的常数一致。设 $|\beta_m|\le\tau(m)^B$、$m\in[V,2V]$，且每个 $\gamma_i$ 都是在 $q_i\asymp Q_i$ 上独立给定、绝对值不超过一的序列，支撑满足 $(q_i,\mathfrak a)=1$。令
\[
 \Delta_{\ell}=\sum_{n,m}\alpha_n\beta_m
 \left({\one}_{nm\equiv\mathfrak a\pmod \ell}
       -\frac{{\one}_{(nm,\ell)=1}}{\varphi(\ell)}\right),
\]
则附录~\ref{app:classical} 中列明的经典输入给出
\begin{equation}
 \sum_{q_1,q_2,q_3}\gamma_1(q_1)\gamma_2(q_2)\gamma_3(q_3)
                    \Delta_{q_1q_2q_3}
       =O_A\bigl(x(\log x)^{-A}\bigr).
 \label{eq:current-medium-conclusion}
\end{equation}
常数对所有满足上述条件的剩余类和系数族一致。端点处的固定倍数由严格的幂次余量吸收。
\end{proposition}

\paragraph{局部估计的应用。}
先固定一个有限可用系数族。附录~\ref{app:coefficients} 利用其完整平方和选定补集平方条件，在完整原子盒上构造三个彼此分离的模数因子。附录~\ref{app:prime-products} 核验命题~\ref{prop:current-medium} 的其余条件，并将其应用于有限素数乘积族。各步骤对应附录~\ref{app:coefficients} 的\hyperlink{coef:medium-checks}{条件核验表}。低阶部分及不含中等长度变量的分支分别采用其相应估计。筛构造共同使用的输入是下一小节的素数乘积分布估计。

命题~\ref{prop:current-medium} 中三个模数系数彼此分离，估计允许系数带符号。应用时，中等长度变量是在单独坐标片上的粗糙 $1$ 序列或 M\"obius 序列，并可带有允许的 Mellin 扭曲。附录~\ref{app:prime-products} 为这些序列证明辅助 Siegel--Walfisz 条件~\eqref{eq:current-aux-sw}。

当剩余类的大小与 $x$ 同阶时，互反变换留下一个振荡相位。附录~\ref{app:medium} 保留该相位，并用 \eqref{eq:current-medium-ranges} 中的严格幂次余量吸收其变差损失。包括 Fourier 归一化在内的谱估计陈述于附录~\ref{app:classical}。

\subsection{从局部估计到素数乘积}

固定一个关于 $k$ 个带标签原始素数的光滑权 $W_N$，各素数至少为 $cN^{a_*}$，其中 $a_*>0$ 固定。令
\[
 X_W^{(q)}=\sum_{(p_1\cdots p_k,q)=1}W_N(p_1,\ldots,p_k),
 \qquad X_W=\sum W_N(p_1,\ldots,p_k).
\]
乘积为 $O(N)$。允许的几何族具有固定的多对数混合 Euler 导数界，且没有额外算术筛选条件。全文使用的结论是
\begin{equation}
 \sum_q\rho_N(q)
 \left\{\sum_{p_1\cdots p_k\equiv N\pmod q}W_N(p_1,\ldots,p_k)
            -\frac{X_W^{(q)}}{\varphi(q)}\right\}
       =O_{A,W}\bigl(N(\log N)^{-A}\bigr).
 \label{eq:current-PP}
\end{equation}
两项中出现的是同一带符号系数 $\rho_N$。其可用类别包括实际的上界树、实际的下界树、有限 $H/h$ 构造的带标记查询及普通后继，以及应用上界系数的非负换位族。

附录~\ref{app:coefficients} 的系数证明首先固定一个有限的前缀决策规则。实际筛级 $N^\lambda$ 决定密度及剩余筛级 $Q/B$。证明分解时可以采用略大的辅助支撑覆盖 $N^{\ell_{\rm aux}}$，但这不改变 $\lambda$ 或整数上的系数。对有序素数前缀，记 $t_i=\log_Np_i$、$S_j=\sum_{i=1}^jt_i$，并设 $I$ 为选入某个模数因子的因子位置子集。完整平方条件
\[
 S_{j-1}+2t_j\le\ell_{\rm aux}
\]
以及每个选定补集平方条件
\[
 \sum_{i\in I}t_i+
       \sum_{\substack{i<j\\i\notin I}}t_i+2t_j
       \le\ell_{\rm aux}\quad(j\notin I)
\]
均予保留。当全部选定位置都已出现后，第二个检验成为完整平方检验。因此，在有限次普通查询中，后继封闭性和实际素数上限始终保留。

估计偏差之前存在一个精确的张量恒等式。对具有重数 $m_I$ 的共同素数标签集合，将各因子位置分成三组，并将每组卷积除以 $M_j=\prod_I m_{I,j}!$，所得 $\gamma_j$ 的绝对值不超过一。记 $M_{\rm all}=\prod_I m_I!$，则完整的归一化卷积为
\[
 \frac{M_1M_2M_3}{M_{\rm all}}
                       (\gamma_1*\gamma_2*\gamma_3).
\]
带标记重数保留在外部标量中。小素数 beta 因子通过唯一的小素数部分与粗糙部分的分解接入。分组在整个原子盒上预先固定，不依赖于随后出现的算术偏差。全局无平方因子条件仍保留在实际系数中；仅为解析比较而去掉该条件时，误差为固定幂次节省
$O((N+N^{\theta_*})N^{-\eta_0+o(1)})$，其中 $N^{\eta_0}$ 是小素数预筛截点，$\theta_*<3/5$ 是有限系数构造的支撑指数。

附录~\ref{app:prime-products} 先光滑截去 $2N^{1-\kappa_0}$ 以下的完整乘积，再作精确的光滑二进分割。局部尺度始终是完整乘积，满足
\[
 cN^{1-\kappa_0}\le x\le CN,\qquad
 c_0x\le N\le x^{1+\kappa}.
\]
而不是剩余商。在尺度为 $R$ 的原始分组上，由 $2R$ 以下的恒等式 $\Lambda*(\delta-1*\mu_Z)^{*3}=0$，取 $Z=(2R)^{1/3}$，得到精确的三阶 Heath--Brown 恒等式
\begin{equation}
 \Lambda=\sum_{j=1}^3(-1)^{j-1}\binom3j
                 \mu_Z^{*j}*1^{*(j-1)}*\log.
 \label{eq:current-HB}
\end{equation}
保留原来的光滑坐标片，不引入分组之间的尖锐耦合边界。几何轮廓的 Fourier 分离所需代价用其可积 Fourier 范数控制。实际的粗糙因子满足辅助 SW 估计。中等长度分支使用命题~\ref{prop:current-medium}；低秩分支使用低于二分之一水平的 Bombieri--Vinogradov 估计。单长因子与双长因子的情形采用附录~\ref{app:brun} 的直接 Brun 论证，保留精确互素均值以及两条 Poisson 零坐标轴上的抵消。特别地，不把可能为负的 Brun 下界多项式当作非负量相乘。

先固定矩的阶数 $K_{\rm mom}$，再缩小参数范围，使
\[
 0<\kappa\le\min\left(\kappa_*,\frac{\sigma}{200},
                           \frac1{2K_{\rm mom}}\right),\qquad
 0<\kappa_0<\frac{\kappa}{4(1+\kappa)}.
\]
其中 $\kappa_*>0$ 由附录~\ref{app:coefficients} 中的严格支撑余量确定。完整局部支撑的指数为 $\Theta_{\rm loc}<3/5$；长分支参数 $\zeta$ 选为满足 $20\zeta+\Theta_{\rm loc}<1$ 和 $40\zeta+3\Theta_{\rm loc}/2<1$。所有有限原子的标量质量以及固定对数数量的坐标片，都由更强的局部对数精度吸收。这就对可用族以及所有充分大的目标偶数证明了 \eqref{eq:current-PP}。

\section{非负上界筛及其密度}
\label{sec:current-upper}

上界构造使用
\[
 \rho_0=\frac{19}{96},\quad \mu_0=\frac{25}{128},\quad
 \Theta=\frac{299}{500},\quad \Lambda_{\rm cap}=\frac{153}{256},\quad
 v(t)=\min\left(\frac{1+t}{2},\frac{57}{64}-\frac{3t}{2}\right).
\]
筛级上限的变更依实际前缀作出，所有保留的祖先决策均不改变。第四、第五和第六次变更依次替换前一规则，其原始密度差构成望远镜求和。其后的八阶变更作用于 $30+15+13+29=87$ 个指定区域。同阶段各区域的内部互不相交；不同阶段均相对于各自父节点的上限取差。上限公式及分离证明见附录~\ref{app:upper}。

\subsection{退出恒等式与实际剩余密度}

对有限的允许素数集合及 $0\le \xi(p)\le1$，首次退出均为奇数阶的上界树满足精确恒等式
\begin{equation}
 \sum_d\mu(d){\one}_{d\ {\rm retained}}\xi(d)
 =\prod_p(1-\xi(p))
   +\sum_{e\ {\rm odd\ first\ exit}}\xi(e)
                      \prod_{p<p_{\rm last}(e)}(1-\xi(p)).
 \label{eq:current-upper-exit}
\end{equation}
其中 $\xi$ 按乘法延拓到无平方因子的乘积，且 $\xi(1)=1$。
按第一个被排除的前缀划分被排除的除数，其余较小素数的子集给出式中的 Euler 乘积。取 $\xi(p)={\one}_{p\mid n}$，即对每个整数 $n$ 得到粗糙性示性函数的上界，包括零和负整数。

计算原始密度时，取 $g_N(p)={\one}_{p>2,\,p\nmid N}/(p-1)$ 及 $V_N(y)=\prod_{2<p<y,\,p\nmid N}(1-1/(p-1))$。对固定阶数，$g_N$ 的归一化退出量介于零和 $g_0(p)={\one}_{p>2}/(p-1)$ 的相应退出量之间，因为每个仍存在的前缀会带上乘子
\[
 \prod_{\substack{p_{\rm last}\le p<y\\p\mid N}}(1-g_0(p))\le1.
\]
在正参数紧集上，先前的平方检验使 $j$ 阶退出的末素数至少为 $N^{\beta_{\min}/(3\cdot2^{j-1})}$。因此，去掉整除目标数的素数，对每个固定阶积分造成的变化为对 $N$ 一致的 $o(1)$。由非负退出级数的收敛、同一原始 Rosser 权的单侧界，再以有限阶上下夹逼，得到
\begin{equation}
 \frac{V_N^+(R,y)}{V_N(y)}=F(s)+o(1)\quad(s>1),\qquad
 \frac{V_N^-(R,y)}{V_N(y)}=f(s)+o(1)\quad(s\ge2),
 \quad s=\frac{\log R}{\log y}.
 \label{eq:current-raw-density}
\end{equation}
当 $1<s\le3$ 时，上界恒等式也可将首个立方条件移至 $R^{1/3}$ 而直接得到。由此得到除数系数本身的密度。普通退出级数所用的输入及其归一化见附录~\ref{app:upper}。

Mertens 定理对目标数一致地给出
\begin{equation}
 V_N(N^u)=\frac{2e^{-\gamma}C_N}{u\log N}(1+o(1)).
 \label{eq:current-mertens}
\end{equation}
例如，与 $2C_N\prod_{p<z}(1-1/p)$ 的精确比值为
\[
 \prod_{p\ge z}\left(1-\frac1{(p-1)^2}\right)^{-1}
 \prod_{\substack{p\mid N\\p\ge z}}\left(1-\frac1{p-1}\right),
\]
当 $z=N^u$、$u>0$ 固定时，其对数一致趋于零。令 $\overline F=F/(2e^\gamma)$ 和 $\overline f=f/(2e^\gamma)$，这一归一化在密度泛函中产生 $4\overline F(s)/u$ 与 $4\overline f(s)/u$。

对一个选定的八阶单元，若 $S_8=t_1+\cdots+t_8$ 且父节点的上限不超过 $R_6$，则其上限 $M$ 满足
\[
 1<\frac{\lambda_{6,\rm parent}-S_8}{t_8}
  \le\frac{R_6-S_8}{t_8}<\frac{M-S_8}{t_8}<3.
\]
故同一个剩余密度恒等式给出如下节省量下界：
\begin{equation}
 \frac{4(M-R_6)}{(R_6-S_8)(M-S_8)\prod_{i=1}^8t_i}.
 \label{eq:current-eighth-saving}
\end{equation}
采用较大的父上限 $R_6$ 只会使估计更保守。此密度不含陈氏权的四分之一因子，也不含排序阶乘。选定区域满足 $t_8>53/13824$，且所有完整与选定补集平方条件仍然有效。

\subsection{有限实现与密度误差}

在 $N$ 增大之前，先用有限的前缀分割逼近理想规则。末素数很小时的尾部趋于零；真正的仿射边界质量趋于零，但开放的继承区域或等式区域予以保留。对偶数粗糙阶使用端点外检验，对奇数阶使用内检验，并保持祖先关系。这样得到高于原始上界系数的 $a_N(d)$。对正的 $a_N(d)$ 接入小素数 beta 上界多项式，对负的 $a_N(d)$ 接入 beta 下界多项式。若 $B^-\le I_w\le B^+$，则逐项有 $a_N(d)B_{\rm sign}(n)\ge a_N(d)I_w(n)$。
所得实际系数 $\rho_N^{\rm up}$ 具有唯一的小素数与粗糙部分分解，其支撑为 $q\le N^{\theta_*}$ 内的无平方因子整数；对固定的 $M$ 有一致界 $|\rho_N^{\rm up}(q)|\le M$，并满足 $(q,N)=1$ 及
\begin{equation}
 R_N(n):=\sum_{q\mid n}\rho_N^{\rm up}(q)
   \ge {\one}_{(n,P_N(N^{\rho_0}))=1}\ge0,
 \qquad P_N(y)=\prod_{\substack{p<y\\p\nmid N}}p.
 \label{eq:current-physical-upper}
\end{equation}
这个有限和也定义了 $R_N(0)$。

将 \eqref{eq:current-upper-exit} 限制在预筛以上的素数上，得到上界投影不等式
\[
 V_N(w)\sum_{d\ {\rm rough}}\lambda_N(d)g_N(d)
       \le\sum_{d\ {\rm full}}\lambda_N(d)g_N(d).
\]
结合 beta 的绝对误差及盒逼近误差，即可从上方控制同一 $\rho_N^{\rm up}$ 的密度。支撑膨胀可以控制在 $1/1000$ 以下，因此在完整尺度的小幅调整之前可取 $\theta_*=599/1000$。第~\ref{sec:current-distribution} 节的辅助支撑覆盖不改变这一密度。

设 $E_{\rm up}$ 为原始上界密度的端点，$u$ 为有限逼近与系数实现各自的误差预留量，其精确有理数值见 \eqref{eq:current-E}--\eqref{eq:current-u}。附录~\ref{app:numeric} 的定向估计给出
\begin{equation}
                     K-E_{\rm up}-2u>9.6522\cdot10^{-8}>0.
                     \label{eq:current-upper-slack}
\end{equation}
因此，\eqref{eq:current-physical-upper} 中的实际系数满足
\begin{equation}
 S_N:=\sum_q\frac{\rho_N^{\rm up}(q)}{\varphi(q)}
       \le(K+o(1))\frac{C_N}{\log N},
 \label{eq:current-upper-density}
\end{equation}
而 \eqref{eq:current-PP} 恰好也适用于这同一组系数。

\begin{proof}[定理~\ref{thm:positive-coefficient} 的证明]
取上述有限构造所得的系数。式 \eqref{eq:current-physical-upper} 给出其有界的无平方因子支撑及 \eqref{eq:core-positivity}，式 \eqref{eq:current-upper-density} 给出密度上界。附录~\ref{app:coefficients} 中的分解与参数选择结合附录~\ref{app:prime-products} 的素数乘积估计，为同一系数及固定有限组中的每个权建立 \eqref{eq:current-PP}。在 $N$ 趋于无穷之前选择共同细分，再取有限多个常数与阈值的最大值，即得所述一致性。
\end{proof}

\begin{proposition}[非负乘积族的推论]
\label{prop:current-positive-product}
固定一个阶数有界、权一致有界的非负有限带标签素数乘积族。设其每个原始素数都至少为 $N^{a_*}$，其中 $a_*>0$ 固定，乘积为 $O(N)$，且仅受可用的固定几何条件限制。假定其未筛加权质量 $X_W$ 满足
\[
 \limsup_{\substack{N\to\infty\\2\mid N}}
             \frac{\log N}{N}X_W\le J<\infty.
\]
用 $W_N(m)$ 表示在乘积 $m$ 上诱导的权，则
\[
 \limsup_{\substack{N\to\infty\\2\mid N}}\frac1{H_N}
       \sum_{\substack{m<N\\N-m\text{ 为素数}}}W_N(m)\le KJ.
\]
即使 $J=0$，误差估计也仍是绝对的。
\end{proposition}

\begin{proof}
对素数 $\ell$，由无平方因子支撑及有界性，得
\[
 \sum_{q:\ell\mid q}\frac{|\rho_N^{\rm up}(q)|}{\varphi(q)}
       \ll\frac{\log N}{\ell-1}.
\]
对固定个数的原始素数位置取并集上界，得到
\begin{equation}
 \left|\sum_q\frac{\rho_N^{\rm up}(q)}{\varphi(q)}
                         (X_W^{(q)}-X_W)\right|
       \ll X_WN^{-a_*}\log N=o(H_N).
 \label{eq:current-coprime-correction}
\end{equation}
这里不引入因子 $N^{\theta_*}$。较小的素数余项代价为 $N^{\rho_0+o(1)}$。完整乘积的光滑截尾在对全部模数求和后代价为 $N^{1-\kappa_0+o(1)}$；其乘积严格小于 $N$，所以不会出现 $\tau(0)$。先用非负控制，再用分布估计，得
\[
 \sum_{m<N,\,N-m\ {\rm prime}}W_N(m)
 \le\sum_m W_N(m)R_N(N-m)+o(H_N)
 \le X_WS_N+o(H_N).
\]
由于对每个整数均有 $R_N(n)\ge0$，外逼近中出现 $m\ge N$ 也无妨。随后应用 \eqref{eq:current-upper-density}。固定几何边界通过非负光滑外逼近处理：先令 $N$ 趋于无穷，再令逼近误差趋于零。在每个精度下，逐点不等式、密度及分布公式都使用同一 $\rho_N^{\rm up}$；不同精度下不要求系数相同。
\end{proof}

\section{完整陈氏权与配对删除}
\label{sec:current-chen-weight}

固定筛截点记为 $a_{\rm s},h_{\rm s},d_{\rm s}$；命题 $(1+a)$ 中的参数 $a$ 仍保留引言中的含义。令
\[
 {a_{\rm s}}=\frac{417}{5000},\qquad {h_{\rm s}}=\frac{19}{125},\qquad
 {d_{\rm s}}=\frac{59}{250},\qquad \mu_0=\frac{25}{128}.
\]
对 $n\le N$ 定义 $Q_d(y;n)=\one_{d\mid n}\one_{P^-(n/d)\ge y}$。以下所有素数前缀的分量均互不相同。省略自变量 $n$，十一项为
\begingroup\allowdisplaybreaks[4]
\begin{align*}
 Y_1&=Q_1(N^{a_{\rm s}}),\\
 Y_2&=Q_1(N^{h_{\rm s}}),\\
 Y_3&=\sum_{N^{a_{\rm s}}\le p<N^{1/3}}Q_p(N^{a_{\rm s}}),\\
 Y_4&=\sum_{N^{a_{\rm s}}\le p<N^{d_{\rm s}}}Q_p(N^{a_{\rm s}}),\\
 Y_5&=\sum_{N^{a_{\rm s}}\le p_1<p_2<N^{h_{\rm s}}}Q_{p_1p_2}(N^{a_{\rm s}}),\\
 Y_6&=\sum_{N^{a_{\rm s}}\le p_1<N^{h_{\rm s}}\le p_2<N^{d_{\rm s}}}Q_{p_1p_2}(N^{a_{\rm s}}),\\
 Y_7&=\sum_{N^{d_{\rm s}}\le p_1<p_2<(N/p_1)^{1/2}}Q_{p_1p_2}(p_2),\\
 Y_8&=\sum_{N^{a_{\rm s}}\le p_1<N^{1/3}\le p_2<(N/p_1)^{1/2}}
                                          Q_{p_1p_2}(p_2),\\
 Y_9&=\sum_{N^{h_{\rm s}}\le p_1<N^{d_{\rm s}}\le p_2<(N/p_1)^{1/2}}
                   Q_{p_1p_2}\bigl((N/(p_1p_2))^{1/2}\bigr),\\
 Y_{10}&=\sum_{N^{a_{\rm s}}\le p_1<p_2<p_3<p_4<N^{h_{\rm s}}}Q_{p_1p_2p_3p_4}(p_2),\\
 Y_{11}&=\sum_{N^{a_{\rm s}}\le p_1<p_2<p_3<N^{h_{\rm s}}\le p_4<N^{d_{\rm s}}}
                                          Q_{p_1p_2p_3p_4}(p_2).
\end{align*}
\endgroup
特别地，$Y_{11}$ 保留完整的最后一个取值区间。设
\begin{equation}
 W=3Y_1+Y_2-Y_3-Y_4+Y_5+Y_6-2Y_7-Y_8-Y_9-Y_{10}-Y_{11}.
 \label{eq:current-full-weight}
\end{equation}

其中 $3Y_1+Y_2$ 给出两个截点处的粗糙数计数基础项；单前缀减项 $Y_3,Y_4$ 去掉多计部分，$Y_5,Y_6$ 则补回选定的二素数前缀贡献。$2Y_7,Y_8,Y_9$ 扣除其余含至少三个素因子的情形；四前缀项 $Y_{10},Y_{11}$ 补偿多个小素因子出现时正的配对计数的增长。各类项合起来给出下面的逐点不等式。后续配对删除用单射关联 $Y_6$ 与 $Y_{11}$，使改动后的权仍满足该不等式。

\begin{lemma}[完整权的逐点不等式]
\label{lem:current-full-weight}
对无平方因子的 $n\le N$，有
\[
                  W(n)\le4{\one}_{\Omega(n)\le2}.
\]
\end{lemma}

\begin{proof}
若 $n$ 不是 $N^{a_{\rm s}}$-粗糙的，则十一项全为零。唯一需要稍加说明的是 $Y_9$，其商的筛除截点超过 $N^{(1-{d_{\rm s}})/4}>N^{a_{\rm s}}$。因此可设 $n$ 无平方因子且为 $N^{a_{\rm s}}$-粗糙的。用 $r,s,t$ 分别计数它在 $N^{h_{\rm s}},N^{d_{\rm s}},N^{1/3}$ 以下的素因子，则前六个计数为
\[
 (Y_1,Y_2,Y_3,Y_4,Y_5,Y_6)
      =\left(1,{\one}_{r=0},t,s,\binom r2,r(s-r)\right).
\]
若 $r\ge2$，在第二个选定素数处对商作筛，迫使 $Y_{10}$、$Y_{11}$ 中前两个素数为 $n$ 的最小两个素因子。因此
\[
 Y_{10}=\binom{r-2}{2},\qquad Y_{11}=(r-2)(s-r),
\]
其中二项式在 $r=2,3$ 时为零。代入前六个加权项再减去 $Y_{10}+Y_{11}$，得到 $s-t\le0$；其余减项只会加强结论。

当 $k=\Omega(n)\ge3$ 时，设 $q_1<q_2<\cdots$ 为 $n$ 的素因子。于是 $q_1q_2^2<n\le N$，从而 $t\ge1$、$Y_7={\one}_{s=0}$ 及 $Y_8={\one}_{t=1}$。当 $r\le1$ 时，$W$ 的其余上界为
\[
\begin{array}{c|c}
 \text{情形}&\text{上界}\\ \hline
 r=1&2-t-{\one}_{t=1}\\
 r=0,\ s\ge2&4-s-t\\
 r=s=0&2-t-{\one}_{t=1}\\
 r=0,\ s=1,\ t\ge3&3-t\\
 r=0,\ s=1,\ t\le2&3-t-1-{\one}_{t=1}.
\end{array}
\]
在最后一种情形，$k\ge4$ 将推出 $n\ge N^{{h_{\rm s}}+{d_{\rm s}}+2/3}>N$，故 $k=3$；此外 $q_1q_2q_3^2\ge N^{{h_{\rm s}}+{d_{\rm s}}+2/3}>N$。因此素数对 $(q_1,q_2)$ 对 $Y_9$ 有贡献，给出式中额外的减项。所有上界均非正。对 $k\le2$ 直接代入即得 $W\le4$；单位元满足 $W(1)=4$。
\end{proof}

设 $\mathscr E=\mathscr E_0\cup \mathscr E_1$ 为附录~\ref{app:numeric} 所指定的有理矩形的不交并。两族分别有 $6091$ 和 $11481$ 个成员，其坐标为
\[
 (t,u)=(\log_Np_2,\log_Np_1),\qquad {a_{\rm s}}\le u<{h_{\rm s}}\le t<{d_{\rm s}}.
\]
矩形内部互不相交；公共边在极限中质量为零，并在有限分割中只分配一次。用 $Y_6(\mathscr E)$ 表示相应的素数对项，用 $Y_{11}(\mathscr E)$ 表示最后两个选定素数 $(p_3,p_4)$ 满足 $(\log_Np_4,\log_Np_3)\in \mathscr E$ 的四素数前缀项。定义权
\begin{equation}
                         W_{\mathscr E}=W-Y_6(\mathscr E)+Y_{11}(\mathscr E).
 \label{eq:current-paired-weight}
\end{equation}

\begin{lemma}[配对单射]
\label{lem:current-paired-injection}
对无平方因子的 $n\le N$，有 $Y_{11}(\mathscr E)\le Y_6(\mathscr E)$，从而 $W_{\mathscr E}(n)\le4{\one}_{\Omega(n)\le2}$。
\end{lemma}

\begin{proof}
将一个有贡献的四素数前缀映到其最后一对 $(p_3,p_4)$。商的粗糙性迫使最先选取的两个素数是 $n$ 的最小两个素因子。像落在规定的矩形内，且除去 $p_3p_4$ 后的余项仍为 $N^{a_{\rm s}}$-粗糙数，故该素数对计入 $Y_6(\mathscr E)$。由于前两个素因子由 $n$ 唯一确定，最后一对相同的两个元组必相同，因而此映射为单射。
\end{proof}

固定的素因子个数界 $\lfloor1/{a_{\rm s}}\rfloor=11$ 使粗糙 $n$ 上的 $|W_{\mathscr E}|$ 不超过 $1016$。对充分大的 $N$，含有至少为 $N^{a_{\rm s}}$ 的重复素因子的 $n\le N$ 至多有 $2N^{1-{a_{\rm s}}}$ 个。限制为满足 $(p,N)=1$ 的素数 $p<N$ 并去掉单位元后，含平方因子情形及单位元的代价至多为 $2032N^{1-{a_{\rm s}}}+4=o(H_N)$。于是
\begin{equation}
 4D_{1,2}(N)\ge
       \sum_{\substack{p<N\\p\ {\rm prime}\\(p,N)=1}}W_{\mathscr E}(N-p)
        -o(H_N).
 \label{eq:current-weight-to-count}
\end{equation}

\section{下界筛与有限双重筛转移}
\label{sec:current-lower}

\subsection{原始素数序列的下界估计}

普通下界树从 $w={a_{\rm s}}$、$w={h_{\rm s}}$ 开始，上限为 $\nu=249/448$，并留固定的小余量。其首次退出为偶数阶。与 \eqref{eq:current-upper-exit} 对应的恒等式是
\begin{equation}
 \sum_d\mu(d){\one}_{d\ {\rm retained}}\xi(d)
 =\prod_p(1-\xi(p))
  -\sum_{e\ {\rm even\ first\ exit}}\xi(e)
                       \prod_{p<p_{\rm last}(e)}(1-\xi(p)).
 \label{eq:current-lower-exit}
\end{equation}
该恒等式同时证明逐点下界和实际原始密度公式。选定的第四个自身立方区域恢复仍是合法的偶数阶退出：上限在第四个立方检验之前选定，$S_3+3t_4<\lambda_4$，且普通剩余筛级超过 $p_4^2$。恢复的薄层与五阶节省量在下界泛函中分配到不交区域。

对预筛 $z_0=N^{\eta_0}$，偶数阶退出给出精确的投影方向
\begin{equation}
 V_N(z_0)R_{\rm rough}-V_{\rm full}
  =\sum_{\substack{e\ {\rm full\ even\ exit}\\p_{\rm last}(e)<z_0}}
                              g_N(e)V_N(p_{\rm last}(e))\ge0.
 \label{eq:current-lower-projection}
\end{equation}
将正的偶数阶系数向内取整，负的奇数阶系数向外取整，并保留前缀决策历史。对正系数接入 beta 下界，对负系数接入 beta 上界。若取整后的系数为 $a_{\rm rnd}(d)$，则逐项有 $a_{\rm rnd}(d)B_{\rm sign}(n)\le a_{\rm rnd}(d)I_0(n)$。故实际下界系数满足
\[
 \sum_{q\mid n}\rho_N^{\rm low}(q)\le
           {\one}_{(n,P_N(N^w))=1},\qquad
 \sum_q\frac{\rho_N^{\rm low}(q)}{\varphi(q)}
              \ge V_{\rm full}-\text{系数实现误差}.
\]
beta 下界本身可以为负。第二个不等式使用 \eqref{eq:current-lower-projection} 及绝对误差控制，而不只依赖系数的大小关系。系数定理对这一实际下界系数给出 \eqref{eq:current-PP}。

原始素数序列在 $N$ 处具有锐截断端点。为将光滑分布估计用于带符号的下界系数，先以绝对值控制端点误差。

\begin{lemma}[原始素数的尖锐端点]
\label{lem:current-sharp-prime}
假设实际系数满足 $|\rho_N^{\rm low}(q)|\le M$，其中 $M$ 固定；其支撑为满足 $\theta<1$、$(q,N)=1$ 的无平方因子 $q\le N^\theta$，且满足 \eqref{eq:current-PP} 中允许的单素数光滑情形。则
\begin{equation}
 \sum_q\rho_N^{\rm low}(q)
    \left(\pi(N;q,N)-\frac{\pi(N)}{\varphi(q)}\right)
      =O_A\bigl(N(\log N)^{-A}\bigr).
 \label{eq:current-sharp-pi}
\end{equation}
\end{lemma}

\begin{proof}
在 $[N-\delta N,N+\delta N]$ 范围内光滑化端点，取 $\delta=(\log N)^{-B}$、$B>A+3$，再应用单素数光滑估计。附录~\ref{app:cross-review} 中的绝对边界估计及精确互素均值估计将差值控制为 $O(\delta N\log N+\log^2N)$，另加完整乘积截尾的固定幂次误差。当 $N>2$ 为偶数时，等式 $p=N$ 不可能成立。所选对数精度足以吸收这些误差。
\end{proof}

将逐点下界系数用于原始素数序列。引理~\ref{lem:current-sharp-prime} 控制余项，式 \eqref{eq:current-raw-density} 与素数定理计算主项，由此得到下文使用的普通素数计数下界。

\subsection{有限改进量与实际后继}

令 $A_+(s)=sF(s)/(2e^\gamma)$、$A_-(s)=sf(s)/(2e^\gamma)$。对非负有限带标记族 $\sigma(B)$ 及一个共同的 $Q=N^\lambda$，采用归一化
\[
 \Theta_Q=4\operatorname{li}(N)
       \sum_B\frac{\sigma(B)C_{BN}}{\varphi(B)\log(Q/B)}.
\]
其中 $C_{BN}$ 是与 \eqref{eq:current-normalization} 相同的奇异因子函数。改进量 $H(s)$ 表示上界 $(A_+(s)-H(s)+o(1))\Theta_Q$；改进量 $h(s)$ 表示下界 $(A_-(s)+h(s)+o(1))\Theta_Q$。这些是对非负带标记族平均后的估计，并非关于 $B$ 的逐点结论，也不是对带符号外层测度的断言。

对剩余指数 $L=\lambda-\log_NB$，记 $\alpha=(1-\log_NB)/L\ge2$。指数为 $Lu$ 的单素数子节点满足
\[
 L'=L(1-u),\qquad \alpha'=1+\frac{\alpha-1}{1-u};
\]
二素数子节点则将 $u$ 替换为 $u+v$。所有子节点保留相同的总筛级 $Q$、每个完整及选定平方条件，以及实际素数上限。参数 $\alpha$ 随转移增大，在换位质量中不会重新设为 $2$。

引理~\ref{sv:finite-seed} 在满足 $1\le s\le11/5$ 的可用带标记状态上证明初始改进量 $J_0=1/50$ 以及非负转移 $J\mapsto0.020249+0.13J$。其假设保留剩余筛级、前缀平方条件和共同的比较截点。应用四次得到 $J_4=0.023273777353>2327/100000$，作为下述有限递推的初值。有限集合包括 $70$ 个单步及 $36$ 个双步不等式，以及非负 Buchstab 转移
\[
 H_{\rm new}(s)=\int_{\max(2,s-1)}^5\frac{h(w)}w\,dw,
 \qquad
 h_{\rm new}(s)=\int_{s-1}^5\frac{H(w)}w\,dw\quad(s\ge2).
\]
单步换位代价为 $KI/8$，其中 $I$ 是附录~\ref{app:lower} 指定作用的非负换位质量。双步作用除已证明的普通项与非负换位基础项外，还有非负改进量贡献
\begin{equation}
 \frac15\left[4H(s_0)+H(\phi_1)
  +\sum_{\mathcal D_1}\int\frac{h(\rho)}{u(1-u)}\,du
  +\sum_{\mathcal D_2}\iint\frac{H(\rho)}{uv(1-u-v)}\,du\,dv\right]
 \label{eq:current-w2-feedback}
\end{equation}
共有三个单素数区域和四个良好素数对区域。令 $a_i=1/\phi_i$、$a_0=1/s_0$，则素数对区域为
\[
 \begin{array}{c|c|c}
 &\text{区域}&\rho_i(u,v)\\ \hline
 D_5^+&a_0<u<v<a_2,\ v+2u<1&s_0(1-u-v)\\
 D_6&(a_0,a_1)\mathbin{\times}(a_2,a_3)&s_0(1-u-v)\\
 D_7&a_0<u<v<a_1&(1-u-v)/u\\
 D_8&(a_0,a_1)\mathbin{\times}(a_1,a_2)&(1-u-v)/u.
 \end{array}
\]
可用性条件包括 $2v<1$ 和 $v+2u<1$。$D_5$ 中被排除的部分只承担非负乘积代价，不获得普通反馈。换位项中的单位商曲面保留在其质量中。每个实际 $\alpha\ge2$ 在 $s\ge1$ 时使用包络 $\omega(s)\le\max(1/s,70893/125000)$；低于 $1$ 的保守延拓见附录~\ref{app:lower}。

\paragraph{两层递推。}
四步标量递推生成用于初始化数组的已证常数，发生在数组递推的第零代之前。$H_m,h_m$ 中的代数，包括 $m=248$、$430$，只计从该初值出发、对两个单元数组作同步更新的次数。因此，标量种子的生成与包络数组的更新具有不同的状态空间和计数起点。

以同步的有限归纳定义包络。每个右端均只使用前一代，在备选估计中取最大值；当自变量跨越单元边界时，在完整闭支撑上取最小值。$200$ 个闭单元的端点为 $1+i/50$，$0\le i\le200$。记 $H_m,h_m$ 为 $m$ 代之后的阶梯下包络，共同端点处取相邻值的较小者。从种子直接出发，$430$ 轮确定全部 $400$ 个最终条目。第 $430$ 代数组仅用于 $U_3,U_4$，其余出现处使用第 $248$ 代。

有限归纳还需要共同逼近，而不只是有限深度。非负转移质量至多为 $12$，根部系数范数至多为 $7/2$。归一化之前，局部定向误差 $\delta$ 的传播至多为
$\delta(12^{441}-1)/11$
将 $\delta$ 选在规定的最终预留量内，并对所有查询采用同一嵌套有限细分。有限转移后的剩余指数仍大于 $101^{-440}$ 的某个固定正倍数，故每个非负叶节点均满足命题~\ref{prop:current-positive-product} 的固定截点假设。零质量单元用绝对误差控制。

\subsection{十一项权的计算}

下界泛函为
\begin{equation}
 \mathcal C=\frac14
   (3L_1+L_2+L_5+L_6-U_3-U_4-2U_7-U_8-U_9-U_{10}-U_{11}).
 \label{eq:current-root-functional}
\end{equation}
下界起点为 ${a_{\rm s}},{h_{\rm s}}$ 处依前缀确定的起点。${h_{\rm s}}$ 的普通估计在从参考表达式中去掉完整的最小素因子修正后代入。在 $U_3,U_4$ 中，$H_{430}$ 的自变量为 $\max(2,L/{a_{\rm s}})$，且增强范围内的 $L=v(t)-t$ 保留。区域 $L/{a_{\rm s}}\le2$ 具有正测度，并使用下包络在 $2$ 处的值。素数对下界项在 $L/{a_{\rm s}}\ge2$ 时使用 $h_{248}$，低于 $2$ 则使用零。非负代价 $U_7,U_8,U_9$ 保留带标签重数，并用命题~\ref{prop:current-positive-product} 估计。

对于 $U_{10},U_{11}$，在四素数前缀区域的一个共同非负分割上处理，坐标按 $x>y>z>w$ 递减。在每个单元上，取换位核 $K\omega((1-S_4)/z)/z$ 与可用的普通核 $4(1-H_{248}(\max(2,L/w)))/L$ 的最小值。每个普通备选均保留全部完整及选定平方条件。细分估计替换其父估计；相交区域上的备选只作比较，不相加。附录~\ref{app:lower} 给出这些状态及其保持前缀决定的共同逼近。

\eqref{eq:current-paired-weight} 的限制除了逐点单射，还需要估计下界泛函。对素数对矩形 $\mathscr B=[t_l,t_r]\times[u_l,u_r]$，令
\[
 R(t,u)=\frac{57}{64}-\frac52t-u,\qquad
 R_{\max}=R(t_l,u_l),\quad s_{\max}=R_{\max}/{a_{\rm s}},\quad
 M_{\mathscr B}=\log(t_r/t_l)\log(u_r/u_l).
\]
用 $\mathcal L_{\mathscr B}$ 表示分配给矩形 $\mathscr B$ 的逐单元下界估计量。它在 $\mathscr E_0$ 上为零。在 $\mathscr E_1$ 中，有 $227$ 个跨界单元，其值为零；另有 $11254$ 个正单元，满足 $2\le s_{\min}\le s_{\max}\le9/4$。由于 $\log(s-1)/s$ 在该处递增，且 $h_{\mathscr B}^{\min}$ 表示完整闭自变量支撑上的逐单元 $h_{248}$ 最小值，所以按陈氏权四分之一归一化，被删除的下界估计量至多为
\begin{equation}
 \sum_{\mathscr B\in \mathscr E_1\ {\rm positive}}
       \frac{M_{\mathscr B}}{R_{\max}}\bigl(\log(s_{\max}-1)+h_{\mathscr B}^{\min}\bigr)
       \le 0.002372094.
 \label{eq:current-crop-loss}
\end{equation}

对负项，在 $48$ 等分的四素数前缀覆盖中对同一组选定父节点换位，并避开 $40000$ 个父单元上的细分。若 $A_{\mathscr B}$ 为参考密度质量，$s_{\mathscr B}$ 为其固定的 $K$ 反馈斜率，附录~\ref{app:numeric} 的逐单元界给出 $0\le s_{\mathscr B}\le A_{\mathscr B}/K_0$、$K_0=6.7619$，因此
\[
 A'_{\mathscr B}=A_{\mathscr B}-(K_0-K)s_{\mathscr B}\ge(K/K_0)A_{\mathscr B}.
\]
除以正的参考质量后，得到换位核的显式逐点密度上界。将这一上界限制到互补的非负族后，按陈氏权归一化，至少可去掉 $KI_{\mathscr E}/8$，其中
\begin{equation}
 I_{\mathscr E}=\int_{(t,u)\in \mathscr E}\int_{{a_{\rm s}}<w<z<u}
                      \frac{dt\,du\,dz\,dw}{tu z^2w}.
 \label{eq:current-crop-integral}
\end{equation}
由于整个族上均有 $(1-S_4)/z\ge77/38>2$，故可使用 $\omega\ge1/2$。对一个矩形，该积分恰为
\[
 \log(t_r/t_l)\left[
 \frac{\log(u_r/u_l)}{a_{\rm s}}+\frac{\log(u_r/{a_{\rm s}})+2}{u_r}
                         -\frac{\log(u_l/{a_{\rm s}})+2}{u_l}\right].
\]
总改善量的有限下界分别为：$\mathscr E_0$ 上的 $0.002095161$ 和 $\mathscr E_1$ 上的 $0.003799514$。减去 \eqref{eq:current-crop-loss} 后，$\mathscr E_1$ 带来的原始增量为
\begin{equation}
          0.003799514-0.002372094
             =0.001427420=\frac{71371}{50000000}.
 \label{eq:current-new-crop-gain}
\end{equation}
配对限制从非负上界中移去质量，并支付同一被删区域上的已分配下界估计量。掩膜 $\mathscr E_0,\mathscr E_1$ 互不相交。

\subsection{初步下界泛函}

记此时的下界端点为 $R_{\rm base}$，有限密度逼近与系数实现的每项误差界为 $B_{\rm base}$，见 \eqref{eq:base-R}--\eqref{eq:base-B}。在 $\mathscr E$ 上还可作两项改进：提高内纤维上 $\omega$ 的下界，以及减小被舍弃下界估计量所需扣除的数值。下面分别处理，再合并其贡献。

\subsection{配对区域的改善}
\label{sec:current-refinement}

\begin{lemma}[Buchstab 函数的全局下界]
\label{lem:current-omega-refinement}
当 $s\ge1$ 时有 $\omega(s)\ge1/2$；当 $s\ge7/3$ 时有 $\omega(s)\ge11/20$。
\end{lemma}

\begin{proof}
下界 $1/2$ 在 $[1,2]$ 上成立，并由 $(s(\omega(s)-1/2))'=\omega(s-1)-1/2$ 向后传播。对 $7/3\le s\le3$，利用 $x\ge1$ 时的 $\log x\ge2(x-1)/(x+1)$，得
\[
 \omega(s)=\frac{1+\log(s-1)}s\ge\frac{3s-4}{s^2}\ge\frac{11}{20}.
\]
最后一个不等式来自凹多项式 $60s-80-11s^2$ 在两个端点均为正。当 $3\le s\le10/3$ 时，第一个下界和 $\log2\ge2/3$ 给出
\[
 s\omega(s)\ge1+\log2+(s-3)/2\ge s/2+1/6\ge11s/20.
\]
从而在长度为一的区间 $[7/3,10/3]$ 上建立较强下界。恒等式 $(s(\omega(s)-11/20))'=\omega(s-1)-11/20$ 随后将其传播到连续的各个单位区间。
\end{proof}

\paragraph{额外的非负纤维贡献。}
在 ${a_{\rm s}}<w<z<u<{h_{\rm s}}<t<{d_{\rm s}}$ 上，已经选定的父节点上界控制 $K\omega((1-t-u-z-w)/z)/z$。对同一掩膜的每个矩形 $\mathscr B=[t_l,t_r]\times[u_l,u_r]$，令
\begin{equation}
 Z_{\mathscr B}=\min\left(u_l,\frac{3(1-t_r-u_r)}{13}\right).
 \label{eq:current-Z}
\end{equation}
有限的精确几何检验表明每个单元都有 $Z_{\mathscr B}>{a_{\rm s}}$。此外，由定义有 $1-t_r-u_r-2Z_{\mathscr B}\ge(7/3)Z_{\mathscr B}$。因此在内纤维 ${a_{\rm s}}<w<z<Z_{\mathscr B}$ 上，满足 $z<u$ 以及
\[
 1-t-u-z-w\ge1-t_r-u_r-2z\ge\frac73z.
\]
引理~\ref{lem:current-omega-refinement} 在 \eqref{eq:current-crop-integral} 已使用的 $1/2$ 之外，再为 $\omega$ 提供 $1/20$。每条纤维的其余部分保留原有下界。这两个非负贡献同时成立。

将新增的 $1/20$ 积分，并计入陈氏因子 $1/4$，则 $\mathscr B$ 至少提供下列额外系数：
\begin{align}
 G_{\mathscr B}
 &=\frac{K}{80}
   \log\frac{t_r}{t_l}\log\frac{u_r}{u_l}
      \int_{a_{\rm s}}^{Z_{\mathscr B}}\frac{\log(z/{a_{\rm s}})}{z^2}\,dz \notag\\
 &=\frac{K}{80}
   \log\frac{t_r}{t_l}\log\frac{u_r}{u_l}
   \left[\frac1{a_{\rm s}}-\frac{1+\log(Z_{\mathscr B}/{a_{\rm s}})}{Z_{\mathscr B}}\right].
 \label{eq:current-extra-fibre}
\end{align}
内层 $w$ 积分为 $\log(z/{a_{\rm s}})$，剩余原函数为 $-(1+\log(z/{a_{\rm s}}))/z$。因此没有排列因子，也没有额外的因子 $2$。\eqref{eq:current-Z} 中的固定有理边界仍属于同一个可用有限非负族。这里使用非负父上界及其指定的单元分解施加限制。

附录~\ref{app:refinement} 的定向区间估计给出
\begin{equation}
 \sum_{\mathscr B\in \mathscr E_0\cup \mathscr E_1}G_{\mathscr B}>
 \frac{558583}{10^9}=0.000558583=:g.
 \label{eq:current-extra-gross}
\end{equation}
下文只使用这个有理下端点。

\paragraph{被删除的非负下界估计量。}
分配给矩形 $\mathscr B$ 的下界估计量，等于普通项的常数估计与弦线估计的最大值，再加上其非负双重筛改进量。用 $L_{\mathscr E}$ 表示在 $\mathscr E$ 上包含陈氏因子 $1/4$ 的完整求和。有理定点计算给出
\begin{equation}
 L_{\mathscr E}=\frac{5645403422587826010135}{2417851639229258349412352}
 <\frac{466977}{200000000}=0.002334885=:\lambda_{\rm loss}.
 \label{eq:current-literal-loss}
\end{equation}
下面直接用有理区间给出这一被删下界泛函的上界。

为明确这一包围，在正单元上采用初步泛函中的剩余函数 $R(t,u)=57/64-(5/2)t-u$，并令
\begin{align*}
 r_{\mathscr B}&=R(t_r,u_r),\quad R_{\mathscr B}=R(t_l,u_l),\quad
 s_l=r_{\mathscr B}/{a_{\rm s}},\quad s_h=R_{\mathscr B}/{a_{\rm s}},\\
 f_l&=\frac{\log(s_l-1)}{s_l},\qquad
 f_h=\frac{\log(s_h-1)}{s_h}.
\end{align*}
这里 $2\le s_l\le s_h\le9/4$。定义
\begin{align*}
 M_{\mathscr B}&=\log(t_r/t_l)\log(u_r/u_l),\\
 M_{R,\mathscr B}
 &=\int_{\mathscr B} R(t,u)\frac{dt\,du}{tu}\\
 &=\frac{57}{64}M_{\mathscr B}
   -\frac52(t_r-t_l)\log(u_r/u_l)
   -(u_r-u_l)\log(t_r/t_l).
\end{align*}
两个普通备选在除以陈氏归一化因子 $4$ 后，分别不超过
\begin{equation}
 \frac{f_lM_{\mathscr B}}{{a_{\rm s}}},
 \qquad
 \frac{f_l(R_{\mathscr B}M_{\mathscr B}-M_{R,\mathscr B})+
       f_h(M_{R,\mathscr B}-r_{\mathscr B}M_{\mathscr B})}{a_{\rm s}(R_{\mathscr B}-r_{\mathscr B})}.
 \label{eq:current-floor-proxy}
\end{equation}
事实上，精确端点值 $f_l,f_h$ 高于泛函实际使用的下端点。因 $r_{\mathscr B}\le R(t,u)\le R_{\mathscr B}$，两个系数积分均非负，而向下的定点取整只会减小所记录的值。逐单元估计量对常数与弦线备选取最大值；对 \eqref{eq:current-floor-proxy} 中的两个上界也取最大值，方向不变。

对单独的改进量项，取 $h_{\mathscr B}^{\min}$ 为与精确区间 $[s_l,s_h]$ 相交的 $h_{248}$ 闭单元值的最小值，节点处包括两侧单元。这个最小值不小于向外取整后查询区间上的最小值。因此 $h_{\mathscr B}^{\min}M_{\mathscr B}/R_{\mathscr B}$ 是逐单元改进量估计的上界，可加到前述最大值上。所有相关 $h_{248}$ 值均非负。若向外查询跨到 $2$ 以下而返回零，同一上界仍成立。在 $227$ 个跨界单元上，常数、可用弦线备选以及查询结果均为零。每个 $\mathscr E_0$ 单元都满足 $s_h\le2$，故其逐单元下界估计量也为零。

这些控制量之和具有严格小于 $\lambda_{\rm loss}$ 的有理上包围。附录~\ref{app:refinement-review} 通过范围归约及带显式余项的交错对数级数得到此界。因此，可用 $\lambda_{\rm loss}$ 替换 \eqref{eq:current-crop-loss} 中扣除的数值，而不破坏下界不等式。

回收的代价恰为
\begin{equation}
 0.002372094-\lambda_{\rm loss}=\frac{37209}{10^9}=0.000037209.
 \label{eq:current-recovered-loss}
\end{equation}
保留的 $L_6$ 下端点等于完整 $L_6$ 下端点减去 $4\lambda_{\rm loss}$，即
\[
 \frac{1421673188323272382945058469141}
      {472236648286964521369600000000}.
\]
按陈氏权归一化，移去的负项总量现在可安全取为 $0.002095161+0.003799514+g=0.006453258$。结合 \eqref{eq:current-extra-gross} 与 \eqref{eq:current-recovered-loss}，得到额外的下界系数
\begin{equation}
 \Delta=g+(0.002372094-\lambda_{\rm loss})
 =\frac{37237}{62500000}=0.000595792.
 \label{eq:current-refinement-delta}
\end{equation}

我们给两项误差预留量各增加 $\Delta/8=37237/500000000=0.000074474$，分别用于有限密度逼近和系数实现。扣除两项新增预留量后的净增量为
\begin{equation}
 \Delta-2(\Delta/8)=\frac{3\Delta}{4}
                  =\frac{111711}{250000000}=0.000446844.
 \label{eq:current-refinement-net}
\end{equation}
边界和查询均为有限个，因此可沿用初步下界泛函所用的共同逼近与系数实现论证，将其局部误差选在这些新增的正预留量内。

\subsection{下界端点与极限次序}
\label{sec:lower-limits}

计入配对区域的改善后，原始端点及各项实现预留量为
\begin{align}
 R_{\rm low}&=R_{\rm base}+\Delta,\label{eq:current-R}\\
 B_{\rm low}&=B_{\rm base}+\Delta/8.\label{eq:current-B}
\end{align}
精确值见附录~\ref{app:numeric}。由 \eqref{eq:current-weight-to-count}、实际下界系数的实现、有限归纳以及配对下界泛函的计算，得
\begin{equation}
 \liminf_{\substack{N\to\infty\\2\mid N}}
         \frac{D_{1,2}(N)}{H_N}\ge R_{\rm low}-2B_{\rm low}.
 \label{eq:current-lower-derived}
\end{equation}
扣除两项预留量后仍有
\begin{equation}
                  R_{\rm low}-2B_{\rm low}>c+4.9061\cdot10^{-9},
                  \label{eq:current-after-budgets}
\end{equation}
特别地，
\begin{equation}
                  R_{\rm low}-2B_{\rm low}-c>4.9061\cdot10^{-9}>0.
                  \label{eq:current-lower-slack}
\end{equation}
因而最终有 $D_{1,2}(N)\ge cH_N$，其中 $c$ 由 \eqref{eq:current-safe-constants} 给出。

先固定有限组估计及其依赖关系、尾部和紧区域上的逼近精度、确定固定紧参数域的截断参数、前缀筛级上限及支撑余量，再选择共同的 beta 与粗糙参数、细标签、固定除数阶数与矩阶数、允许的 $\kappa$ 和局部对数精度。此后才令 $N$ 沿偶数趋于无穷。取有限多个阈值的最大值，使全部估计同时成立。由于 $C_N\ge C_2>0$，固定幂次节省及 $A>2$ 时的 $O_A(N/\log^AN)$ 均一致为 $o(H_N)$。因此，有限归纳给出对每个充分大偶数目标成立的结论。

\section{有限数值估计}
\label{sec:current-numeric}

上述上、下界构造将证明中的数值部分归结为以下有限不等式。表~\ref{num:certificate-table} 提前列明这些量与核验所用证书字段的对应关系。附录~\ref{app:numeric} 定义相应的积分、分割、定向包围及误差预留量。

\begin{lemma}
\label{lem:numerical-endpoints}
上界密度泛函具有 \eqref{eq:current-E} 中的上端点 $E_{\rm up}$。它允许有限密度实现和系数实现，两项误差各至多为 $u$，其中 $u$ 由 \eqref{eq:current-u} 给出，且 $E_{\rm up}+2u<K$。配对删除后的十一项下界泛函具有 \eqref{eq:current-R} 中的下端点 $R_{\rm low}=R_{\rm base}+\Delta$。相应两项误差各至多为 $B_{\rm low}$，其中 $B_{\rm low}$ 由 \eqref{eq:current-B} 给出，且 $R_{\rm low}-2B_{\rm low}>c$。$U_3,U_4$ 中使用的第 $430$ 代包络，由所给初值同时施加非负有限作用映射 $430$ 次得到。其余指定项使用的第 $248$ 代包络，具有附录~\ref{app:numeric} 所述的独立包围。
\end{lemma}

\begin{proof}
对上界，计算四个代价项
\[
 E_{\rm up}=\frac4\Theta+I_1+I_{3,\mathrm{compact}}+I_{3,\mathrm{tail}}
       -\sum_{j=1}^{22}G_j,
\]
其中四项代价取上端点，非负节省量 $G_j$ 取下端点。每次连续变更都相对于其父节点计算；重叠备选在积分之前先取最大值。非紧尾部在其完整范围上估计。

对下界，在 \eqref{eq:current-root-functional} 中对 $L_1,L_2,L_5,L_6$ 取下端点，对 $U_3,U_4,U_7,\ldots,U_{11}$ 取上端点。下界起点的变化、退出区域的恢复和有限表修正，均在规定的不交区域上插入，或作为完整的父节点替换。再用 \eqref{eq:current-crop-integral}、\eqref{eq:current-crop-loss}、\eqref{eq:current-extra-gross} 和 \eqref{eq:current-refinement-delta} 计算配对区域项，得到有理端点 $R_{\rm low}$。

全部标量运算均为有理运算。对数用带显式余项的级数包围；筛函数用有方向的插值估计包围；仿射单元利用精确顶点和精确多项式矩进行三角剖分。非负节省量可以只在子集上估计，并放弃其余部分；负项代价的上界则必须覆盖完整保留区域。有限递推只读取前一代，并将每个非负贡献向下取整。具体的包围及覆盖条件见附录~\ref{app:numeric}。有理数相减即得 \eqref{eq:current-upper-slack} 和 \eqref{eq:current-lower-slack}。
\end{proof}

\paragraph{数值数据。}
附录~\ref{app:numeric} 所用的有限分割、递推及定向积分界收录于 \emph{Theorem (1+1.75) on the Goldbach Conjecture --- Numerical Supplement}，版本 20261106~\cite{ChanSupplement}。该附录说明数据及其证书与引理~\ref{lem:numerical-endpoints} 的对应关系。

\section{不满足限制的有序半素数与最终相减}
\label{sec:current-transfer}

现在估计 $B_b(N)$。其对 $b$ 的依赖包含在有序半素数的未筛质量中，随后用非负乘积估计加入补数为素数的条件。

固定 $1/2<b<1$，令 $\beta=b/(1+b)$。低于 $2N^{1-\kappa_0}$ 的乘积计数可忽略。对其余不满足限制的乘积 $m=rq$，不等式 $r>q^b$ 推出 $r>m^\beta$，故两个素数都有固定的正指数。采用坐标
\[
                       t=\frac{\log r}{\log N},\qquad y=\frac{rq}{N}.
\]
在固定径向条带 $\eta\le y\le C$ 上，先对 $q=Ny/r$ 使用素数定理，再对 $r$ 作素数倒数的分部求和，得到未筛质量元
\begin{equation}
                         \frac{N}{\log N}\,
                              \frac{dt\,dy}{t(1-t)}.
 \label{eq:current-pair-measure}
\end{equation}
素数定理误差在正指数紧集上一致。素数上界将 $y<\eta$ 的贡献控制为 $O_b(\eta N/\log N)$；先令 $N$ 趋于无穷，再令 $\eta$ 趋于零。

不满足限制且满足排序的精确条件为
\[
 0<y<1,\qquad
 t\le\frac12+\frac{\log y}{2\log N},\qquad
 t>\beta+\frac{b\log y}{(1+b)\log N}.
\]
其极限未筛质量是
\begin{equation}
 J_b=\int_0^1dy\int_\beta^{1/2}\frac{dt}{t(1-t)}
     =\left[\log\frac{t}{1-t}\right]_\beta^{1/2}
     =\log(1/b).
 \label{eq:current-bad-mass}
\end{equation}
排序 $r\le q$ 使每个合数只计一次。$J_b$ 中既无因子 $2$，也无额外奇异因子。两个素数相等的对角线大小为 $O(\sqrt N)$，可以保留，或在无平方因子的论证后恢复。

将 $y$ 中的完整乘积边界及 $t$ 中的其余边界光滑化。对任意固定 $\epsilon>0$，存在极限质量至多为 $J_b+\epsilon$ 的非负外逼近，且仍属于可用有限族。单位分解不重复计入质量。应用命题~\ref{prop:current-positive-product}，先令 $N\to\infty$，再令 $\epsilon\downarrow0$，得到
\begin{equation}
 \limsup_{\substack{N\to\infty\\2\mid N}}
             \frac{B_b(N)}{H_N}\le K\log(1/b).
 \label{eq:current-bad-bound}
\end{equation}

唯一分解给出精确恒等式
\begin{equation}
                         D_b(N)=D_{1,2}(N)-B_b(N).
 \label{eq:current-exact-subtraction}
\end{equation}
\eqref{eq:current-d12} 已排除单位元；即使采用包含单位元的约定，计数也至多改变一。$p<N$ 排除了零。当 $b<1$ 时，平方余项不满足限制，素数余项则满足限制。结合 \eqref{eq:current-lower-derived}、\eqref{eq:current-lower-slack} 与 \eqref{eq:current-bad-bound}，即证明定理~\ref{thm:current-main}。

最后验证推论~\ref{cor:three-quarters}、\ref{cor:current-explicit} 的严格标量不等式。对有理数 $b$，令 $z=(1-b)/(1+b)$。正项反双曲正切级数及其几何尾界给出
\begin{equation}
 2\sum_{j=0}^{95}\frac{z^{2j+1}}{2j+1}
 <\log(1/b)
 <L_{96}(b):=
 2\sum_{j=0}^{95}\frac{z^{2j+1}}{2j+1}
       +\frac{2z^{193}}{193(1-z^2)}.
 \label{eq:current-log-certificate}
\end{equation}
所有端点均为有理数。代入 \eqref{eq:current-safe-constants}，得
\begin{align}
 c-KL_{96}(3/4)&>\frac{16532064}{10^9}>0.0165,\notag\\
 c-KL_{96}(74816964/10^8)
     &>\frac{1001149712517311}{10^{20}}>10^{-5},\notag\\
 c-KL_{96}(74816865/10^8)
     &>\frac{106416285358807}{10^{20}}>10^{-6},
                 \label{eq:current-rational-margins}\\
 c-KL_{96}(9/10)
     &>\frac{124934244400949969098}{10^{20}}>1.24889.\notag
\end{align}
相应剩余量分别超过 $3.2064\cdot10^{-5}$、$1.1497\cdot10^{-8}$、$6.4162\cdot10^{-8}$ 和 $4.5244\cdot10^{-4}$。先将固定实现误差和非负光滑化误差选在这些严格余量内，再令其余一致渐近误差趋于零。取四个所得阈值的最大值，即证两个推论。

\paragraph{固定常数给出的范围。}
对 $1/2<b<1$，\eqref{eq:current-general-theorem} 右端为正当且仅当 $b>\exp(-c/K)$。该值描述当前这组系数界给出的范围。若要扩大此范围，需要在保留变动剩余类及系数类别条件的同时，加强适用的下界或上界估计。

\appendix
\section{受限中等长度卷积估计}
\label{app:medium}

本附录证明命题~\ref{prop:current-medium}。卷积长度为 $U,V$，记 $x=UV$，非零剩余类记作 $\mathfrak{a}$；正文中的目标偶数记作 $N$。令 $\vartheta=7/32$，并记
\[
 z=x^{1/(\log\log x)^3},\qquad y=x^{1/\log\log x}.
\]
系数 $\alpha_n$ 在全部模数求和过程中固定，并满足辅助互素 Siegel--Walfisz 条件 \eqref{eq:current-aux-sw}。该条件中的常数对系数族一致。采用记号 $e(t)=\exp(2\pi it)$；指数函数中的逆元始终按所显示的分母取模。

\subsection{光滑控制函数与除数估计}

\begin{lemma}\label{ana:majorant}
存在固定函数 $\psi\in C_c^\infty((0,\infty))$，使 $0\le\psi\le1$，在 $[1,2]$ 上有 $\psi=1$，并满足 $\supp\psi\subset[1/2,5/2]$ 及
\[
 \|\psi^{(j)}\|_\infty\le C^j(j!)^2,
 \qquad |\widehat\psi(t)|\ll \exp(-c\sqrt{|t|}).
\]
因此，对具有多项式上界的正数 $\ell,V$ 和 $H\ge C_1\ell(\log x)^5/V$，有
\begin{equation}
 \frac V\ell\sum_{|h|>H}
       \left|\widehat\psi\left(\frac{hV}{\ell}\right)\right|
       \ll_D x^{-D}
 \label{ana:poisson-tail}
\end{equation}
其中 $D>0$ 为任意固定值。
\end{lemma}

\begin{proof}
选择满足 $\sum \eta_j=1/2$ 的 $\eta_j=c_0/j^2$，令 $u_j$ 为 $[-\eta_j/2,\eta_j/2]$ 上的均匀概率测度，并取 $\phi=u_1*u_2*\cdots$，其支撑为 $[-1/4,1/4]$。将前 $j$ 个因子作为带符号测度求导，总变差至多为 $\prod_{i\le j}(2/\eta_i)=C^j(j!)^2$。其余卷积为概率测度；再保留足够多未求导因子，即可在每个固定导数阶数上得到连续密度。函数 $\psi=\one_{[3/4,9/4]}*\phi$ 满足所述支撑与导数界。分部积分 $j$ 次，其中 $j$ 为 $\sqrt{|t|}$ 的充分小倍数，即得 Fourier 估计。用递减函数积分比较，将 \eqref{ana:poisson-tail} 中的尾部控制为一个多项式因子乘以 $\exp(-c_1(\log x)^{5/2})$。若 $H<1$，每个被略去的整数本就满足 $|h|V/\ell\gg(\log x)^5$，同一论证仍然适用。
\end{proof}

记 $n=\log_xU$、$q=\log_xQ_1$、$u=\log_xQ_2$ 和 $v=\log_xQ_3$。命题~\ref{prop:current-medium} 的假设推出
\begin{equation}
 n>\sigma,\qquad n\le\frac12-4\sigma,\qquad
 q+u+v<\frac7{10}-\frac{13\sigma}2.
 \label{ana:lengths}
\end{equation}
特别地，当 $x$ 充分大时有 $Q_1Q_2Q_3\le x^{4/5}$。将 $q_2=u_0s$、$q_3=v_0t$ 分成 $z$-粗糙部分和 $z$-光滑部分。满足 $s>y$ 或 $t>y$ 的项有任意大的对数节省。事实上，先利用固定的除数界合并模数系数，再应用除数有界函数在本原等差数列中的估计~\cite{Shiu}，随后以指数 $1/\log z$ 使用 Rankin 方法，得到
\begin{equation}
 \sum_{\substack{\ell\le x^{4/5}\\ \operatorname{sm}(\ell;z)\ge y}}
       \frac{\tau(\ell)^C}{\ell}
 \ll \exp\bigl(- (\log\log x)^2+O_C(\log\log x)\bigr).
 \label{ana:smooth-tail}
\end{equation}
这里 $\operatorname{sm}(\ell;z)$ 是 $\ell$ 中所有素因子均不超过 $z$ 的最大除数。使用该界时配合除数均值，不引入逐点因子 $x^\epsilon$。

我们还使用以下平移后的等差数列估计。固定 $C,K,\delta>0$，设区间 $I\subset[-x^K,x^K]$ 的长度至多为 $L$，且 $f\le Lx^{-\delta}$，则
\begin{equation}
 \sum_{\substack{v\in I\\v\equiv r\pmod f\\v\ne0}}
        \tau(|v|)^C
       \ll \frac Lf\tau(f)^C(\log x)^{O(1)}.
 \label{ana:shifted-divisor}
\end{equation}
事实上，令 $g=(r,f)$、$v=gw$，所得模 $f/g$ 的等差数列为本原数列，其长度与模数之比为 $L/f\ge x^\delta$。范围 $|w|\le(f/g)x^{\delta/2}$ 用一个固定的除数估计控制。其余部分先按正负分开，再分成二进区间，并细分为长度与其尺度及 $L/g$ 两者中较小者同阶的区间。在每个区间上，长度和模数都留有固定的正余量，可使用本原除数估计；求和即得 \eqref{ana:shifted-divisor}。Cauchy 不等式还给出关于 $\tau(v)^B\tau(|v-\mathfrak{a}|)^C$ 的混合版本，其中排除 $v=\mathfrak{a}$。

在一个二进块上，令 $d=st\asymp D$、$w=u_0v_0\asymp\mathcal R$，并定义
\begin{equation}
 \Lambda_d(w)=\sum_{\substack{st=d\\u_0v_0=w}}
       \gamma_2(u_0s)\gamma_3(v_0t)
       \one_{\text{固定的各变量区间及支撑}}.
 \label{ana:Lambda}
\end{equation}
于是 $\Lambda_d$ 不依赖于 $q_1$，并满足 $|\Lambda_d(w)|\le\tau(dw)$、$D\le y^2=x^{o(1)}$、$D\mathcal R\asymp Q_2Q_3$。由 $Q=Q_1$、$Y=QD$、$\delta=\sigma/3$，假设 \eqref{ana:lengths} 推出
\begin{equation}
 U,V\ge x^\delta,\qquad Y\le Ux^{-\delta},\qquad
 Y\mathcal R\le x^{1-\delta}.
 \label{ana:dispersion-lengths}
\end{equation}
用 $\mathcal S$ 表示限制在该块上的原始分解系数偏差和。施加 Cauchy 不等式之前，$d\asymp D$ 的所有值均保留在同一个求和中。

在应用离散法之前，先归约到模不超过 $1$ 的 $\beta$。对于除数有界的 $\beta$，保留 $\tau(m)\le(\log x)^C$ 的部分，并以 $(\log x)^{BC}$ 归一化。在 $v=nm$ 处被舍弃的系数至多为 $(\log x)^{-C}\tau(v)^{B+2}$。当 $v\ne\mathfrak a$ 时，对模数求和的代价至多为 $\tau(|v-\mathfrak a|)^3$，因此固定阶的平移除数矩给出 $O(x(\log x)^{E-C})$，其中 $E$ 与 $C$ 无关。精确乘积 $v=\mathfrak a$ 的贡献另以 $x^{4/5+o(1)}$ 控制。普通除数均值与倒数 Euler 函数和控制均值项。先选择 $C$，再为归一化估计要求足够的对数精度。

\subsection{离散法与主方差}

离散步骤将共同主项与非零频率分离。剩余的非负主方差先用辅助 Siegel--Walfisz 假设估计，再处理指数和。

令 $\ell=qd$、$\eta_r=\Lambda_d(r)$、$h_0=\int\psi$。对固定 $(q,d)$，定义
\begin{align*}
 B_r&=\sum_{(n,\ell r)=1}\alpha_n,
 &A_r(b)&=\sum_{\substack{n\equiv b\pmod\ell\\(n,r)=1}}\alpha_n,\\
 C_r(m)&=\frac{\one_{(m,\ell r)=1}}{\varphi(\ell r)}B_r,
 &X_r(m)&=\sum_{(n,\ell r)=1}\alpha_n
       \bigl(\one_{mn\equiv \mathfrak{a}\pmod{\ell r}}-\one_{mn=\mathfrak{a}}\bigr).
\end{align*}
所有 $n$ 求和均保留原区间。先去掉 $mn=\mathfrak{a}$，再应用 Cauchy 不等式。一阶矩代价为 $Y\mathcal R x^{o(1)}$，由 \eqref{ana:dispersion-lengths} 可知它具有幂次节省。用 $q$、$m$ 的非负控制函数，得到 Cauchy 乘子 $VY(\log x)^C$ 以及精确方差
\begin{align}
 \mathcal V'&=S_1-2\Re S_2'+S_3',\notag\\
 S_1&=\sum_{q,d}\psi(q/Q)\sum_m\psi(m/V)
                 \left|\sum_r\eta_rC_r(m)\right|^2,\notag\\
 S_2'&=\sum_{q,d}\psi(q/Q)\sum_m\psi(m/V)
       \left(\sum_r\eta_rX_r(m)\right)
       \overline{\left(\sum_s\eta_sC_s(m)\right)},\notag\\
 S_3'&=\sum_{q,d}\psi(q/Q)\sum_m\psi(m/V)
                 \left|\sum_r\eta_rX_r(m)\right|^2.
 \label{ana:variance}
\end{align}

共同主项由核
\[
 c_\ell(r,s)=\frac{\varphi(\ell rs)}{
       \ell rs\varphi(\ell r)\varphi(\ell s)}.
\]
控制。若 $L$ 是 $\ell$ 与有限个支撑内 $r$ 的任意公倍数，则
\begin{equation}
 c_\ell(r,s)=\frac1L\sum_{v\bmod L}
 \frac{\one_{(v,\ell r)=1}}{\varphi(\ell r)}
 \frac{\one_{(v,\ell s)=1}}{\varphi(\ell s)}.
 \label{ana:gram}
\end{equation}
这是一个非负 Gram 核。$L$ 的大小无关紧要，因为它仅用于有限恒等式。光滑互素计数给出
\begin{equation}
 S_1=\mathcal M+O\left(\frac{U^2(\log x)^C}{Y}\right),\qquad
 \mathcal M=Vh_0\sum_{q,d}\psi(q/Q)
        \sum_{r,s}\eta_r\overline{\eta_s}
               c_\ell(r,s)B_r\overline{B_s}.
 \label{ana:mean}
\end{equation}

对 $S_2'$，在 $mn\ne \mathfrak{a}$ 时先去掉 $(m,s)>1$。粗糙公因子 $f>z$ 带来因子 $\mathcal R/f$ 和 $x/f$；应用 \eqref{ana:shifted-divisor} 并对 $\sum_{f>z}\tau(f)^C/f^2$ 求和，误差至多为 $VU^2(\log x)^C/(Yz)$。随后才为 Poisson 求和恢复 $mn=\mathfrak{a}$，代价为 $\mathcal RUx^{o(1)}$。零频系数为 $1/(\ell r\varphi(\ell s))$，其与 $c_\ell(r,s)$ 之比为 $1+O((\log x)/z)$。所以
\begin{equation}
 S_2'=\mathcal M+VE_1+
 O_A\left(\frac{VU^2}{Y(\log x)^A}+\mathcal RUx^{o(1)}\right),
 \label{ana:cross}
\end{equation}
其中保留原有互素条件，有
\begin{align}
 E_1={}&\sum_{q,d}\psi(q/Q)\sum_{r,s}
     \frac{\eta_r\overline{\eta_s}}{\ell r\varphi(\ell s)}
 \sum_{\substack{n_1,n_2\\(n_1,\ell r)=(n_2,\ell s)=1}}
       \alpha_{n_1}\overline{\alpha_{n_2}}\notag\\
 &\hspace{12mm}\times
 \sum_{0<|h|\le H_1}\widehat\psi\left(\frac{hV}{\ell r}\right)
               e\left(\frac{\mathfrak{a}h\overline{n_1}}{\ell r}\right),
 \qquad H_1\asymp\frac{Y\mathcal R}{V}(\log x)^5.
 \label{ana:E1-definition}
\end{align}
引理~\ref{ana:majorant} 保证可以作这一 Fourier 截断。

对 $S_3'$，在所有除数截断及公因子消去过程中都保留两项精确乘积排除条件。整理后的集合满足
\begin{equation}
 \begin{gathered}
 (r,s)=(s,\ell)=1,\quad (n_1,\ell r)=(n_2,\ell s)=1,
 \quad(n_1,n_2)=1,\\
 n_1\equiv n_2\pmod\ell,\qquad
 |n_1-n_2|\ge U(\log x)^{-C_{\rm sep}}.
 \end{gathered}
 \label{ana:clean-set}
\end{equation}
移去粗糙公因子的代价为
\[
VU^2(\log x)^C/(Yz)+x(\log x)^C,
\]
邻近数对的代价为
\[
VU^2(\log x)^{C-C_{\rm sep}}/Y+x(\log x)^C.
\]
这两项分别通过计数两个倍数，或一个倍数与一个本原等差数列得到；附带的单位项用 $Y/U\le x^{-\delta}$ 吸收。先选择除数截断，再令 $C_{\rm sep}$ 充分大。

此时 $n_1\ne n_2$。恢复一个等式 $mn_i=\mathfrak{a}$ 的代价仅为 $\mathcal R(U+Y)x^{o(1)}=\mathcal RUx^{o(1)}$：固定其分解后，另一个 $n$ 只有 $O(U/Y+1)$ 种选择，另一个粗糙模数则整除一个非零的平移乘积。两个等式不能同时成立。仅在恢复之后才使用中国剩余定理与 Poisson 求和。在 \eqref{ana:clean-set} 上有
\[
 \varphi(\ell)c_\ell(r,s)=\frac1{\ell rs}.
\]
令 $d_r(b)=A_r(b)-B_r/\varphi(\ell)$。由于 $\sum_b^*A_r(b)=B_r$，零频为 $\mathcal M+\mathcal P$，其中
\begin{equation}
 \mathcal P=Vh_0\sum_{q,d}\psi(q/Q)\varphi(\ell)
   \sum_{r,s}\eta_r\overline{\eta_s}c_\ell(r,s)
                     \sum_b^*d_r(b)\overline{d_s(b)}\ge0.
 \label{ana:principal-variance}
\end{equation}
非负性由 \eqref{ana:gram} 得到。特别地，这个主方差必须估计，不能声称其恒等抵消。

辅助 Siegel--Walfisz 假设和本原特征大筛给出：对每个固定 $A,w>0$，有
\begin{equation}
 \sum_{\ell\le8Y}\tau(\ell)^w\sum_b^*|d_r(b)|^2
       \ll_A\frac{\tau(r)^{2C_0}U^2}{(\log x)^A}.
 \label{ana:filtered-BDH}
\end{equation}
为明确起见，将 $\alpha_n\one_{(n,r)=1}$ 除以 $\tau(r)^{C_0}$ 归一化。因 $\tau(kr)\le\tau(k)\tau(r)$，其辅助 Siegel--Walfisz 常数仍一致。利用特征正交性，把未加权方差表示为非主特征的和。写每个模数为 $fk$，其中 $f>1$ 为本原导子，并在应用大筛之前把 $(n,k)=1$ 保留在系数内。在导子二进区间 $f\asymp R_0$ 上，固定 $k$ 时的界为 $O(U(U/R_0+R_0)/\varphi(k))$。对 $f>(\log x)^b$ 的导子及诱导因子求和，得 $O(U^2(\log x)^{2-b}+UY)$。小导子由辅助 Siegel--Walfisz 条件控制，带上代价 $\tau(k)^{C_0}$，总量为 $O(U^2(\log x)^{3b+C_1-2B})$。先选 $b$，再选 $B$，即可得到任意指定的对数节省。带权 $\tau(\ell)^w$ 的结论由 Cauchy 不等式、普通除数均值及方差的初等界 $O(U^2/\ell)$ 推出。

利用 \eqref{ana:gram}，对 $r$ 变量使用 Cauchy 不等式，并应用 $\varphi(\ell r)^{-1}\ll(\log x)/(Y\mathcal R)$，对固定除数阶数 $B_0$ 得到
\begin{align*}
 \mathcal P\ll{}&\frac{V(\log x)^2}{Y\mathcal R}
  \sum_{r\asymp\mathcal R}\tau(r)^{2B_0}
  \sum_{\ell\le8Y}\tau(\ell)^{2B_0+1}\sum_b^*|d_r(b)|^2\\
 \ll_A{}&\frac{VU^2}{Y(\log x)^A}.
\end{align*}
\eqref{ana:filtered-BDH} 中的因子 $\tau(r)^{2C_0}$ 用除数均值求和，不放入未说明的一致常数。结合 \eqref{ana:variance}，得
\begin{equation}
 |\mathcal S|\ll_A\frac{x}{(\log x)^A}
  +V\sqrt Y(\log x)^C\bigl(\sqrt{|E_1|}+\sqrt{|E_2|}\bigr),
 \label{ana:dispersion-bound}
\end{equation}
其中 $E_2$ 是 \eqref{ana:clean-set} 上由中国剩余定理产生的非零频率和。恢复的单零项贡献至多为 $\sqrt{xY\mathcal R}\,x^{o(1)}$，具有幂次节省。

\subsection{第一个指数和}

对固定光滑函数 $w$，由 Poisson 求和及完整 Kloosterman 和的 Weil 界，对每个整数 $A$ 有
\begin{equation}
 \left|\sum_{(n,\ell)=1}w(n/U)e(A\overline n/\ell)\right|
 \ll_w \tau(\ell)^2\sqrt{\ell}+\frac U{\ell}(A,\ell).
 \label{ana:inverse-completion}
\end{equation}
零 Fourier 频率为 Ramanujan 和，给出第二项；其余频率使用
\[
 \sqrt{(k,A,\ell)}\le\sum_{\substack{d\mid(A,\ell)\\d\mid k}}\sqrt d,
 \qquad
 \sum_{\substack{k\ne0\\d\mid k}}(1+|k|/V)^{-2}\le\frac{2V}{d}.
\]
因此，即使频率长度小于一，\eqref{ana:inverse-completion} 仍然成立。

固定 \eqref{ana:E1-definition} 中的 $d$，令 $Q'=dQ$，并将 $q'$ 限制为 $d$ 的倍数。该限制可放入依赖于 $q'$ 的有界系数。以下计算暂将 $Q'$ 记为 $Q$，将 $\mathcal R$ 记为 $R$。把有界光滑因子吸收到 $r,h$ 系数，应用 Cauchy 不等式，同时保留 $r,s$ 原有的未加限制支撑。展开平方得到
\begin{equation}
 |E_1(d)|^2\ll x^\epsilon\frac{U^3}{Q^3R^2}\,T,
 \label{ana:E1-cauchy}
\end{equation}
其中
\[
 T=\sum_{\substack{q\asymp Q, r,r'\asymp R\\(\mathfrak{a},qrr')=1}}
       \sum_{0<|h|,|h'|\le H_1}
 \left|\sum_{(n,qrr')=1}w(n/U)
 e\left(\frac{\mathfrak{a}(hr'-h'r)\overline n}{qrr'}\right)\right|.
\]
令 $j=hr'-h'r$。因 $hr'\ne0$，且 $h'r=hr'$ 是一个除数方程，$j=0$ 元组共有 $RH_1x^\epsilon$ 个，贡献为 $\ll x^\epsilon QRH_1U$。对 $j\ne0$，\eqref{ana:inverse-completion} 的平方根项贡献为 $\ll x^\epsilon Q^{3/2}R^3H_1^2$。对 Ramanujan 项，记 $r=gu$、$r'=gv$、$(u,v)=1$，则
\[
 (j,rr')\le g^2(h,r)(h',r'),\qquad
 \sum_{q\asymp Q}(j,q)\le3Q\tau(|j|),
\]
以及
\[
 \sum_{0<|h|\le H_1}(h,r)\le2H_1\tau(r),\qquad
 \sum_{r,r'\asymp R}(r,r')^2\ll R^3.
\]
由于 $(\mathfrak{a},qrr')=1$，$\mathfrak{a}$ 从这些最大公因数中消去。Ramanujan 项对 $T$ 的总贡献为 $\ll x^\epsilon UH_1^2R$。由 $H_1\ll QRU(\log x)^5/x$，\eqref{ana:E1-cauchy} 化为
\begin{equation}
 |E_1(d)|^2\ll x^\epsilon\frac{U^4}{Q^2}
  \left(\frac{UQ}{x}+\frac{UQ^{5/2}R^3}{x^2}
                         +\frac{U^2QR}{x^2}\right).
 \label{ana:E1-polynomial}
\end{equation}

回到 $Q=Q_1$ 和 \eqref{ana:lengths} 中的对数变量。所需余量来自
\begin{align*}
 n+q&<1-9\sigma,\\
 2n+q+u+v&<\frac{17}{10}-\frac{29\sigma}2,\\
 n+\frac52q+3u+3v
 &=\frac{35}{16}(q-n)+\frac{23}{16}(2n+u+2v)
   +\frac5{16}(n+q+5u+2v)-\frac v2\\
 &<\frac{1021}{512}-\frac{147\sigma}{8}+O(1/\log x).
\end{align*}
最后一个不等式使用 $|\mathfrak{a}|\ge c_0x$、$\vartheta=7/32$。此时 $\Lambda_d$、$D=x^{o(1)}$ 的除数界均可由幂次节省吸收。利用 $\sum_{d\asymp D}1/d\ll1$ 对 $d\asymp D$ 中的 $U^2/(dQ)$ 求和，再留出吸收一个额外 $D$ 的余地，得到
\begin{equation}
 |E_1|\ll\frac{U^2}{DQ}x^{-\sigma/2}.
 \label{ana:E1-final}
\end{equation}

\subsection{第二个指数和与互反相位}

记 $E_2=\sum_{d\asymp D}E_3(d)/d$。固定 $d$，则 $E_3(d)$ 为对 $q$、$w_1,w_2\asymp\mathcal R$、$n_1,n_2\asymp U$ 以及满足 $H\asymp QD\mathcal R^2(\log x)^5/V$ 的 $0<|h|\le H$ 求和，其被加项为
\begin{align}
 &\frac{\Lambda_d(w_1)\overline{\Lambda_d(w_2)}
               \alpha_{n_1}\overline{\alpha_{n_2}}}{qw_1w_2}
 \psi(q/Q)\widehat\psi\left(\frac{hV}{qdw_1w_2}\right)\notag\\
 &\hspace{8mm}\times
 e\left(\frac{\mathfrak{a}h\overline{n_1w_2}}{qdw_1}
          +\frac{\mathfrak{a}h\overline{n_2qdw_1}}{w_2}\right).
 \label{ana:E3-definition}
\end{align}
限制条件为
\[
 \begin{gathered}
 (q,\mathfrak{a})=1,\quad(w_1,\mathfrak{a}w_2)=1,\quad(w_2,\mathfrak{a}qdw_1)=1,\\
 n_1\equiv n_2\pmod{qd},\quad
 (n_1,n_2qdw_1)=(n_2,n_1qdw_2)=1,\quad
 |n_1-n_2|\ge U(\log x)^{-C_{\rm sep}}.
 \end{gathered}
\]
外层 Cauchy 步骤后，不再有独立的有界 $q$ 系数。\eqref{ana:E3-definition} 的分母为 $qw_1w_2$，而因子 $1/d$ 位于 $E_3(d)$ 外部。

精确互反公式给出
\begin{align}
 &e\left(\frac{\mathfrak{a}h\overline{n_1w_2}}{qdw_1}
          +\frac{\mathfrak{a}h\overline{n_2qdw_1}}{w_2}\right)\notag\\
 &\qquad=
 e\left(\frac{\mathfrak{a}h(n_1-n_2)\overline{n_2qdw_1}}{n_1w_2}\right)
 e\left(\frac{\mathfrak{a}h}{qdw_1w_2n_1}\right).
 \label{eq:current-retained-phase}
\end{align}
式中的互素条件保证全部逆元存在。将 $h$ 按正负分开，并把 $q$ 的支撑分成有界个尺度为 $q_0\asymp Q$ 的区间。用 $(q_0,\mathcal R,\mathcal R,U,H)$ 归一化五个变量 $(q,w_1,w_2,n_1,h)$，则最后一个相位的振幅为
\[
 \frac{\mathfrak{a}}{x}\frac Q{q_0}\frac Dd(\log x)^5.
\]
在每个坐标上至多求导一次，即将五维总变差控制为
\begin{equation}
 \bigl(1+(\log x)^5+(|\mathfrak{a}|/x)(\log x)^5\bigr)^5
       =x^{5\kappa+o(1)}.
 \label{ana:phase-variation}
\end{equation}
反复分部求和后只留下锚定矩形和。其端点在估计之前固定，$h$ 系数为单个固定的带符号区间示性函数 $\omega_h$，满足 $|\omega_h|\le1$，不依赖其余任何算术变量。

按 \cite[Lemma 13.2]{MaynardI} 的交叉互素分割划分粗糙互素数对 $(n_1w_2,n_2)$，部分数为 $\exp(O((\log\log x)^4))$。独立展开两份 \eqref{ana:Lambda}，包括分解 $st=d$ 的两种选择。令 $R\asymp v_0$、$S\asymp u_0$，总乘积区间保留为
\[
 \one_{A_i<r_is_i\le B_i}
       =\one_{r_is_i\le B_i}-\one_{r_is_i\le A_i}.
\]
于是它成为两个允许的乘积上截点之差，无需在单变量系数中引入耦合条件。对双线性型，复极化可将两组不同系数化为对角系数，代价至多为对角界的四倍。依次对两组 $R$、两组 $S$ 及两组 $n$ 使用该方法，总常数代价至多为 $4^3$。

下面明确写出所得未加权和。固定 $d$、交叉互素类 $\mathcal C$、端点 $q_0\le q_1\le2q_0$、两个乘积上截点 $B_1,B_2$，以及绝对值不超过一的序列 $\gamma_r,\lambda_s,\alpha_n,\omega_h$。所有这些数据均在其余求和之前固定。设 $\mathscr C$ 为满足
\begin{equation}\label{ana:bare-domain}
\begin{gathered}
 q_0\le q\le q_1,\quad (q,\mathfrak{a})=1,\quad
 r_i\asymp R,\quad s_i\asymp S,\quad n_i\asymp U,\\
 (r_1s_1,\mathfrak{a}r_2s_2)=(r_2s_2,\mathfrak{a}qdr_1s_1)=1,\quad
 r_1s_1\le B_1,\quad r_2s_2\le B_2,\\
 n_1\equiv n_2\pmod{qd},\quad
 (n_1,n_2qdr_1s_1)=(n_2,n_1qdr_2s_2)=1,\\
 (n_1r_2s_2,n_2)\in\mathcal C,\qquad
 |n_1-n_2|\ge U(\log x)^{-C_{\rm sep}}.
\end{gathered}
\end{equation}
的元组 $(q,r_1,r_2,s_1,s_2,n_1,n_2)$ 的集合。定义
\begin{equation}\label{ana:bare-definition}
\mathcal W_{\rm bare}=
 \sum_{\mathscr C}
 \frac{\gamma_{r_1}\lambda_{s_1}
           \overline{\gamma_{r_2}\lambda_{s_2}}
           \alpha_{n_1}\overline{\alpha_{n_2}}}
      {q r_1r_2s_1s_2}
 \sum_{0<|h|\le H}\omega_h
 e\left(\frac{\mathfrak{a}h(n_1-n_2)\overline{n_2r_1s_1dq}}
               {n_1r_2s_2}\right).
\end{equation}
此定义不含光滑互反相位。上述端点分离和极化给出
\begin{equation}\label{ana:E3-bare-reduction}
 |E_3(d)|\ll x^{5\kappa+o(1)}
       \sup|\mathcal W_{\rm bare}|,
\end{equation}
上确界取遍刚才固定的数据。$x^{o(1)}$ 包括交叉互素类及 $d$ 的两次独立分解，但五变量相位损失只出现一次。剩下只需估计 \eqref{ana:bare-definition}。

\paragraph{第一次 Cauchy 不等式。}
令 $f=(n_1-n_2)/(dq)$，它是带符号的非零整数。新变量满足 $0<|f|\le2U/(dq_0)$，且
\begin{equation}\label{ana:f-phase}
 e\left(\frac{\mathfrak{a}h(n_1-n_2)\overline{n_2r_1s_1dq}}
                    {n_1r_2s_2}\right)
 =e\left(\frac{\mathfrak{a}hf\overline{n_2r_1s_1}}
                    {n_1r_2s_2}\right).
\end{equation}
此时仍保留 $q_f=(n_1-n_2)/(df)$ 的整数性以及 $[q_0,q_1]$。$s_1$ 中的乘积截点具有精确有限 Fourier 展开
\[
 \one_{r_1s_1\le B_1}=
 \int_0^1D_{r_1}(\xi)e(s_1\xi)\,d\xi,
 \qquad
 D_{r_1}(\xi)=\sum_{1\le j\le B_1/r_1}e(-j\xi).
\]
存在共同的非负控制函数 $L(\xi)\ll\min(1+B_1/R,1+\|\xi\|^{-1})$，满足 $\int_0^1L(\xi)\,d\xi\ll\log x$ 和 $|D_{r_1}(\xi)|\le L(\xi)$；令 $c_{r_1}(\xi)=D_{r_1}(\xi)/L(\xi)$。固定 $\xi$，记 $\lambda_s'=S\lambda_s e(s\xi)/s$，这些系数在 $s\asymp S$ 上绝对值不超过一。代入 $q=q_f$ 后，用 $\mathscr D$ 表示 \eqref{ana:bare-domain} 中不依赖于 $s_1$ 的条件，则精确 Fourier 被积项为
\begin{align}\label{ana:W2-definition}
 \mathcal W_2(\xi)={}&\frac1S\sum_{\mathscr D}
 \frac{\gamma_{r_1}c_{r_1}(\xi)
              \overline{\gamma_{r_2}\lambda_{s_2}}
              \alpha_{n_1}\overline{\alpha_{n_2}}}
      {q_f r_1r_2s_2}\notag\\
 &\quad\times
 \sum_{\substack{s_1\asymp S\\(s_1,\mathfrak{a}n_1r_2s_2)=1}}
 \lambda_{s_1}'\sum_{0<|h|\le H}\omega_h
 e\left(\frac{\mathfrak{a}hf\overline{n_2r_1s_1}}{n_1r_2s_2}\right).
\end{align}
特别地，$|\mathcal W_{\rm bare}|\ll(\log x)\sup_\xi|\mathcal W_2(\xi)|$。对每个 $f$，满足 $n_1\equiv n_2\pmod{df}$ 的数对 $n_1,n_2\asymp U$ 有 $O(U^2/(d|f|)+U)$ 个。先对 $f$ 再对 $r_1,r_2,s_2$ 求和，外层元组共有 $O(U^2R^2S\log x)$ 个。由于 \eqref{ana:W2-definition} 中的外部系数绝对值至多为 $O(1/(q_0R^2S^2))$，Cauchy 不等式给出
\begin{equation}\label{ana:first-Cauchy}
 |\mathcal W_2(\xi)|^2
 \ll\frac{U^2\log x}{q_0^2R^2S^3}\,\mathcal W_3,
\end{equation}
其中以下非负和保留同余条件，去掉其他外层限制：
\begin{align}\label{ana:W3-definition}
 \mathcal W_3={}&
 \sum_{0<|f|\le2U/(dq_0)}
 \sum_{\substack{n_1,n_2\asymp U\\n_1\equiv n_2\pmod{df}}}
 \sum_{r_1,r_2\asymp R}
 \sum_{\substack{s_2\asymp S\\(n_2r_1,n_1r_2s_2)=1}}\notag\\
 &\quad\times
 \left|\sum_{\substack{s_1\asymp S\\(s_1,\mathfrak{a}n_1r_2s_2)=1}}
 \lambda_{s_1}'\sum_{0<|h|\le H}\omega_h
 e\left(\frac{\mathfrak{a}hf\overline{n_2r_1s_1}}{n_1r_2s_2}\right)
 \right|^2.
\end{align}
所去掉的每个限制都位于内层平方和之外。$\mathcal C$ 的交叉互素性质，保证在移去原有类别限制之前，式中互素条件已成立。

\paragraph{补全非负平方。}
现在可去掉 \eqref{ana:W3-definition} 中的同余条件并扩大 $|f|\le2U/q_0$，因为所有被加项均非负。令 $b=n_2r_1$、$c=n_1r_2s_2$，其重数至多为 $\tau(b)\tau_3(c)$，可在求和前控制。将 $b,c$、$|f|$ 分入二进区间。每个块上有 $B\ll UR$、$C\ll URS$、$F\ll U/q_0$，取固定光滑非负控制函数 $g(b,c)=\psi(b/B)\psi(c/C)$，得
\begin{equation}\label{ana:second-Cauchy}
 \mathcal W_3\ll x^\epsilon\sum_{B,C,F,\pm}\mathcal W_4,
\end{equation}
其中 $f$ 的符号吸收到 $\mathfrak{a}$ 中，且
\begin{equation}\label{ana:W4-definition}
 \mathcal W_4=
 \sum_{\substack{b,c\\(b,c)=1}}g(b,c)
 \sum_{f\asymp F}
 \left|\sum_{\substack{s\asymp S\\(s,\mathfrak{a}c)=1}}
 \lambda_s'\sum_{0<|h|\le H}\omega_h
 e\left(\frac{\mathfrak{a}hf\overline{bs}}c\right)\right|^2.
\end{equation}
二进块数只有对数量级。结合 \eqref{ana:first-Cauchy}、\eqref{ana:second-Cauchy} 并使用 $q_0\asymp Q$，外部因子可写为
\begin{equation}\label{ana:bare-Cauchy-factor}
 |\mathcal W_{\rm bare}|
 \ll x^\epsilon\frac{U}{QRS^{3/2}}
       \sup_{B,C,F}\mathcal W_4^{1/2}.
\end{equation}

展开 \eqref{ana:W4-definition}，必要时交换两个 $s$ 标签，得到
\begin{equation}\label{ana:W4-expansion}
 \mathcal W_4=
 \sum_{\substack{s_1,s_2\asymp S\\(s_1s_2,\mathfrak{a})=1}}
 \sum_{0<|h_1|,|h_2|\le H}\lambda_{s_2}'\overline{\lambda_{s_1}'}
       \omega_{h_1}\overline{\omega_{h_2}}
 \sum_{f\asymp F}\sum_{\substack{b,c\\(c,bs_1s_2)=1}}
 g(b,c)e\left(\frac{\mathfrak{a}f(h_1s_1-h_2s_2)\overline{bs_1s_2}}c\right).
\end{equation}
满足 $h_1s_1=h_2s_2$ 的项总绝对值至多为 $x^\epsilon HSBFC$。对其余项，令 $z=s_1s_2$、$y=f(h_1s_1-h_2s_2)\ne0$，并定义系数
\begin{equation}\label{ana:beta-definition}
 \beta(y,z)=
 \sum_{\substack{s_1s_2=z,\ s_i\asymp S\\
       f\asymp F,\ 0<|h_i|\le H\\
       f(h_1s_1-h_2s_2)=y}}
 \lambda_{s_2}'\overline{\lambda_{s_1}'}
       \omega_{h_1}\overline{\omega_{h_2}}.
\end{equation}
保留 $(z,\mathfrak{a})=1$。变量 $b,c$ 不出现在 \eqref{ana:beta-definition} 中；涉及它们的限制仍在核内。因此非零部分恰为
\begin{equation}\label{ana:off-diagonal-kernel}
 \sum_{\substack{y\ne0,z\\(\mathfrak{a},z)=1}}\beta(y,z)
 \sum_{\substack{b,c\\(c,bz)=1}}g(b,c)
                         e(\mathfrak{a}y\overline{zb}/c).
\end{equation}
特别地，对下一引理中的非负元组计数有 $|\beta(y,z)|\le b(z,y)$。在 $y,z$ 上再作任意二进限制后，此控制仍成立，且不依赖于尚未引入的 Poisson 频率。

\begin{lemma}\label{ana:multiplicity}
令 $b(z,y)$ 计数满足 $z=s_1s_2$、$y=f(h_1s_1-h_2s_2)\ne0$ 的元组 $s_i\asymp S$、$f\asymp F$、$0<|h_i|\le H$。则在具有多项式上界的尺度上，对每个 $\nu\ge1$ 有
\begin{equation}
 \sum_z\sum_{u\ne0}b(z,\nu u)^2
 \ll x^\epsilon(1+H/S)\frac{(HS)^2F}{\nu}
 \label{ana:multiplicity-bound}
\end{equation}
当 $\nu\ge8FHS$ 时左端为零。
\end{lemma}

\begin{proof}
固定 $(z,y)$ 时，分解 $(s_1,s_2,f,k)$ 至多有 $\tau(z)\tau(|y|)$ 个，且 $k=h_1s_1-h_2s_2\ne0$。因此，由 Cauchy 不等式，只需对具有相同 $(s_1,s_2,k)$ 的 $h$ 数对碰撞计数，再作一次 $f$ 计数。按 $c\mid\nu$、$\nu/c\mid k$ 分组，$f$ 的个数至多为 $2F/c$。相同 $h$ 的数对贡献为 $\ll x^\epsilon(HS)^2F/\nu$。对不同数对，令 $g=(s_1,s_2)$、$s_1=gu$、$s_2=gv$、$(u,v)=1$，其差为 $(h_1-h_1',h_2-h_2')=(jv,ju)$，且 $0<|j|\le2Hg/S$。若 $b=(g,\nu/c)$，则 $k/g$ 的条件为被 $\nu/(cb)$ 整除。从零起计数非零倍数而不增加单位项，至多得到 $O(HScb/(g\nu))$ 个值。分解其余非零整数 $h_2v$ 只损失除数因子，因此不同数对的贡献为
\[
 \ll x^\epsilon\frac{H^3SF}{\nu}
             \sum_{g\le2S}\frac{(g,\nu)}g
 \ll x^{2\epsilon}\frac{H^3SF}{\nu}.
\]
最后一个估计来自 $\sum_{g\le2S}(g,\nu)/g\le\tau(\nu)(1+\log(2S))$。重新命名 $2\epsilon$ 即得结论。支撑界由 $0<|y|<8FHS$ 得到。
\end{proof}

在当前应用中，$H/S\le x^{o(1)}(Q/U)D_0<1$，其中 $D_0=U^2R^2S/x$，所以因子 $1+H/S$ 有界。

\subsection{互反核与谱估计}

我们估计互反核时，令系数独立于之后的 Poisson 频率。由此可在每个谱层级使用同一频率区间，满足异常谱层级平均估计的要求。

本小节用 $Y$ 表示系数长度，$Z$ 表示筛级长度，$P$ 表示非零 Poisson 频率块的长度。原卷积长度仍为 $U,V$。设 $\beta(u,z)$ 的支撑满足 $u\asymp Y$、$z\asymp Z$、$(\mathfrak{a},z)=1$，令
\begin{equation}
 \mathcal K=\sum_{u,z}\beta(u,z)
   \sum_{\substack{c,b\\(c,bz)=1}}
   \psi(c/C)\psi(b/D)e(\mathfrak{a}u\overline{zb}/c).
 \label{ana:kernel}
\end{equation}
$u$ 的正负分别处理。定义 $B^2=\sum_{u,z}|\beta(u,z)|^2$ 和 $B_\nu^2=\sum_{u,z}|\beta(\nu u,z)|^2$；再加入互素示性函数或模为一的相位，不会增大这些范数。

对 $b$ 作 Poisson 求和，得到精确恒等式
\[
 \sum_{(b,c)=1}\psi(b/D)e(\mathfrak{a}u\overline z\,\overline b/c)
 =\frac Dc\sum_{m\in\mathbb Z}\widehat\psi(mD/c)
                   S(m\overline z,\mathfrak{a}u;c).
\]
因 $|S(0,\mathfrak{a}u;c)|\le(\mathfrak{a}u,c)$、$\sum_{c\asymp C}(\mathfrak{a}u,c)\ll C\tau(|\mathfrak{a}u|)$，零频贡献至多为 $x^\epsilon D\sqrt{YZ}\,B$，其平方为 $x^\epsilon D^2YZB^2$。对非零 $m$，在共同截点 $(C/D)x^\epsilon$ 处截断，再划分为共同的带符号二进区间。取固定且充分高的导数阶数，普通光滑 Fourier 衰减即可控制被删除的尾部。

在实际区间 $m\asymp P$ 上，令
\[
 Y_0=\frac{4\pi\sqrt{|\mathfrak{a}|PY}}{C\sqrt Z},\quad
 v=m/P,\quad w=u/Y,\quad \zeta=z/Z,\quad
 t=\frac{4\pi\sqrt{|\mathfrak{a}|mu}}{c\sqrt z\,Y_0}.
\]
则 $c/C=\sqrt{vw/\zeta}/t$。此变量替换从归一化振幅中移去 $\mathfrak{a}$。在固定紧盒上作 Fourier 展开，将振幅分离为一个 Bessel 测试函数、共同因子 $\one_I(m)e(m\omega)$ 以及任意复 $u,z$ 系数。取 Fourier 衰减阶数高于固定的 Bessel 导数阶数，即可使系数范数可求和。因此 $m$ 区间 $I$ 及相位 $\omega$ 在每个筛级 $z$ 上均相同。

尖点转换与旧形式族操作的精确定义见附录~\ref{app:classical}。它们把 $u$ 系数变为
\begin{equation}
 \beta_\nu(n,z)=\one_{(n,\mathfrak{a})=1}\beta(\nu n,z)e(-n/z),
 \qquad \nu\mid \mathfrak{a}^\infty,
 \label{ana:packet-coefficient}
\end{equation}
代价为 $(|\mathfrak{a}|\nu)^{\vartheta/2}$，并保留共同的 $m$ 泛函。令
\[
 P_0=\frac{PC^4Z}{|\mathfrak{a}|PY+C^2Z},\qquad T_0=\frac{|\mathfrak{a}|PY}{C^2}.
\]
正规谱贡献与异常谱估计给出
\begin{align}
 |\mathcal L_P|^2\ll x^\epsilon P_0\bigg\{
 &\sum_{\nu\mid \mathfrak{a}^\infty}(|\mathfrak{a}|\nu)^\vartheta
       (Y/\nu+Z+T_0)(P+Z+T_0)B_\nu^2\notag\\
 &+(C^2Z)^\vartheta(Y+Z)^{1-\vartheta}
                    (P+Z)^{1-2\vartheta}B^2\bigg\}.
 \label{ana:spectral-polynomial}
\end{align}
其中 $\mathcal L_P$ 为分离后的 Kuznetsov 和，此时尚未恢复 Poisson 因子 $D/C$。

下面说明得到这一多项式的谱求和。此步骤中，正规谱范数始终为非负直和范数 \eqref{ana:regular-norm}：Maass 原子的权为 $1/\cosh(\pi t_j)$；按指定 Fourier 归一化后，全纯原子的权为 $1/(2\pi)$；连续部分的测度为 $dt/\pi$。先在同一测度的各个分量内部施加 Cauchy 不等式，再将三个分量相加。固定 $z,\nu$，以 $A(t),B(t)$ 分别表示 $d\mu_z$ 上的 $n$ 端和 $m$ 端谱系数数组，其条目分别为 $\sum_n\beta_\nu(n,z)R_{t,1}(n)$ 与 $\sum_{m\in I}e(m\omega)R_{t,\infty}(m)$。此处 $t$ 是正规谱直和中的条目指标，包含分量和基向量标签；$R_{t,\mathfrak c}$ 表示 \eqref{ana:regular-norm} 中相应的 Fourier 系数。令 $A_0,B_0\ge0$ 为范数界的幅度常数，并取 $u_0,v_0\ge0$。若两个累计正规谱平方范数分别不超过 $(T^2+u_0)A_0^2$ 和 $(T^2+v_0)B_0^2$，则 Bessel 界
\[
 |\mathcal H(t)|\ll\frac{1+|\log Y_0|}{1+Y_0},\qquad
 |\mathcal H(t)|\ll |t|^{-5/2}+Y_0|t|^{-3}
 \quad(|t|\ge\max(1,2Y_0))
\]
给出
\[
 \int|\mathcal H(t)A(t)B(t)|\,d\mu_z(t)
 \ll \frac{1+|\log Y_0|}{1+Y_0}
 \sqrt{(1+Y_0^2+u_0)(1+Y_0^2+v_0)}A_0B_0.
\]
事实上，在 $\max(1,2Y_0)$ 以上，累计范数之积在第 $j$ 个壳层至多增长 $2^{2j}$，留下可求和因子 $2^{-j/2}+2^{-j}$。筛级上的 Cauchy 因子为 $\sqrt Z$。取 $L=Y/\nu$ 即得 \eqref{ana:spectral-polynomial} 的第一项，其中包括旧形式族代价。

对异常形式，令 $\tau_j=2\sqrt{1/4-\lambda_j}\le\vartheta$、$X_0=1+Z/L$、$X=1/Y_0$，使用
\[
 1+X^{\tau_j}\le
 2(1+X/\sqrt{X_0})^{\tau_j}(1+X_0)^{\tau_j/2}.
\]
固定筛级的 $n$ 范数为 $\ll x^\epsilon(1+L/Z)B_\nu^2$，共同区间的 $m$ 范数为
\[
 \ll x^\epsilon P(P+Z)
 \left[1+\left\{
  \frac{P(1+X/\sqrt{X_0})^2}{(P+Z)^2}\right\}^{\vartheta}\right].
\]
将这些界与旧形式族代价的平方相乘，得到可吸收入正规谱表达式的项以及 \eqref{ana:spectral-polynomial} 的第二项。仅在后一项中，两个范数相乘后因子 $|\mathfrak{a}|\nu$ 才被消去。同一论证也适用于高度为 $k-1$ 的全纯部分，以及 \eqref{ana:regular-norm} 中实高度为 $t$ 的连续部分；整个壳层求和过程中，其测度均不变。在异常谱上，因 $|\operatorname{Im}t_j|\le7/64$，有 $1/\cosh(\pi t_j)=1/\cos(\pi|\operatorname{Im}t_j|)\asymp1$，故以上所用未加权异常谱估计只相差一个绝对常数。

恢复 $D/C$ 后，由初等不等式
\[
 \frac{Z(L+Z+T_0)(P+Z+T_0)}{Z+T_0}
       \le(L+Z)(P+Z)+ZT_0
\]
及 $P\le C/D$，得到
\[
 |\mathcal K|^2\ll x^\epsilon\mathcal J^2,
\]
其中完整的非负多项式为
\begin{align}
 \mathcal J^2={}&\sum_{\nu\mid \mathfrak{a}^\infty}(|\mathfrak{a}|\nu)^\vartheta
 \bigl\{C(Y/\nu+Z)(C+DZ)+|\mathfrak{a}|YZ\bigr\}B_\nu^2
 \notag\\
 &+\bigl\{
 C(CD\sqrt Z)^{2\vartheta}(Y+Z)^{1-\vartheta}
                         (C+DZ)^{1-2\vartheta}+D^2YZ\bigr\}B^2.
 \label{ana:returned-polynomial}
\end{align}
通过缩小最初选择的 $\epsilon$，即可兼容实际截点 $P\le x^\epsilon C/D$。若 $C/D<1/2$，只对实际非空区间应用估计；实数比较值 $C/D$ 不被用作并不存在的整数区间的长度。

\paragraph{回到卷积尺度。}
先在 \eqref{ana:off-diagonal-kernel} 的实际二进长度上应用 \eqref{ana:returned-polynomial}。因 $0\le\vartheta\le1/2$，其系数均非负，且对四个长度单调递增，故随后可将这些长度替换为
\begin{equation}\label{ana:scale-substitution}
 D=B_*=UR,\quad C=C_*=URS,\quad Z=S^2,\quad
 Y=\frac{UR^2S^3}{V}.
\end{equation}
相应计数参数满足 $F\le U/Q$、$H\le x^\epsilon H_0$，其中 $H_0=QR^2S^2/V$。先将输入指数选得更小，这些可调幂次损失即可纳入最终的 $x^\epsilon$。令
\[
 D_0=\frac{U^2R^2S}{x}=\frac YZ<1.
\]
则 $B_*Z\ge C_*$、$B_*^2<x$、$C_*B_*Z=xY$。由引理~\ref{ana:multiplicity} 及 $H/S\le x^\epsilon(Q/U)D_0<1$，有
\begin{equation}\label{ana:norm-insertion}
 B_\nu^2\le x^\epsilon\frac{(HS)^2F}{\nu},\qquad
 B^2\le x^\epsilon(HS)^2F.
\end{equation}
辅助除数和满足
\[
 \sum_{\nu\mid\mathfrak a^\infty}\nu^{\vartheta-1}
 =\prod_{p\mid\mathfrak a}(1-p^{-(1-\vartheta)})^{-1}
 =x^{o(1)}.
\]
把素因子在 $\log x$ 处分开，可知其对数为 $O((\log x)^\vartheta)$，并对具有多项式上界的 $\mathfrak a$ 一致。因此 \eqref{ana:returned-polynomial} 给出
\begin{equation}\label{ana:J-reduced}
 \mathcal J^2\ll x^\epsilon(HS)^2F
 \left\{
 |\mathfrak a|^\vartheta(C_*B_*Z+|\mathfrak a|Y)Z
 +C_*^{1+2\vartheta}B_*Z^{2-2\vartheta}
 +B_*^2YZ\right\}.
\end{equation}
为明确起见，右端四个单项式乘以
\[
 \frac{(H_0S)^2(U/Q)}{U^4R^4S^6}
                  =\frac{Q}{U^3V^2}:
\]
后所得比值如下：
\begin{equation}\label{ana:ratio-table}
\begin{array}{c|c}
 \text{单项式}&\text{归一化贡献}\\ \hline
 |\mathfrak a|^\vartheta C_*B_*Z^2
     &D_1=|\mathfrak a|^\vartheta UQR^2S^5/x^2\\
 |\mathfrak a|^{1+\vartheta}YZ
     &(|\mathfrak a|/x)D_1\\
 C_*^{1+2\vartheta}B_*Z^{2-2\vartheta}
     &D_2=U^{1+2\vartheta}QR^{2+2\vartheta}S^{5-2\vartheta}/x^2\\
 B_*^2YZ&QR^4S^5/V^3\le D_1.
\end{array}
\end{equation}
最后一项比较来自 $B_*^2<x$、$|\mathfrak a|\ge1$。\eqref{ana:W4-expansion} 中的零差贡献满足
\[
 \left(\frac{U}{QRS^{3/2}}\right)^2H_0SB_*\frac UQ C_*
       =\left(\frac{U^2}{Q}\right)^2D_0.
\]
对非零部分，\eqref{ana:off-diagonal-kernel} 给出 $|\mathcal W_{4,\ne0}|\ll x^\epsilon\mathcal J$。将此平方根与 \eqref{ana:bare-Cauchy-factor} 中的平方根结合，明确得到
\[
 \left(\frac{U}{QRS^{3/2}}\right)^4
 \frac{\mathcal J^2}{(U^2/Q)^4}
 =\frac{\mathcal J^2}{U^4R^4S^6}
 \ll x^\epsilon\bigl\{(1+|\mathfrak a|/x)D_1+D_2\bigr\}.
\]
两部分共同证明
\begin{align}
 |\mathcal W_{\rm bare}|&\ll\frac{U^2}{Q}x^\epsilon
   \left\{D_0^{1/2}
    +\bigl[(1+|\mathfrak{a}|/x)D_1+D_2\bigr]^{1/4}\right\},
 \label{eq:current-bare-bound}\\
 D_0&=\frac{U^2R^2S}{x},\qquad
 D_1=\frac{|\mathfrak{a}|^\vartheta UQR^2S^5}{x^2},\qquad
 D_2=\frac{U^{1+2\vartheta}QR^{2+2\vartheta}S^{5-2\vartheta}}{x^2}.
 \notag
\end{align}

\subsection{卷积估计的完成}

尺度假设在每个小因子块上给出 $D_0\ll x^{-8\sigma}$、$D_1,D_2\ll x^{-15\sigma}$。结合 $|\mathfrak{a}|/x\le x^\kappa$ 以及只出现一次的变差代价 \eqref{ana:phase-variation}，\eqref{eq:current-bare-bound} 给出
\[
 |E_3(d)|\ll\frac{U^2}{Q}
 \left(x^{-4\sigma+5\kappa+\epsilon}
       +x^{-15\sigma/4+21\kappa/4+\epsilon}\right).
\]
全部可选的 epsilon 损失、分割产生的趋零次幂，以及一个额外归一化因子 $D$，均可纳入 $x^{\sigma/10}$。由于 $\kappa\le\sigma/100$，剩余两个节省量都超过 $\sigma/2$。故对 $d$ 作调和求和，得到
\begin{equation}
 |E_2|\ll\frac{U^2}{DQ}x^{-\sigma/2}.
 \label{ana:E2-final}
\end{equation}
将 \eqref{ana:E1-final}、\eqref{ana:E2-final} 代入 \eqref{ana:dispersion-bound}，指数和贡献为 $O(x^{1-\sigma/4}(\log x)^C)$。先前恢复单零项的代价为 $O(x^{1-\sigma/6+o(1)})$，仍有幂次节省。选择辅助 Siegel--Walfisz 精度和对数精度，以吸收有限的二进选择、固定对数数量的块，以及 $\beta$ 的归一化。被略去的光滑部分由 \eqref{ana:smooth-tail} 处理。这就证明了 \eqref{eq:current-medium-conclusion}。

\section{谱输入的归一化}
\label{app:classical}

这些输入用于 \eqref{ana:spectral-polynomial}；packet 系数与共同频率区间见 \eqref{ana:packet-coefficient} 及该估计之前的讨论。

下面明确归一化，以确定谱 Cauchy 不等式中使用的正测度。常规谱与异常谱分别在各自范围内处理。

下面明确列出前面证明使用的经典结论。其特征为平凡特征，筛级为 $z$，尖点为 $\infty$、$1$，且 $(\mathfrak{a},z)=1$。用 $R(n)=\sqrt{|n|}\rho(n)$ 表示 Deshouillers--Iwaniec 归一化的 Fourier 系数。尖点伸缩矩阵为
\[
 \sigma_\infty=I,\qquad
 \sigma_1=\begin{pmatrix}\sqrt z&0\\\sqrt z&1/\sqrt z\end{pmatrix}.
\]
\cite[Lemma 4.3]{Drappeau} 中的尖点转换引入 $e(\mathfrak{a}y/z)$，因此普通 Kloosterman 系数带有 $e(-\mathfrak{a}y/z)$。

对 $y=\nu n$、$\nu\mid \mathfrak{a}^\infty$、$(n,\mathfrak{a})=1$，\cite[Lemmas 3.1 及 3.2]{ABL} 中互素旧形式族的恒等式在每个归一化的旧形式伸缩上给出
\[
 R_v(\mathfrak{a}\nu n)=\lambda_v(\mathfrak{a}\nu)
                   e\bigl(n(\mathfrak{a}\nu-1)/z\bigr)R_v(n).
\]
乘以 $e(-\mathfrak{a}\nu n/z)$ 后恰好留下 $e(-n/z)$。计入除数损失后，\cite[(3.3)]{ABL} 中族系数的 Frobenius 范数为 $z^\epsilon$。因此，族内的 Cauchy 不等式以及 Kim--Sarnak 指数界~\cite{KimSarnak} 带来代价 $(|\mathfrak{a}|\nu)^{7/64+\epsilon}$，并给出共同系数 \eqref{ana:packet-coefficient}。不假设所有伸缩具有同一特征值。\cite[Lemma 3.2]{ABL} 中的全纯及连续二次型使用相应恒等式，取 $\vartheta=0$。连续项按形式 Plancherel 测度积分，不把单个 Eisenstein 级数当作平方可积函数取范数；另见 \cite[Proposition 8.2]{Young}。

为完整起见，下面明确给出每次正规谱 Cauchy 不等式所用的非负范数。设 $\mathfrak c$ 为任一尖点，采用 \cite[Proposition E]{Pascadi} 的记号，写作
\begin{align*}
 R^M_{j,\mathfrak c}(n)&=
       \rho^{\rm DI}_{j,\mathfrak c}(n)
       =\sqrt{|n|}\rho^{\rm ABL}_{j,\mathfrak c}(n),\\
 R^H_{jk,\mathfrak c}(n)&=
       \left(\frac{(k-1)!}{(4\pi)^{k-1}}\right)^{1/2}
       n^{-(k-1)/2}\psi_{jk,\mathfrak c}(n),\\
 R^E_{\mathfrak d,\mathfrak c}(n,t)&=
       n^{it}\varphi_{\mathfrak d\mathfrak c}
                       (n;\tfrac12+it).
\end{align*}
旧形式族公式中的符号 $R_j$ 表示 $R_j^M$。对支撑在 $n\asymp L$、$L\ge1/2$ 上的 $T\ge1$、$b$，定义
\begin{align}\label{ana:regular-norm}
 \mathfrak Q_{z,\mathfrak c}(T;b)
 ={}&\sum_{\substack{t_j\in\mathbb R\\|t_j|\le T}}
 \frac1{\cosh(\pi t_j)}
 \left|\sum_{n\asymp L}b_nR^M_{j,\mathfrak c}(n)\right|^2\notag\\
 &+\frac1{2\pi}\sum_{\substack{k\in2\mathbb Z_{>0}\\k-1\le T}}
      \sum_j\left|\sum_{n\asymp L}b_nR^H_{jk,\mathfrak c}(n)\right|^2
       \notag\\
 &+\frac1\pi\sum_{\mathfrak d}\int_{-T}^T
      \left|\sum_{n\asymp L}b_n
                 R^E_{\mathfrak d,\mathfrak c}(n,t)\right|^2dt.
\end{align}
这就是前文使用的直和测度 $d\mu_z$。常数 $1/(2\pi)$、$1/\pi$ 来自 Kuznetsov 公式；因子 $1/\cosh(\pi t_j)$ 属于 Maass 测度，不会因从 $\rho^{\rm ABL}$ 改为 $R^M$ 而消失。由 \cite[Proposition 4.7]{Drappeau}、\cite[Lemma H]{Pascadi}，有
\begin{equation}\label{ana:regular-large-sieve}
 \mathfrak Q_{z,\mathfrak c}(T;b)
 \ll (zLT)^\epsilon\left(T^2+\frac Lz\right)\sum_n|b_n|^2.
\end{equation}
从 $k\le T$ 改为 $k-1\le T$ 只损失固定因子。此约定中的连续项测度为 $dt/\pi$，与 \cite[(3.27)]{Pascadi} 相同，不再附加双曲因子。配合该范数的测试函数界为 \cite[Lemma D]{Pascadi}。对异号 Kuznetsov 和，全纯项不存在，同一非负范数仍给出上界。

在异常谱上令 $\tau_j=2\sqrt{1/4-\lambda_j}\le7/32$。尖点 $1$ 处的固定筛级估计 \cite[Theorem A]{Pascadi} 给出
\begin{equation}
 \sum_{\rm exc}X^{\tau_j}
       \left|\sum_{n\asymp L}b_nR_{j,1}(n)\right|^2
 \ll(zL)^\epsilon(1+L/z)\sum_n|b_n|^2,
 \qquad X\ll\max(1,z/L).
 \label{ana:exceptional-fixed}
\end{equation}
系数可依赖 $z$，但不依赖谱高度。筛级平均估计 \cite[Theorem J]{Pascadi} 给出
\begin{equation}
 \sum_{z\asymp Z}\sum_{\rm exc}X^{\tau_j}
       \left|\sum_{m\in I}e(m\omega)R_{j,\infty}(m)\right|^2
 \ll(ZV)^\epsilon(Z+V)V,
 \quad X\ll\max(V,Z^2/V),
 \label{ana:exceptional-average}
\end{equation}
其中 $I\subset(V,2V]$、$\omega$ 在每个筛级均相同。对部分区间 $I$，选择整数 $P>4V$，将 $\one_I(m)=\sum_{k=0}^{P-1}c_ke(km/P)$ 按模 $P$ 的剩余类展开，其系数范数为 $\sum|c_k|\ll\log(2+V)$。在以 $(z,j)$ 为指标的 Hilbert 范数中使用 Minkowski 不等式，再以相位 $\omega+k/P$ 应用定理 J，即得 \eqref{ana:exceptional-average}；仅有一个对数平方损失，可吸收入 $(ZV)^\epsilon$。这一论证要求每个筛级使用相同区间。对更大的 $X$，由界 $\tau_j\le\vartheta$ 得到延拓因子 $1+(VX/(V+Z)^2)^\vartheta$。

变换后的估计适用于每个长度至少为 $1/2$ 的非空区间，包括含有 $n=1$ 的端点区间。$m$ 的区间对所有层级共同，$n$ 的系数不依赖谱高度。每次谱包操作均保留 $(\mathfrak a,z)=1$。这正是前述谱估计所使用的区间与系数限制。
\section{有限筛系数的分解}
\label{app:coefficients}

本附录证明，实际的上、下界筛系数具有分布估计所需的分解。全文以 $N$ 表示目标偶数。有限分割、筛级和逼近精度均在 $N$ 趋于无穷之前选择，算术系数本身则可以依赖 $N$。记 $\omega_{\mathrm{pf}}(n)$ 为 $n$ 的不同素因子个数。

\subsection{有限决策规则与剩余平方条件}

记
\[
 \rho_0=\frac{19}{96},\qquad \mu_0=\frac{25}{128},\qquad
 \Lambda_{\rm cap}=\frac{153}{256},\qquad \Theta=\frac{299}{500}.
\]
$\Theta$ 控制完整支撑，$\Lambda_{\rm cap}$ 控制高阶筛级。有限决策规则在有限多个素数前缀上分别指定筛级 $N^\lambda$；祖先决策在子节点细分之前固定。对递减指数 $t_1,\ldots,t_j$，令 $S_j=t_1+\cdots+t_j$。辅助筛级 $N^\ell$ 上的完整平方条件为
\begin{equation}
 S_{j-1}+2t_j\le\ell.
 \label{coef:full-squares}
\end{equation}
若 $I$ 是带标记位置的选定子集，令 $b_I=\sum_{i\in I}t_i$。除 \eqref{coef:full-squares} 外，还保留所有条件
\begin{equation}
 b_I+\sum_{\substack{i<j\\i\notin I}}t_i+2t_j\le\ell
 \qquad(j\notin I).
 \label{coef:selected-squares}
\end{equation}
在最后一个选定指标之后，对 $j$，\eqref{coef:selected-squares} 的左端恰为 $S_{j-1}+2t_j$。因此，只要完整条件保留，接入更小素数就会保留选定条件。一个可用的立方条件还蕴含下一个平方条件，因为
\[
 S_j+2t_{j+1}\le S_{j-1}+3t_j.
\]
这一观察对普通筛树的两种奇偶性都适用：上界首次退出为奇数阶，下界首次退出为偶数阶。

实际筛级 $\lambda$ 和辅助覆盖筛级 $\ell$ 的用途不同。系数及其密度使用 $\lambda$，而分解可以采用略大的 $\ell$；获得该覆盖时不改变 $\lambda$。以下给出三因子备选所需的选择。令
\[
 A(t)=\frac{1+t}{2},\qquad W(z)=\Lambda_{\rm cap}-\frac38z.
\]
在粗盒上，设 $\Gamma=\inf\max(A,W)$、$\lambda=\Gamma-2e$，其中变动预留量满足 $\chi<e/10$。若 $\inf A\ge\lambda+e/2$，取 $\ell=\lambda+e/4$；否则有 $\sup A<\lambda+11e/20$，故整个盒上都有 $W\ge\Gamma$，此时取 $\ell=\lambda+e$。这样在下边界以上至少留有 $9e/20$，在上边界以下至少留有 $e$。其余初始化预留量包括
\[
 1-\Lambda_{\rm cap}-2\mu_0=\frac3{256},\qquad
 \frac{525}{8704}+\frac{49z}{272}>0.
\]
高阶变更限制在各条件均有正余量的紧区域内，再令精细端点扩张小于该余量。

在普通下界起点 $t_1\le {h_{\rm s}}=19/125<\mu_0$，端点可用条件满足
\[
 \max\left(\frac{1+t_1}{2},\Lambda_{\rm cap}-\frac38t_1\right)
 \ge\frac{249}{448}.
\]
因此，比 $249/448$ 低一个固定余量的筛级具有严格的覆盖区间。第四阶退出的恢复，必须先通过其自身立方条件及全部完整和选定平方条件。对 $\mu_0$ 以上的上界前缀，函数
\[
 v(x)=\min\left(\frac{1+x}{2},\frac{57}{64}-\frac32x\right)
\]
在 $x<\rho_0$ 时满足 $v(x)>3x$。选取 $y<x$ 给出 $2x+y<3x$；选取 $y+z<x$ 给出 $2x+y+z<3x$。这就是两个附加备选的初始平方界。

在带标记状态中，记 $S=\log_N B$、$L=\lambda-S$、$\alpha=(1-S)/L\ge2$。指数为 $Lu$ 的单素数将剩余长度变为 $L(1-u)$；指数为 $Lu,Lv$ 的素数对将其变为 $L(1-u-v)$。两种情形下 $\alpha$ 都增大。条件 $2u<1$，或 $2v<1$ 与 $v+2u<1$，恰好是新增的完整平方条件。选定条件随后由 \eqref{coef:selected-squares} 推出，父节点的素数上限也保留。有限构造至多含 $440$ 次基本转移。对某个 $\ell_0>0$，其剩余指数至少为 $L_{\min}=\ell_0 101^{-440}$，接入素数的指数至少为 $L_{\min}/6$。特别地，全部换位叶节点具有共同的正素数截点。

\subsection{有限逼近与严格余量}

对全部带标记查询共同选择有限分割。在第三个上界退出处，被略去的末素数尾部为
\[
 O\!\left(\eta\left(1+\log^2\frac{\rho_0}{\eta}\right)\right),
 \qquad\eta\downarrow0.
\]
在 $r=3,4,5$ 阶下界退出处，相应界为 $O(\eta(1+\log^{r-1}({h_{\rm s}}/\eta)))$。更高阶变更具有正的末素数截点；特别地，八阶区域满足 $t_8>53/13824$。在其余紧区域上，可用任意小质量裁去真正的仿射边界。若某区域上的标签恒为继承，则保留该区域，不把它当作边界。先固定粗层决策，再细分子节点，使计数和密度使用同一有限规则。

对于第三阶包络，粗参数 $e,\chi$ 给出 $0\le L_{\mathrm{end}}-\lambda\le4e+6\chi$。固定有限前缀规则及其正余量后，可用性函数在有限多个紧盒上连续，故其严格余量具有正的最小值。因此，存在 $\kappa_*>0$，使得当 $\log_xN\in[1,1+\kappa_*]$ 且端点扩张充分小时，全部不等式仍成立。这里的余量是在除去极限边界区域之后，针对固定紧区域取得的。

\subsection{三因子表示}

由于分解必须对完整系数成立，下面说明端点构造。设 $x$ 为完整局部尺度，$q_0=\log_x N$、$\vartheta=7/32$，$\alpha_0\ge q_0$ 为局部可用条件中的剩余类参数。对固定的小 $\zeta>0$，令
\[
 \theta_{\rm cov}=q_0\ell+2\zeta,\qquad D=x^{\theta_{\rm cov}-2\zeta},
 \qquad A=Ux^{-\zeta},
\]
其中 $U$ 是选定的中等长度。定义
\begin{align*}
 v_0&=2\theta_{\rm cov}-1,\\
 u_0&=\min\left\{
 \frac{2-(2+2\vartheta)\theta_{\rm cov}}{3-4\vartheta},
 \frac{2-2\theta_{\rm cov}-\alpha_0\vartheta}{3}\right\}.
\end{align*}
本小节各位置的指数以及 $b_I$ 都以 $x$ 为底，即为先前以 $N$ 为底的值乘以 $q_0$。因此，完整与选定平方条件中的覆盖指数为 $q_0\ell$。选定子集的可用条件为 $v_0\le b_I\le u_0$，并对有序上端点施加 \eqref{coef:full-squares}--\eqref{coef:selected-squares}。重标度后保留该条件的全部各行，包括在 $\alpha_0=1$ 时被其他行支配的一行。

先设 $C=D/(Ax^{b_I})\ge1$。将选定位置放入第二组，再把其余位置以贪心方式放入容量 $A,C$。每次放入时，选定平方条件使新位置的平方不超过剩余容量的乘积，因此至少一个容量可以容纳它。所得分组端点满足
\[
 B_1\le A,\qquad B_2=x^{b_I},\qquad B_3\le C
\]
以及
\begin{align}
 A^2B_2B_3^2&\le x^{1-4\zeta},\notag\\
 N^\vartheta AB_1B_2^5B_3^2&\le x^{2-4\zeta},
 \label{coef:three-monomials}\\
 A^{1+2\vartheta}B_1B_2^{5-2\vartheta}B_3^{2+2\vartheta}
 &\le x^{2-(4+4\vartheta)\zeta}.\notag
\end{align}
若 $C<1$，则利用完整平方条件将所有位置放入 $A,D/A$。此时 $1<D/A<x^{b_I}$。令 $B_3=1$，并在最后两个界中将 $x^{b_I}$ 换为 $D/A$。第一个界来自 $A^2B_2\le AD\le x^{23/24-5\zeta/2}$。所以不会引入端点低于一的因子。

对三位置 $W$ 备选，用前两个位置初始化容量 $A,x^{v_0},D/(Ax^{v_0})$。其中一个可放入 $x^{v_0}$，另一个的平方不超过其余两个容量的乘积 $x^{1-\theta_{\rm cov}-2\zeta}$。继续逐个放入完整位置。若首次出现一个位置不能放入任何容量，则将其并入第二组。新乘积 $B$ 位于 $(x^{v_0},x^{u_0}]$，因为该位置不超过第三个位置，且 $p_3\le x^{u_0-v_0}$。第一容量无法容纳它，加上完整平方条件，给出 $ABc<D$，故其余两个剩余容量可用。继续用完整平方贪心论证完成分组。若从未失败，则容量乘积直接给出 \eqref{coef:three-monomials}。这些算法只使用整个原子的端点数据及中等长度，绝不依赖具体算术元组或等差数列偏差的符号。

\subsection{系数恒等式与无平方因子限制}

将区间标签细化为共同的不交半开标签集。若 $m_I$ 个位置具有标签 $I$，令 $M=\prod_I m_I!$。位置列表的未加限制素数卷积 $N_J$ 在其支撑上取值 $M/\prod_p\nu_p(n)!$，故相应有序无平方因子系数为 $\one_{\mathrm{sf}}N_J/M$。对一个固定三分组，令 $M_j=\prod_I m_{I,j}!$，并设 $\gamma_j$ 为组卷积除以 $M_j$。则在每个整数上 $0\le\gamma_j\le1$，且有精确恒等式
\begin{equation}
 \frac{N_J}{M}
 =\frac{M_1M_2M_3}{M}\,
       (\gamma_1*\gamma_2*\gamma_3).
 \label{coef:tensor-identity}
\end{equation}
空组是在一处的单位质量。带标记卷积已经包含其标记重数，只将有序剩余位置除以相应阶乘：对带符号原子
$\epsilon\one_{\mathrm{sf}}N_{\mathrm{all}}/M_{\mathrm{res}}$
重新分组后，外部标量为 $\epsilon M_1M_2M_3/M_{\mathrm{res}}$。这一标量固定，但不必小于或等于一。

选择低于每个带标记素数指数的 $w=N^{\eta_0}$。对 $p<w$，$p\mid BN$ 当且仅当 $p\mid N$。因此，只需一对支撑低于最小剩余预留量的小素数 beta 权，即可适用于有限带标记族的全部状态。原子的符号决定接入其中哪一个。小素数集合和粗糙素数集合互不相交，分解唯一。将小素数 beta 权接入 \eqref{coef:tensor-identity} 的一个组，仍保持绝对值不超过一。

全局无平方因子条件保留在实际系数及其密度中，仅在解析张量展开时去掉。若支撑为 $D_*\le N^{\theta_*}$ 且 $\theta_*<1$，偏差的变化至多为一个固定除数因子乘以
\begin{align}
 &\sum_{\substack{q\le D_*\\p^2\mid q\text{ 对某个 }p\ge N^{\eta_0}}}
 \left(\frac{N}{q}+1+\frac{N}{\varphi(q)}\right)
 \ll (N+D_*)N^{-\eta_0+o(1)}.
 \label{coef:collision}
\end{align}
其中 $1$ 包括完整乘积恰为 $N$ 的情形。没有任何除数函数在零处取值。

\subsection{参数次序与中等长度估计}

接入小素数、扩张端点及作局部重标度之前，上界族、普通下界族和带标记族的粗糙阶数至多为二的完整支撑指数分别为
\[
 \frac{19}{48},\qquad \frac{38}{125},\qquad
 \frac{1527}{3584}
\]
它们低于 $1/2$ 的余量分别是 $5/48$、$49/250$、$265/3584$。最后一项用到 $x<57/224$、$x+2y\le\Lambda_{\rm cap}$，从而有 $x+y\le(\Lambda_{\rm cap}+57/224)/2=1527/3584$。固定的系数范数允许对这些部分使用低于二分之一水平的 Bombieri--Vinogradov 估计。高阶系数则需要前述分解。

先固定有限组估计及其依赖关系、粗层规则、分解条件及正预留量。选取低于这些预留量和素数截点的 $\zeta$，并令 $\sigma=\zeta/2000$。再选共同的 beta 支撑，使其对密度和支撑容差都足够小。若粗糙阶数界为 $J$，则选择端点网格 $\tau$，使 $(J+2)\tau$ 小于每一剩余预留量。仅在固定的除数阶数与矩阶数确定后，才选择
\begin{equation}
 0<\kappa\le\min\left\{\kappa_*,\frac{\sigma}{200},
                         \frac1{2K_{\mathrm{mom}}}\right\},
 \qquad
 0<\kappa_0<\frac{\kappa}{4(1+\kappa)}.
 \label{coef:parameter-order}
\end{equation}
其中 $K_{\mathrm{mom}}$ 控制 Fourier 尾部和互补系数截断中的固定矩阶数。必要时再缩小 $\kappa_0$，使由 $\theta_*=599/1000$ 安全控制的完整支撑满足 $\theta_* /(1-\kappa_0)<3/5$，并使三个低阶支撑都保持低于二分之一。

在 \eqref{coef:three-monomials} 中将 $A$ 替换为 $U$ 后，仍有节省量 $\zeta,2\zeta,9\zeta/16$。扣除两项 $\zeta/16$ 的预留量后，最小节省量为 $7\zeta/16>15\sigma$。小素数只接入第三组，其代价 $2\ell_\beta,2\ell_\beta,
(2+2\vartheta)\ell_\beta$ 已纳入预留量。因而二进细分产生满足命题~\ref{prop:current-medium} 全部四个大小条件的 $Q_1,Q_2,Q_3$。严格余量吸收固定端点常数。在 $N$ 增大之前，原子只有有限个，绝对标量总范数固定。由此建立了实际有限系数族的一致三因子表示。

\hypertarget{coef:medium-checks}{}
\paragraph{中等长度分支的条件核验。}
下表将命题~\ref{prop:current-medium} 的条件与命题~\ref{coef:prime-product-insertion} 中的实际应用对应起来。核验在固定共同有限规则及其严格支撑预留量之后进行。
\begingroup
\renewcommand{\LTcaptype}{}
\stepcounter{LT@tables}
\small
\setlength{\tabcolsep}{4pt}
\begin{longtable}{@{}>{\raggedright\arraybackslash}p{.21\textwidth}>{\raggedright\arraybackslash}p{.45\textwidth}>{\raggedright\arraybackslash}p{.28\textwidth}@{}}
\toprule
 条件或输入 & 应用中的实现 & 核验位置 \\
\midrule
\endhead
\bottomrule
\endfoot
完整尺度与剩余类 & $x=UV$ 是完整乘积的尺度，$\mathfrak a=N$。截断及局部重标度给出 $c_0x\le N\le x^{1+\kappa}$。 & \eqref{coef:complete-scales}、\eqref{coef:parameter-order}\\[3pt]
粗糙中等长度序列 & 单独区间上的粗糙 $1$ 序列或 M\"obius 序列带有界权及允许的 Mellin 扭曲；保留辅助互素条件下的估计。 & \eqref{coef:rough-SW} 及其证明\\[3pt]
互补序列 & 唯一分解给出固定除数界。应用局部估计前，任何对数截断的代价均由混合矩估计控制。 & 引理~\ref{coef:mixed-moment}；命题~\ref{coef:prime-product-insertion} 的证明\\[3pt]
三个模数因子 & 归一化组卷积彼此分离、绝对值不超过一，并保留与剩余类互素的支撑。带标记重数留在外部标量中。 & \eqref{coef:tensor-identity}；完整与选定平方条件 \eqref{coef:full-squares}--\eqref{coef:selected-squares}\\[3pt]
四项严格大小检验 & 三个单项式界、中等长度容量、小素数接入及局部重标度共同给出 \eqref{eq:current-medium-ranges} 的全部检验。 & \eqref{coef:three-monomials}；\eqref{coef:parameter-order} 后的预留量\\[3pt]
变动剩余类相位与谱输入 & 保留精确互反相位。其变差以及 packet 与谱范数均使用指定的局部余量。 & \eqref{eq:current-retained-phase}、\eqref{ana:phase-variation}；附录~\ref{app:classical}\\[3pt]
一致求和误差 & 固定原子标量、Fourier 范数及对数个坐标片由更强的局部对数精度覆盖。素数幂、平方因子碰撞及完整乘积恰为 $N$ 的情形分别计入。 & \eqref{coef:collision}；命题~\ref{coef:prime-product-insertion} 的证明\\[3pt]
\end{longtable}
\endgroup

从节省量 $\zeta,2\zeta,9\zeta/16$ 中扣除已述支撑预留量后，四项大小检验仍留有严格余量，可吸收固定的二分端点倍数。此项核验针对中等长度分支；命题~\ref{coef:prime-product-insertion} 的其余分支由低于二分之一水平的估计及附录~\ref{app:brun} 覆盖。上界系数、下界系数和普通带标记后继均以其实际带符号系数 $\rho_N$ 通过该命题进入应用。

\section{光滑素数乘积的分布}
\label{app:prime-products}

本附录为附录~\ref{app:coefficients} 的有限系数族证明 \eqref{eq:current-PP}。中等长度分支使用命题~\ref{prop:current-medium}；其余分支使用低于二分之一水平的估计或附录~\ref{app:brun}。

设光滑权 $W_N$ 依赖至多 $k_*$ 个带标签素数位置，满足 $p_i\ge cN^{a_*}$，乘积为 $O(N)$。定义
\[
 E_q(N;W)=\sum_{p_1,\ldots,p_k}W_N(p_1,\ldots,p_k)
 \left(\one_{p_1\cdots p_k\equiv N\ (q)}
       -\frac{\one_{(p_1\cdots p_k,q)=1}}{\varphi(q)}\right).
\]
几何权的混合 Euler 导数满足固定多对数界，不带附加算术筛选条件。

\subsection{两个辅助均值估计}

用 $d_j$ 表示 $j$ 重除数函数。由于对数截断要求固定矩阶数，下面记录初等矩估计及其对变动平移量的依赖。

\begin{lemma}\label{coef:mixed-moment}
固定非负整数 $R,S$ 和常数 $C\ge1$。令 $K_{\rm mom}=2^{2\max(R,S,1)}$、$0\le\kappa\le1/(2K_{\rm mom})$。对 $|a|\le Cx^{1+\kappa}$ 及区间 $J\subseteq[1,Cx]$，一致有
\[
 \sum_{\substack{b\in J\\b\ne a}}
       \tau(b)^R\tau(|b-a|)^S
 \ll_{C,R,S} x(1+\log x)^{K_{\rm mom}-1}.
\]
\end{lemma}

\begin{proof}
设区间 $I\subseteq[1,Z]$ 的长度为 $Y$。对乘积落在 $I$ 中的有序 $j$ 元组，选取最大坐标，得到
\begin{align*}
 \sum_{n\in I}d_j(n)
 &\le j\sum_{d\le Z^{1-1/j}}d_{j-1}(d)(Y/d+1)\\
 &\le j\{Y(1+\log Z)^{j-1}
              +Z^{1-1/j}(1+\log Z)^{j-2}\}.
\end{align*}
事实上，其余 $j-1$ 个坐标的乘积至多为 $Z^{1-1/j}$；去掉最大坐标限制只会增大和。最后两个界由把除数函数展开为有序乘积、再对末因子求和得到。对每个整数 $e\ge0$，有 $\tau(n)^e\le d_{2^e}(n)$。在素数幂上，可将 $\{0,\ldots,v\}^e$ 中的 $e$ 元组映到达到各个位置 $1,\ldots,v$ 的坐标子集的频数表；该表是把 $v$ 分为 $2^e$ 份的弱组合，并唯一决定原元组。

非零值 $|b-a|$ 以至多二的重数落在至多两个区间内，每个区间长度为 $O_C(x)$，上端点为 $O_C(x^{1+\kappa})$。对 $d_{K_{\rm mom}}$ 应用前面的区间界，得
\[
 \sum_{\substack{b\in J\\b\ne a}}\tau(|b-a|)^{2S}
 \ll x(1+\log x)^{K_{\rm mom}-1}
       +x^{(1+\kappa)(1-1/K_{\rm mom})}(1+\log x)^{K_{\rm mom}-2}.
\]
第二个指数至多为 $1-1/(2K_{\rm mom})-1/(2K_{\rm mom}^2)<1$。未平移的 $2R$ 矩也满足同样的界。Cauchy--Schwarz 不等式即得结论。在应用除数函数之前，已去掉 $b=a$ 项。
\end{proof}

\begin{lemma}[低于二分之一水平的卷积估计]
\label{coef:low-rank-BV}
固定 $\eta>0$、$c\ge1$。设 $x=UV$、$x^\eta\le U,V\le x^{1-\eta}$，且 $\alpha,\beta$ 支撑在 $[U,cU]$、$[V,cV]$ 上，满足
\[
 \sum_m|\alpha(m)|^2\ll U(\log x)^C,\qquad
 \sum_n|\beta(n)|^2\ll V(\log x)^C.
\]
它们不依赖于求和模数。假定对每个固定 $A_1,B>0$、每个 $k\ge1$、$r\le(\log x)^B$、$(b,r)=1$，都有
\begin{equation}
 \left|\sum_{\substack{m\equiv b\ (r)\\(m,k)=1}}\alpha(m)
 -\frac1{\varphi(r)}\sum_{(m,rk)=1}\alpha(m)\right|
 \ll_{A_1,B}\tau(k)^{C_0}U(\log x)^{-A_1},
 \label{coef:small-aux-SW}
\end{equation}
其中指数 $C_0$ 固定。则对每个 $A>0$，存在 $D>0$，使 $Q\le\sqrt x(\log x)^{-D}$ 蕴含
\[
 \sum_{q\le Q}\max_{(a,q)=1}
 \left|\sum_{mn\equiv a\ (q)}\alpha(m)\beta(n)
 -\frac1{\varphi(q)}\sum_{(mn,q)=1}\alpha(m)\beta(n)\right|
 \ll_A x(\log x)^{-A}.
\]
\end{lemma}

\begin{proof}
令 $L=\log x$。本原乘法大筛 \cite[定理~18.2]{Kedlaya} 为
\begin{equation}
 \sum_{r\le R}\frac r{\varphi(r)}\sum_{\chi\ (r)}^{*}
 \left|\sum_nz(n)\chi(n)\right|^2
 \le(H-1+R^2)\sum_n|z(n)|^2,
 \label{coef:multiplicative-LS}
\end{equation}
这里 $z$ 支撑在至多含 $H$ 个连续整数的区间上。特征正交性精确消去主特征。由本原特征 $\chi^*$ 诱导的模 $q=rk$ 特征满足 $\chi(m)=\chi^*(m)\one_{(m,k)=1}$，包括 $(r,k)>1$ 的情形。令 $\alpha_k=\alpha\one_{(m,k)=1}$、$\beta_k=\beta\one_{(n,k)=1}$，用 $A_k(\chi),B_k(\chi)$ 表示其特征和。由于 $\varphi(rk)\ge\varphi(r)\varphi(k)$，所求和不超过
\begin{equation}
 \sum_{k\le Q}\frac1{\varphi(k)}
 \sum_{1<r\le Q/k}\frac1{\varphi(r)}
              \sum_{\chi\ (r)}^{*}|A_k(\chi)B_k(\chi)|.
 \label{coef:induced-characters}
\end{equation}
两个互素限制都不可省略。它们不会增加 $L^2$ 范数，且固定 $k$ 时不依赖于 $r$。

设 $R_0=L^{B_0}$。对 $R<r\le2R$、$R\ge R_0$，由 Cauchy--Schwarz 及 \eqref{coef:multiplicative-LS}，内和至多为
\[
 \ll L^{C'}\left\{\frac xR+
 \sqrt x(\sqrt U+\sqrt V)+R\sqrt x\right\}.
\]
先作二进求和，再在 \eqref{coef:induced-characters} 中对 $k$ 求和，得到
\begin{equation}
 \ll L^{C'}\left\{\frac{xL}{R_0}
       +L^2\sqrt x(\sqrt U+\sqrt V)+Q\sqrt x\right\}.
 \label{coef:large-conductors}
\end{equation}
这里 $\sum_{k\le Y}1/\varphi(k)\ll1+\log Y$、$\sum_{k\ge1}1/(k\varphi(k))<\infty$，后一求和出现在最后一项中。两者均由 $n/\varphi(n)=\sum_{d\mid n}\mu(d)^2/\varphi(d)$、$1/(d\varphi(d))\le\tau(d)/d^2$ 得到。因 $\sqrt x(\sqrt U+\sqrt V)\le2x^{1-\eta/2}$，\eqref{coef:large-conductors} 的中间项有固定幂次节省。取 $B_0,D$ 充分大即可控制其余两项。

对 $r\le R_0$，将 \eqref{coef:small-aux-SW} 与本原非主特征相乘求和，均值为零，得到 $|A_k(\chi)|\ll\varphi(r)\tau(k)^{C_0}UL^{-A_1}$。$\beta_k$ 的 $L^1$ 界为 $O(VL^{C/2})$，故 \eqref{coef:induced-characters} 的小导子部分至多为
\[
 \ll xL^{C/2-A_1}
       \left(\sum_{r\le R_0}\varphi(r)\right)
       \left(\sum_{k\le Q}\frac{\tau(k)^{C_0}}{\varphi(k)}\right).
\]
选择整数 $J\ge C_0+1$。由 $n/\varphi(n)\le\tau(n)$、$\tau(n)^J\le d_{2^J}(n)$，末和至多为 $(1+\log Q)^{2^J}$。固定 $B_0$ 后，取 $A_1$ 充分大即完成证明。所有除数阶数都在所要求的对数精度之前固定。
\end{proof}

对原始素数位置，经典素数 Siegel--Walfisz 定理 \cite[Exercise~64]{TaoSW} 配合分部求和，在 $k=1$ 时以任意指定节省给出 \eqref{coef:small-aux-SW}。对任意 $k$，改变互素条件仅删除整除 $k$ 的素数，代价至多为 $2\|\alpha\|_\infty\omega_{\mathrm{pf}}(k)$；其固定对数范数由正幂尺度 $U$ 吸收。因此 $C_0=1$ 适用于全部所需指数。固定光滑轮廓、区间端点及具有固定多对数变差的 Mellin 扭转都保留这一估计。对多个原始素数，将其中一个放入 $\alpha$，其余因子放入 $\beta$。唯一分解给出固定除数界，从而给出 $\beta$ 所需的 $L^2$ 范数。

\subsection{素数乘积估计的证明}

\begin{proposition}\label{coef:prime-product-insertion}
对附录~\ref{app:coefficients} 构造的有限上界、下界及带标记系数，命题~\ref{prop:current-medium} 的受限中等长度估计与低于二分之一水平的 Bombieri--Vinogradov 估计推出
\[
 \sum_q\rho_N(q)E_q(N;W)=O_{A,W}\bigl(N(\log N)^{-A}\bigr)
 \qquad(A>0).
\]
计数和精确互素均值使用同一系数。该估计对固定有限族及全部充分大的偶数 $N$ 一致。
\end{proposition}

\begin{proof}
选取光滑 $\eta$，使其在 $[0,1]$ 上为零、在 $[2,\infty)$ 上为一，并将 $W_N$ 乘以 $\eta(p_1\cdots p_k/N^{1-\kappa_0})$。被删乘积小于 $2N^{1-\kappa_0}$，对充分大的 $N$ 严格低于 $N$。其完整模数除数和及精确均值的误差为 $O(N^{1-\kappa_0+o(1)})$。这一截尾的全部 Euler 导数均有界。

对每个原始素数插入非负光滑二进单位分解。在一个原始坐标片上，令 $R_i=2^{j_i}$、$X=\prod_iR_i$、$X^+=2^kX$。取 $f(v)=v^{1/(\log\log v)^3}$，在应用卷积恒等式之前先固定共同粗糙截点
\[
 z=f(2^{6k_*+2}X^+)
\]
每个原始素数最终都大于 $z$。将其素数示性函数替换为 $\Lambda(n)/\log n$，并删除真素数幂，代价具有固定幂次节省：$cN^{a_*}$ 以上的相应倒数质量为 $N^{-a_*/2+o(1)}$。非零完整差值的除数代价为 $N^{o(1)}$；精确乘积 $N$ 单独贡献至多 $N^{\theta_*+o(1)}$。

对尺度 $R$ 的分组，令 $Z=(2R)^{1/3}$、$\mu_Z=\mu\one_{[1,Z]}$。若 $\delta$ 表示卷积恒等元，则 $\delta-1*\mu_Z$ 支撑在 $Z$ 以上，故 $2R$ 以下有 $\Lambda*(\delta-1*\mu_Z)^{*3}=0$。展开得到
\begin{equation}
 \Lambda=\sum_{j=1}^3(-1)^{j-1}\binom3j
             \mu_Z^{*j}*1^{*(j-1)}*\log.
 \label{coef:HB}
\end{equation}
保留原始光滑坐标片。粗糙性投影可分解到 \eqref{coef:HB} 中的所有变量上；使用 $P^-(1)=\infty$。不增加尖锐的耦合分组边界。

每个原始组至多有六个变量。与其光滑支撑相交的各个矩形，具有满足下式的完整局部尺度：
\begin{equation}
 2^{-7k-1}X\le x\le2^{6k-1}X^+,
 \qquad f(x)\le z\le x^{2/(\log\log x)^3}.
 \label{coef:complete-scales}
\end{equation}
所以，由 \eqref{coef:parameter-order} 有 $cN^{1-\kappa_0}\le x\le CN$、$c_0x\le N\le x^{1+\kappa}$。尺度 $x$ 始终对应完整乘积。

为分离几何轮廓，采用坐标 $u_j=\log(n_j/U_j)$，并取轮廓的光滑紧支撑延拓 $w$。写作 $D_w=\sum_{S\subseteq\{1,\ldots,m\}}\|\partial_{2S}w\|_1$，分部积分给出
\[
 |\widehat w(\xi)|\le D_w\prod_j(1+\xi_j^2)^{-1},
 \quad \frac1{(2\pi)^m}\int|\widehat w(\xi)|\,d\xi
 \le\frac{D_w}{2^m}.
\]
$|\xi|_\infty\le H$ 以外的尾部为 $O(mD_w/H)$，其中 $D_w$ 为固定对数幂。算术和可与绝对可积的 Fourier 表示交换。对完整模数求和后，非零差值 $b-N$ 处的尾部由固定混合除数矩控制为
\[
 \sum_{b\asymp x}\tau(b)^R\tau(|b-N|)^S
 \ll x(\log x)^C.
\]
应用引理~\ref{coef:mixed-moment}，其固定阶数纳入 $K_{\mathrm{mom}}$。精确差值为零的情形单独用 $N^{\theta_*+o(1)}$ 估计。选取充分大的对数参数 $H$。对保留的频率，支付的是 Fourier 积分范数，而非频率个数。

若存在中等长度位置，就在 $X^{3\zeta}\le U\le X^{1/3+\zeta/4}$ 中选取一个。由 \eqref{coef:complete-scales}，其局部长度处于 $x^{2\zeta}\le U\le x^{1/3+\zeta/2}$；其序列具有形式
\[
 \alpha_n=e_n\one_I(n)\one_{P^-(n)>z}h(n)n^{it},
 \qquad e_n=1\text{ 或 }\mu(n),\quad |h|\le1,
\]
且变差与频率具有固定多对数界。中等长度估计所需的辅助 Siegel--Walfisz 性质为
\begin{equation}
 \left|\sum_{\substack{n\equiv b\ (q)\\(n,k)=1}}\alpha_n
 -\frac1{\varphi(q)}\sum_{(n,kq)=1}\alpha_n\right|
 \ll_A U\tau(k)(\log x)^{-A}.
 \label{coef:rough-SW}
\end{equation}
为完整起见，可在偶数 Brun 阶数 $h=2\lfloor(\log\log x)^2\rfloor$ 插入粗糙性；下一层除数支撑为 $z^{h+1}<U^{1/4}$，而上下差的密度具有任意对数衰减。对无平方因子的 $\ell$，有 $\mu(\ell v)=\mu(\ell)\mu(v)\one_{(v,\ell)=1}$。商变量 $v$ 的尺度记为 $V_\ell=U/\ell$。恒等式 \[
 \mu(v)\one_{(v,\ell)=1}
 =\sum_{\substack{w\mid v\\p\mid w\Rightarrow p\mid\ell}}\mu(v/w)
\] 将小模数计算化为 $\mu$ 的未加粗糙限制的 Siegel--Walfisz 估计。在 $w\le V_\ell^{1/4}$ 处截断后留下 $O(V_\ell^{7/8}\tau(\ell)^2)$，最小内层尺度为 $U^{9/16}$。对 $\ell\le U^{1/4}$ 求和后为 $O(U^{29/32}\log^4U)$。大模数由绝对值不超过一的区间估计控制。插入 $(n,k)=1$ 仅删除高于 $z$ 的素数 $p\mid k$，代价至多为 $2(U/z+1)\omega_{\mathrm{pf}}(k)$。先取更强的初始对数精度，分部求和即可吸收允许的变差与扭转。这就对实际因子证明了 \eqref{coef:rough-SW}。

互补卷积由固定除数幂控制。必要时，仅将这一互补系数在固定对数阈值处截断：对非零差值使用同一混合除数矩界，在零处使用 $N^{\theta_*+o(1)}$。阈值在固定阶数确定后选择，再以更强局部对数估计恢复归一化。中等长度位置不作截断。现在对精确张量表示 \eqref{coef:tensor-identity} 应用命题~\ref{prop:current-medium}。

若没有中等长度变量，则每个原始分组均含有长变量：由于 $\zeta<a_*/100$，六个小于 $X^{3\zeta}$ 的变量之积不可能与 $R_i$ 同阶。长变量至多有两个，否则三个长变量将超过完整尺度。截断 M\"obius 变量为 $O(X^{1/3})$，不能是长变量。因此，剩下一或两个无权的长整数变量，同时保留其对数、光滑扭曲、独立区间及共同粗糙性投影。两种情形的互补短乘积分别不超过 $x^{20\zeta}$ 和 $x^{40\zeta}$。附录~\ref{app:brun} 在保留精确互素均值的条件下估计这两种情形。

最后，低阶系数部分的完整支撑与 $x^{1/2}$ 之间有固定余量。对一个原始素数，使用带最大值形式的素数 Bombieri--Vinogradov 估计；对两个或更多原始素数，按前述方法取一个原始素数位置，应用引理~\ref{coef:low-rank-BV}。有界的系数总范数允许使用其绝对值形式。全部固定原子标量、对数数量的坐标片及 Fourier 积分范数，均以更强的局部对数精度吸收。移去的素数幂、碰撞及截尾都具有固定幂次节省。因此，对这一有限族存在统一的最终阈值。
\end{proof}

\section{一个与两个长粗糙变量}
\label{app:brun}

本附录的估计处理命题~\ref{coef:prime-product-insertion} 证明中不含中等长度变量的分支。

对于一或两个长变量的情形，我们使用直接的等差数列估计及经典的完全 Kloosterman 和界。论证保留精确互素均值，并对本原剩余类的大小一致。

\subsection{具有一致密度差的 Brun 多项式}

设 $L=\log\log x$、$h=2\lfloor L^2\rfloor$、$r=h+1$、$\log z\le2\log x/L^3$。令 $P(z)=\prod_{p\le z}p$、$y=z^r=x^{o(1)}$。定义
\[
 \lambda^+(d)=\mu(d)\one_{d\mid P(z),\,\omega_{\mathrm{pf}}(d)\le h},
 \qquad
 \lambda^-(d)=\mu(d)\one_{d\mid P(z),\,\omega_{\mathrm{pf}}(d)\le h+1},
\]
及
\[
 U_z(n)=\sum_{d\mid n}\lambda^+(d),\quad
 L_z(n)=\sum_{d\mid n}\lambda^-(d),\quad
 G_z=U_z-L_z,\quad r_z(n)=\one_{(n,P(z))=1}.
\]
两个除数和的系数绝对值均不超过一，支撑低于 $y$。交错二项式恒等式给出
\begin{equation}
 L_z\le r_z\le U_z,\qquad U_z\ge0,\qquad
 G_z(n)=\binom{\omega_{\mathrm{pf}}((n,P(z)))}{r}\ge0.
 \label{coef:Brun-order}
\end{equation}
这也包括 $(n,P(z))=1$ 的情形，此时两个多项式均为一。下界多项式不必非负。

对每个 $q$，令
\[
 A_q=\sum_{(d,q)=1}\frac{\lambda^+(d)}d,\quad
 g_q=\sum_{(d,q)=1}\frac{\lambda^+(d)-\lambda^-(d)}d,\quad
 V_q=\prod_{\substack{p\le z\\p\nmid q}}\left(1-\frac1p\right).
\]
将有限容斥应用于均值为 $1/p$ 的独立 Bernoulli 变量，得到
\begin{equation}
 0\le A_q-V_q\le g_q
 \le\frac{(\sum_{p\le z}1/p)^r}{r!}=:\epsilon_x.
 \label{coef:Brun-gap}
\end{equation}
该界对整除 $q$ 的素数集合一致。初等素数调和界 $\sum_{p\le z}1/p\ll L$ 与 $r!\ge(r/e)^r$ 最终给出 $\epsilon_x\le\exp(-cL^2\log L)$。因此 $\epsilon_x$ 小于 $\log x$ 的每个固定负幂，而 $y=x^{o(1)}$ 仅与固定幂次节省相乘出现。

\subsection{光滑等差数列的核心估计}

对光滑紧支撑轮廓，定义 $N_1(f)=\|f\|_1+\|f'\|_1$、$N_2(f)=\|f\|_1+\|f''\|_1$。等差数列求和及 M\"obius 反演给出
\begin{align}
 \sum_{n\equiv b\ (q)}f(n/U)&=\frac Uq\int f+O(N_1(f)),\notag\\
 \sum_{(n,q)=1}f(n/U)&=U\frac{\varphi(q)}q\int f
                        +O(\tau(q)N_1(f)).
 \label{coef:one-core}
\end{align}
这些估计对 $U>0$ 及本原剩余类 $b$ 一致。因此，对固定且与 $q$ 互素的 $m$，$mn\equiv a\pmod q$ 相对于精确互素均值的偏差为 $O((1+\tau(q)/\varphi(q))N_1(f))$。若 $(m,q)>1$ 且 $(a,q)=1$，则两项均为零。

对两个变量，使用 $e(t)=e^{2\pi it}$、$\widehat f(t)=\int f(s)e(-ts)\,ds$。设 $c_q$ 为 Ramanujan 和，$S(j,k;q)=\sum_{u\bmod q}^{*}e((ju+k\bar u)/q)$。当 $b=a\bar m\pmod q$ 时，双重周期化给出精确偏差
\begin{equation}
 \frac{UV}{q^2}\sum_{j,k\in\mathbb Z}
 \widehat f(jU/q)\widehat g(kV/q)
 \left(S(bj,k;q)-\frac{c_q(j)c_q(k)}{\varphi(q)}\right).
 \label{coef:two-Poisson}
\end{equation}
两条坐标轴均恒等为零，因为 $S(0,k;q)=c_q(k)$、$S(bj,0;q)=c_q(j)$、$c_q(0)=\varphi(q)$。对 $jk\ne0$，使用
\[
 |S(bj,k;q)|\le\tau(q)\sqrt q\,(j,k,q)^{1/2},\qquad
 |\widehat f(t)|\ll N_2(f)(1+|t|)^{-2}.
\]
对每个 $H>0$ 和 $d\ge1$，有 $\sum_{j\ne0,\,d\mid j}(1+|j|/H)^{-2}\le2H/d$。展开公因子并使用 $\sum_{d\mid q}d^{-3/2}\le\zeta(3/2)$，因此将 \eqref{coef:two-Poisson} 控制为
\begin{equation}
 N_2(f)N_2(g)
 \left(\sqrt q\,\tau(q)+\frac{\tau(q)^2}{\varphi(q)}\right).
 \label{coef:two-core-bound}
\end{equation}
Ramanujan 项使用 $|c_q(j)|\le(j,q)$、$\sum_{d\mid q}\varphi(d)/d\le\tau(q)$。当 $q=1$ 时偏差为零。

设短系数 $\alpha_m$ 满足 $\sum|\alpha_m|\ll M\log^C x$，其中 $x\asymp MU$ 或 $MUV$。若 $q\le D$、$|\rho(q)|\le\tau(q)^J\log^B x$ 且 $J,B$ 固定，则单变量与双变量核心误差之和分别为
\begin{equation}
 O(MD\log^{C'}x),\qquad O(MD^{3/2}\log^{C'}x),
 \label{coef:core-summed}
\end{equation}
插入 Brun 除数 $d_1,d_2\le y$ 后，轮廓尺度为 $U/d_1,V/d_2$，而其在 \eqref{coef:two-core-bound} 中的范数不变。仅在本原剩余类中将 $m$ 替换为 $md_1d_2$。固定乘积的界对该剩余类一致，故对除数对求和的代价至多为 $y^2$。

\subsection{以精确均值恢复粗糙条件}

对非负光滑 $F$，展开 \eqref{coef:one-core} 中的除数系数，得到
\begin{align}
 \sum_{(n,q)=1}F(n/U)U_z(n)
 &=U\frac{\varphi(q)}q A_q\int F+O(y\tau(q)N_1(F)),\notag\\
 \sum_{(n,q)=1}F(n/U)G_z(n)
 &=U\frac{\varphi(q)}q g_q\int F+O(y\tau(q)N_1(F)).
 \label{coef:exact-Brun-means}
\end{align}
两个左端均非负，且为精确互素和；因子 $\varphi(q)/q$ 始终保留。

对两个变量，记 $U_i=U_z(n_i)$、$r_i=r_z(n_i)$、$G_i=G_z(n_i)$。有效的乘积比较为
\begin{equation}
 0\le U_1U_2-r_1r_2\le G_1U_2+U_1G_2.
 \label{coef:rough-product-gap}
\end{equation}
因为 $U_1U_2-r_1r_2=(U_1-r_1)U_2+r_1(U_2-r_2)$。这一论证不把两个可能为负的下界多项式相乘。对复轮廓，选择非负光滑控制函数 $F\ge|f|$、$G\ge|g|$，仅在估计差的绝对值时使用 $|\alpha|$。若 $C_q$ 表示同余计数，$M_q$ 表示其精确互素均值，且 $B=G_1U_2+U_1G_2$，则
\begin{align*}
 |(C_q-M_q)(\alpha fg(r_1r_2-U_1U_2))|
 &\le C_q(|\alpha|FGB)+M_q(|\alpha|FGB)\\
 &\le |(C_q-M_q)(|\alpha|FGB)|+2M_q(|\alpha|FGB).
\end{align*}
第一项通过展开 Brun 除数并使用 \eqref{coef:core-summed} 控制，第二项用 \eqref{coef:exact-Brun-means} 估计。固定的除数与 Euler 函数均值界给出双变量总误差
\begin{equation}
 \ll\log^{C'}x\left\{
 MUV\epsilon_x+M(U+V)y+My^2+My^2D^{3/2}\right\}.
 \label{coef:two-rough-error}
\end{equation}
单变量版本为
\begin{equation}
 \ll\log^{C'}x\{MU\epsilon_x+My+MyD\}.
 \label{coef:one-rough-error}
\end{equation}
$\epsilon_x$ 只乘以固定对数因子。对长变量情形，$M\le x^{20\zeta}$ 或 $x^{40\zeta}$，且 $D\le x^{\Theta_{\mathrm{loc}}}$，其中 $\Theta_{\mathrm{loc}}<3/5$；另有 $\min(U,V)\ge x^\nu$，其中 $\nu>0$ 固定。选择 $\zeta$，使
\[
 20\zeta+\Theta_{\mathrm{loc}}<1,\qquad
 40\zeta+\frac32\Theta_{\mathrm{loc}}<1.
\]
则 \eqref{coef:two-rough-error}、\eqref{coef:one-rough-error} 中每个其余项都有固定幂次节省。因此，对任意固定 $A$，误差均为 $O_A(x\log^{-A}x)$。

\subsection{各变量的尖锐区间}

设 $\delta=(\log x)^{-B}$。固定正尺度上的区间具有光滑逼近，其误差可用端点处归一化宽度为 $O(\delta)$ 的非负光滑条带控制。若区间短于 $\delta$，单个条带即可覆盖。条带积分为 $O(\delta)$，一阶导数范数为 $O(1)$，二阶导数范数由 $\delta^{-1}$ 的某个固定幂控制。再乘以其余轮廓的固定包络。

对这些非负条带应用已证明的粗糙光滑估计。在双变量情形，其精确均值至多为
\[
 \ll\log^{C'}x\{x\delta+M(U+V)y+My^2\}
\]
单变量情形则至多为 $\log^{C'}x\{x\delta+My\}$。第一项的指数取决于条带积分而非导数，且不依赖 $B$。依赖 $B$ 的导数损失只影响固定幂次误差。在确定所需对数精度之后，再选择 $B$。由此同时控制计数及精确均值，并对每个本原剩余类一致，包括剩余类等于某个中间完整乘积的情形。这就证明了所需的各变量尖锐区间推广；没有使用关于任意尖锐耦合区域的结论。

\section{依前缀确定的上界筛规则}\label{app:upper}

本附录给出第~\ref{sec:current-upper} 节所用上界规则的密度和系数实现；相应有限节省量在附录~\ref{app:numeric} 中计算。

我们通过对素数前缀作有限多个决定来构造除数系数。前缀按素数递减排列：若 $p_i=N^{t_i}$，则 $t_1>\cdots>t_r>0$。记 $S_j=t_1+\cdots+t_j$，$S_0=0$。删去前缀的最后一个素数即得其父前缀；添加更小素数时，保留该父前缀处已经作出的每个决定。

\subsection{可用上限与首次退出的奇偶性}

对非空子集 $I\subseteq\{1,\ldots,r\}$，记
\begin{equation}\label{sv:selected-gate}
 b_I=\sum_{i\in I}t_i,\qquad
 G_I(t)=\max_{\substack{j\le r\\j\notin I}}
 \left(b_I+\sum_{\substack{i<j\\i\notin I}}t_i+2t_j\right),
 \qquad \max\varnothing=0.
\end{equation}
完整平方条件为
\begin{equation}\label{sv:full-gate}
 G_{\mathrm{full}}(t)=\max_{1\le j\le r}(S_{j-1}+2t_j).
\end{equation}
候选选定乘积端点为
\begin{equation}\label{sv:tent}
 v(b)=\min\left\{\frac{1+b}{2},\frac{57}{64}-\frac{3b}{2}\right\}.
\end{equation}
紧前缀区域上的一个选定分支仅在其常数上限 $\lambda$ 在整个区域上都以固定正余量满足下式时才可采用：
\begin{equation}\label{sv:branch-licence}
 \max\{1/2,G_{\mathrm{full}}(t),G_I(t)\}
       <\lambda<v(b_I)
\end{equation}
端点为 $\Lambda_{\rm cap}-3t_3/8$ 的独立三位置分支，还须满足附录~\ref{app:coefficients} 的整盒分解条件。对候选端点取最大值时，始终只考虑通过相应条件的分支。

总指数为 $\lambda$ 的普通上界剩余树，在 $r$ 阶前缀后施加立方检验
\begin{equation}\label{sv:upper-cubes}
 S_{j-1}+3t_j<\lambda\qquad(j>r,\ j\text{ 奇数阶}).
\end{equation}
若 $r$ 为偶数，这些是普通上界剩余检验；若 $r$ 为奇数，则是普通下界剩余检验。此前的成员资格决策保持不变。在偶数前缀处，上限变更发生在接纳该前缀之后，因而不会在该阶引入排除。在奇数前缀处，可在自身奇数阶检验之前选择上限，但决策只能依赖于该前缀和先前固定的规则。两种操作都使首次排除保持在完整奇数阶。

\begin{lemma}[首次退出恒等式]\label{sv:first-exit}
设 $\mathcal T$ 是由有限素数集合 $\mathcal P$ 生成、包含空乘积且对前缀封闭的无平方因子乘积集合。假定所有首次被排除的前缀均为奇数阶。对每个满足 $0\le \xi(p)\le1$ 的乘法函数 $\xi$，有
\begin{equation}\label{sv:upper-exit}
 \sum_{d\in\mathcal T}\mu(d)\xi(d)
 =\prod_{p\in\mathcal P}(1-\xi(p))
 +\sum_{e\text{ 首次退出}}\xi(e)
           \prod_{\substack{p\in\mathcal P\\p<p_{\mathrm{last}}(e)}}
                            (1-\xi(p)).
\end{equation}
特别地，对每个整数 $n$，$\sum_{d\mid n,\ d\in\mathcal T}\mu(d)$ 都控制 $\one_{(n,\prod_{p\in\mathcal P}p)=1}$。
\end{lemma}

\begin{proof}
每个被排除的除数都有唯一的首次被排除前缀 $e$，其余素因子均小于 $p_{\mathrm{last}}(e)$。对完整的后续部分求和，得到 $\mu(e)\xi(e)\prod_{p<p_{\mathrm{last}}(e)}(1-\xi(p))$。从完整 Euler 乘积中减去这些互不相交的后续部分。由于 $\mu(e)=-1$，\eqref{sv:upper-exit} 中符号为正。最后一个结论取 $\xi(p)=\one_{p\mid n}$ 即得。同一有限恒等式也适用于 $n=0$ 和负整数。
\end{proof}

平方检验在后继中持续有效。保留的奇数阶立方条件蕴含其自身平方条件，随后的偶数阶平方条件由排序推出：
\begin{equation}\label{sv:descendant-square}
 S_{j-1}+2t_j\le S_{j-2}+3t_{j-1}<\lambda.
\end{equation}
全部选定位置出现后，对任意新接入位置，\eqref{sv:selected-gate} 中的表达式等于其完整平方。因此，在选定前缀处检验的 \eqref{sv:branch-licence} 会继续控制完整的后继向量。

\subsection{上限变更及其密度差}

令
\[
 \rho_0=\frac{19}{96},\qquad \mu_0=\frac{25}{128},\qquad
 \Theta=\frac{299}{500},\qquad \Lambda_{\rm cap}=\frac{153}{256}.
\]
第一个素数小于 $N^{\rho_0}$。$t_1=\mu_0$ 以下的三位置端点是下列候选在满足各自条件时的包络：
\begin{equation}\label{sv:third-envelope}
 v(t_1),\ v(t_1+t_2),\ v(t_1+t_3),\ v(t_2+t_3),\ v(S_3),\
                         \Lambda_{\rm cap}-\frac{3t_3}{8};
\end{equation}
$\mu_0$ 以上的初始三位置端点为 $v(t_1)$。\eqref{sv:third-envelope} 中每个分支保留独立的分解条件。特别地，所列候选并不意味着可以无条件选取任何非首部素数对。

在四阶处，选定候选为满足 $4\in I$ 的 $M=v(b_I)$，还需满足 \eqref{sv:branch-licence} 且超过继承的三阶上限。仅对已经通过第三立方检验的前缀施加这种变更。直到四阶的成员资格不变，后续奇数阶立方检验采用新上限。在区域 $\mathcal C$ 上，极限密度差为
\begin{equation}\label{sv:even-gain}
 4\int_{\mathcal C}
 \frac{\overline F((L-S_r)/t_r)
                   -\overline F((M-S_r)/t_r)}
      {t_1\cdots t_{r-1}t_r^2}\,dt_1\cdots dt_r,
 \qquad \overline F=\frac{F}{2e^\gamma},
\end{equation}
其中 $r=4$，$L$ 为继承上限，$M$ 为选定端点。事实上，偶数前缀 $P_r$ 对其完整子树的贡献为 $g_N(P_r)V_N^+(N^\lambda/P_r,p_r)$。将两棵子树相减即得 \eqref{sv:even-gain}；这是取渐近极限之前的原始密度恒等式。

在五阶处，将候选 $v(b_I)$（$5\in I$）与继承上限及第五立方 $S_4+3t_5$ 比较。其上限在该立方检验之前固定。因此，在较小上限处被拒绝的五阶前缀可在较大上限处保留，而所有更早决策不变。负的五阶子树具有普通下界剩余项，其节省量为
\begin{equation}\label{sv:odd-gain}
 4\int_{\mathcal C}
 \frac{\overline f((M-S_5)/t_5)
                   -\overline f((L-S_5)/t_5)}
      {t_1t_2t_3t_4t_5^2}\,dt_1\cdots dt_5,
 \qquad \overline f=\frac{f}{2e^\gamma},
\end{equation}
其中 $\overline f$ 在 $2$ 以下延拓为零。继承的前五阶上限不超过
\begin{equation}\label{sv:fifth-envelope}
 R_5(t)=\max\left\{\Lambda_{\rm cap}-\frac{3t_3}{8},\
         v(b_I):\varnothing\ne I\subseteq\{1,\ldots,5\}\right\}.
\end{equation}
在非负节省量中用继承上限的上界，只会得到保守估计，不会改变实际上限及其密度。

六阶变更是在偶数前缀处采用 \eqref{sv:even-gain} 的形式，并在保留第三及第五接纳检验后施加。特别地，在第五立方由 $L_3$ 控制的区域上，有 $L_3-S_6>2t_5-t_6>t_6$。在 \eqref{sv:fifth-envelope} 的每个帐篷函数中选择一张定义它的仿射上平面，即得凸的分段仿射控制函数 $U_5\ge R_5$。对新的凹端点 $M$，令 $d=M-U_5$，则在 $d>0$ 上，节省量至少为
\begin{equation}\label{sv:sixth-lower}
 \frac{4d_+}{R_oR_n\prod_{i=1}^6t_i},
\end{equation}
其中 $R_n\ge M-S_6$、$R_o\ge\min\{U_5-S_6,M-S_6\}>0$ 是整个区域上的上界。先取正部，再提取非负素数密度测度的下界。对矩形区域，由凹性有 $\int d_+/\operatorname{Vol}(\mathcal C)
\ge 2^{-6}\sum_{v\text{ 顶点}}d(v)$。

两个附加变更位于 $\mu_0<x<\rho_0$ 的不交部分。对二阶变更，令 $a(x)=2v(x)-1=25/32-3x$ 并使用
\[
 a(x)<y<x,\qquad L=v(x),\quad M=v(y).
\]
所需条件由 $2x,x+2y,2x+y< v(x)<v(y)$ 推出，选定位置为 $2$。节省量是 \eqref{sv:even-gain} 在 $r=2$ 时的情形。对三阶变更，令 $t=y+z$，并施加
\[
 a(x)<t<x,\quad z=tw/2,\quad 5/9\le w<1,
 \qquad L=v(x),\quad M=v(t).
\]
这里 $y<3/16$，而二阶变更满足 $y\ge3/16$。选定位置 $\{2,3\}$ 给出附加条件 $2x+t<M$。保留第三节点并改变其下界剩余项，得到 \eqref{sv:odd-gain} 的 $r=3$ 类似式。其自变量位于 $[2,4]$，此处有 $\overline f(s)=\log(s-1)/s$。

在八阶处，将选定端点 $M$ 与继承的六阶上限的有效上界 $R_6$ 比较。保留的第七立方条件以及附加条件 $M-S_8<3t_8$ 给出
\begin{equation}\label{sv:eighth-range}
 1<\frac{\lambda_6-S_8}{t_8}
 \le\frac{R_6-S_8}{t_8}<\frac{M-S_8}{t_8}<3.
\end{equation}
因此，\eqref{sv:even-gain} 至少为下列函数的积分：
\begin{equation}\label{sv:eighth-density}
 \frac{4(M-R_6)}{(R_6-S_8)(M-S_8)\prod_{i=1}^8t_i}.
\end{equation}
每个选定区域必须保留全部完整及选定补集平方条件，以及所有祖先检验。同一阶段作用的区域内部互不相交。不同阶段通过精确的望远镜恒等式 $V_3-V_6=(V_3-V_4)+(V_4-V_5)+(V_5-V_6)$ 合并。每个密度差采用一个下包围。

\subsection{实际普通密度与有限实现}

设 $g_N(p)=\one_{p>2,\ p\nmid N}/(p-1)$、$V_N(y)=\prod_{p<y}(1-g_N(p))$。普通立方 Rosser 系数在所述范围的紧子集上一致满足
\begin{equation}\label{sv:literal-density}
 \frac{V_N^+(N^\beta,N^t)}{V_N(N^t)}=F(\beta/t)+o(1)
       \quad(\beta/t>1),\qquad
 \frac{V_N^-(N^\beta,N^t)}{V_N(N^t)}=f(\beta/t)+o(1)
       \quad(\beta/t\ge2).
\end{equation}
这里的密度是实际系数的密度，而不只是任意筛序列的上、下界。

为说明对 $N$ 的一致性，将归一化的固定阶退出与 $g_0(p)=\one_{p>2}/(p-1)$ 的相应退出比较。保留的前缀获得乘子
\[
 \prod_{\substack{p_{\rm last}\le p<y\\p\mid N}}(1-g_0(p))\le1;
\]
其余前缀消失。在正参数紧范围上，$j$ 阶退出的末素数至少为 $N^{\beta_{\min}/(3\cdot2^{j-1})}$。因此，删除 $N$ 的素因子使每个固定阶积分一致改变 $o(1)$。非负 Rosser 退出级数在参数紧范围上一致收敛，而原始单侧界提供有限阶夹逼的另一侧，见 \cite[命题~11.8 及 引理~14.7]{Suzuki}。先选阶数，再令 $N$ 增大，即证明 \eqref{sv:literal-density}。当 $1<s\le3$ 时，首个立方处的恒等式将上界计算归约至截点 $N^{\beta/3}$，给出 $F(s)=2e^\gamma/s$。

归一化来自 Mertens 公式。特别地，与 $2C_N\prod_{p<z}(1-1/p)$ 的精确比值是
\begin{equation}\label{sv:euler-ratio}
 \frac{V_N(z)}{2C_N\prod_{p<z}(1-1/p)}
 =\prod_{p\ge z}\left(1-\frac1{(p-1)^2}\right)^{-1}
  \prod_{\substack{p\mid N\\p\ge z}}\left(1-\frac1{p-1}\right).
\end{equation}
其对数为 $O((1+\log N/\log z)/(z-2))$，故 $V_N(N^t)\sim2e^{-\gamma}C_N/(t\log N)$ 对 $N$ 一致。\eqref{sv:even-gain} 和 \eqref{sv:odd-gain} 中的因子 $4$ 恰为 $2e^{-\gamma}\cdot2e^\gamma$。

有限单元数据由有理区域、相应阶数、选定子集、端点及父标签组成。每个有效区域的准入与支撑不等式均具有正余量；继承区域保留父前缀的决定。密度是精确子树差之和。对于正的节省量，任取一个子族都给出下界；对于负的代价，上界必须覆盖完整的保留区域。

\begin{proposition}[有限规则的实现]\label{sv:finite-realization}
固定这样的一个有限可用规则及正误差 $\epsilon$，并假设其选定分支满足附录~\ref{app:coefficients} 的系数分解条件。则存在固定的有限细分及有界无平方因子系数 $\rho_N^{\rm up}$，支撑于 $(q,N)=1$，具有下列性质：其除数和为非负的粗糙性示性函数上界；带符号分布估计 \eqref{eq:current-PP} 适用于同一系数；其密度相对于理想规则密度，在所需上界方向上的差至多为 $\epsilon C_N/\log N+o(C_N/\log N)$。
\end{proposition}

\begin{proof}
先截去所有末素数很小的尾部。在第三次激活处，可取 $u=\rho_0/2$、$0<h<u$ 并使用
\begin{equation}\label{sv:activation-tail}
 T_3(h)\le\frac{18h}{(u-h)^2}
   \left(\log^2\frac{\rho_0}{h}
                   +2\log\frac{\rho_0}{h}+2\right)=o(1).
\end{equation}
这由 $s\ge3$ 时的 $\overline F(s)-\overline f(s)\le9/s^2$、剩余量下界 $\lambda-t_1-t_2\ge u$，以及质量为 $\tfrac12\log^2(\rho_0/t_3)$ 的有序素数倒数三角形得到。更高阶选定变更有正的末素数截点。真正的仿射边界具有任意小的条带体积；恒等条件和开放继承区域则保留。在剩余紧区域上选择具有严格余量的粗层上限，再按前缀阶数递增地细化。任何祖先决策都不能依赖于后来接入的素数。

第三端点是分段仿射 Lipschitz 函数。对宽度 $\chi<e/10$ 的盒，形如 $\lambda_{\mathscr B}=\Gamma_{\mathscr B}-2e$ 的上限以至多 $4e+6\chi$ 的误差逼近它。在端点为 $W=\Lambda_{\rm cap}-3t_3/8$ 的三位置分支附近，辅助分解覆盖可以超过实际上限。若 $\inf_{\mathscr B} A\ge\lambda_{\mathscr B}+e/2$，其中 $A=(1+t_1)/2$，取 $\ell_{\rm aux}=\lambda_{\mathscr B}+e/4$；否则 $\sup_{\mathscr B} A<\lambda_{\mathscr B}+11e/20$，而 $\ell_{\rm aux}=\lambda_{\mathscr B}+e$ 比必要边界至少高 $9e/20$，比 $W$ 至少低 $e$。另一个初始化平方条件具有余量 $3/256$。辅助覆盖用于证明系数可用，绝不在密度或剩余筛级中替代实际上限。

选择预筛 $z_0=N^{\eta_0}$。偶数粗糙阶使用端点外检验，奇数粗糙阶使用内检验，保留所有祖先标签，得到原始上界系数的 $a_N(d)\ge\lambda_N(d)$。若 $B^-\le I_0\le B^+$ 为一维小素数 beta 多项式，则对正的 $a_N(d)$ 接入 $B^+$，对负的 $a_N(d)$ 接入 $B^-$，逐项有 $a_N(d)B_{\rm sign}\ge a_N(d)I_0$。唯一的小素数与粗糙部分分解定义所得 $\rho_N^{\rm up}$，并保持无平方因子性和有界性。

密度方向还需一个恒等式。将 \eqref{sv:upper-exit} 投影到至少为 $z_0$ 的素数，得
\begin{equation}\label{sv:upper-projection}
 V_N(z_0)\sum_{d\text{ 粗糙}}\lambda_N(d)g_N(d)
       \le\sum_{d\text{ 完整}}\lambda_N(d)g_N(d).
\end{equation}
左端只保留末素数至少为 $z_0$ 的非负退出。对固定粗层规则，beta 与精细盒的绝对误差相对于根部 Euler 乘积为
$O(e^{-c/\eta}\eta_0^{-2}+C_{\rm policy}\tau\eta_0^{-5})$
先选 $\eta,\eta_0$，再选 $\tau$，使这些误差及截去的部分落在 $\epsilon$ 内。

以上选择均先于 $N$。精细阶数由 $J_0>\theta_*/\eta_0+2$ 控制；支撑膨胀可保持低于 $1/1000$，从而得到 $\theta_*=599/1000$。最后选择固定矩阶数及完整尺度截尾，使 $\theta_*/(1-\kappa_0)<3/5$ 和全部局部不等式具有严格余量。附录~\ref{app:coefficients} 随即对这一系数给出 \eqref{eq:current-PP}。逐点界来自引理~\ref{sv:first-exit}；密度界来自 \eqref{sv:upper-projection} 与绝对误差控制。
\end{proof}

\section{下界规则与有限反馈}\label{app:lower}

本附录给出第~\ref{sec:current-lower} 节所用下界规则的实现及非负带标记转移不等式。对应有限映射和积分由附录~\ref{app:numeric} 计算。

\subsection{普通下界树与四素数前缀恢复}

令 ${a_{\rm s}}=417/5000$、${h_{\rm s}}=19/125$、$\nu=249/448$。对 $w={a_{\rm s}}$ 或 $w={h_{\rm s}}$，从总指数 $\nu_\delta=\nu-\delta$ 的普通下界支撑开始，其中 $0<\delta$ 固定且 $2{h_{\rm s}}<\nu_\delta$。在 $w$ 以下，首个平方条件 $2t_1<\nu_\delta$ 自动成立；后续检验为
\begin{equation}\label{sv:lower-cubes}
 S_{j-1}+3t_j<\nu_\delta\qquad(j\text{ 偶数阶}).
\end{equation}
在可接纳的三阶前缀处，固定 $\lambda_3\ge\nu_\delta$，供后续偶数阶立方检验使用。普通四阶变更只能在 $S_3+3t_4<\lambda_3$ 之后施加，并固定 $\lambda_4\ge\lambda_3$。随后接纳第五前缀，且可为后续偶数阶检验固定 $\lambda_5\ge\lambda_4$。每个活动选定上限均低于其可用端点，高于父上限，并满足全部条件 \eqref{sv:selected-gate}--\eqref{sv:full-gate}。对下界起始族，选定端点为
\begin{equation}\label{sv:lower-endpoint}
 v_-(b)=\min\left\{\frac{1+b}{2},
 \frac{32-34b}{39},\frac{57}{64}-\frac{3b}{2}\right\}.
\end{equation}
后继保留选定分支。下文引入的带标记状态剩余大小比，与分布估计中的局部剩余类指数是不同参数。

还可以恢复第四前缀。在原本会首次被排除的第四前缀上，在其自身立方检验之前选择 $\lambda_4$，满足
\begin{equation}\label{sv:restore-gates}
 \lambda_3\le S_3+3t_4<\lambda_4,
 \qquad \max\{G_{\rm full},G_I\}<\lambda_4<v_-(b_I).
\end{equation}
在 $p_4$ 以下接入筛级为 $N^{\lambda_4}/P_4$ 的普通下界剩余树。由于 $\lambda_4-S_4>2t_4$，其首个平方检验不会失败。其后首次退出的剩余阶数为偶数，因而完整阶数也为偶数；恢复后的树仍为下界树。密度增量恰为
\begin{equation}\label{sv:restoration-density}
 4\int_{\mathcal B}
 \frac{\overline f((\lambda_4-S_4)/t_4)}
      {t_1t_2t_3t_4^2}\,dt_1dt_2dt_3dt_4.
\end{equation}
该首次退出原先没有保留剩余树，因此 \eqref{sv:restoration-density} 不需要减去父剩余项。

下面定义恢复区域的一个有限有理族。设 $r_0=(\nu-{h_{\rm s}})/9$、$n\in\{400,800\}$，以及满足 $r_0\le\ell_4<\ell_3<\ell_2<\ell_1<{h_{\rm s}}$、$u_1+3u_2<\nu$ 的 $\mathscr B=\prod_{i=1}^4[\ell_i,\ell_i+1/n]$，其中 $u_i=\ell_i+1/n$。设 $U_3$ 为 $\nu$、全部三位置选定端点上界以及 $\Lambda_{\rm cap}-3\ell_3/8$ 的最大值，它从上方控制继承的三阶上限。取 $D=\operatorname{lcm}(448,256,78n,64n,8n)$，仅当
\[
 c_I=\inf_{\mathscr B} v_-(b_I)-D^{-1}
       >\max\{U_3,\sup_{\mathscr B} G_{\rm full},\sup_{\mathscr B} G_I\}.
\]
成立时，才保留含有 $4$ 的掩膜 $I$。选择最大的保留 $c_I$，记为 $c$。对 $A=t_1+t_2+t_3+3t_4$，将盒裁至
\[
 \max(A_-,U_3+D^{-1})<A<\min(A_+,c-D^{-1}),
\]
并将该区间等分为有限个有理层。每层都以正余量满足 \eqref{sv:restore-gates}，且位于先前保留的四阶前缀之外。在裁剪区域上，
\[
 2<\frac{c-S_4}{t_4}
       <2+\frac{9(\Lambda_{\rm cap}-\nu)}{\nu-{h_{\rm s}}}
       =2+\frac{84375}{90452}<3.
\]
因此 $\overline f(s)=\log(s-1)/s$ 全程递增。若一层上有 $A<a_j$，且 $u_4$ 是第四个上端点，则在下界
$4\overline f(s_j)\operatorname{Vol}(E_j)/(u_1u_2u_3u_4^2)$
中使用 $s_j=2+(c-a_j)/u_4$。$\mathscr B\cap\{A<t\}$ 的精确体积为
\begin{equation}\label{sv:slab-volume}
 \frac1{72}\sum_{e\in\{0,1\}^4}(-1)^{|e|}
 \left(t-A_--\frac{e_1+e_2+e_3+3e_4}{n}\right)_+^4.
\end{equation}
这就是加权盒宽 $(1,1,1,3)/n$ 的容斥公式。\eqref{sv:slab-volume} 之差精确给出各层体积。

\subsection{下界系数及其原始素数序列}

对首次退出为偶数阶的前缀封闭树，与 \eqref{sv:upper-exit} 对应的公式是
\begin{equation}\label{sv:lower-exit}
 \sum_{d\in\mathcal T}\mu(d)\xi(d)
 =\prod_p(1-\xi(p))
 -\sum_{e\text{ 首次退出}}\xi(e)
               \prod_{p<p_{\rm last}(e)}(1-\xi(p)).
\end{equation}
证明相同，只需使用 $\mu(e)=1$。取 $\xi(p)=\one_{p\mid n}$ 即得粗糙性示性函数的下界。取 $\xi=g_N$ 及预筛 $z_0$，得
\begin{equation}\label{sv:lower-projection}
 V_N(z_0)R_{\rm rough}-V_{\rm full}
 =\sum_{\substack{e\text{ 完整偶数阶退出}\\p_{\rm last}(e)<z_0}}
                      g_N(e)V_N(p_{\rm last}(e))\ge0.
\end{equation}
这正是下界系数实现所需的密度方向。

将正偶数阶系数向内取整，负奇数阶系数向外取整，保持前缀决策历史。对正的取整系数接入 beta 下界，对负的接入 beta 上界。逐项有 $a_{\rm rnd}(d)B_{\rm sign}\le a_{\rm rnd}(d)I_0$，因而实际系数满足
\begin{equation}\label{sv:physical-lower}
 \sum_{q\mid n}\rho_N^{\rm low}(q)\le
       \one_{(n,P_N(N^w))=1},\qquad
 \sum_q\frac{\rho_N^{\rm low}(q)}{\varphi(q)}
       \ge V_{\rm full}-\text{密度绝对误差}.
\end{equation}
第二个不等式使用 \eqref{sv:lower-projection} 及 beta 与盒的绝对误差估计，不仅由系数逐项大小关系推出；beta 下界本身也不必非负。附录~\ref{app:cross-review} 中的尖锐素数结论适用于同一 $\rho_N^{\rm low}$，将 \eqref{sv:physical-lower} 化为 $\{N-p:p<N,\ p\text{ 为素数}\}$ 的普通下界估计。

对于可用的下界前缀，第三、第四和第五阶段的密度变化依次是上界父项减上界子项、下界子项减下界父项、上界父项减上界子项。各项均由 \eqref{sv:literal-density} 计算。在 $r=3,4,5$ 阶，普通上、下界密度之差抵消最后坐标中表面的逆平方奇性，留下尾界 $O(\epsilon(1+\log^{r-1}({h_{\rm s}}/\epsilon)))$。余下紧区域可用一个共同的有限细分实现，并保留每个继承决定。

\subsection{非负带标记族上的有限反馈}\label{sv:finite-feedback}

归一化状态是对具有共同总筛级 $Q=N^\lambda$ 的非负有限带标记族 $\sigma(B)$ 作平均的估计。令
\[
 L=\lambda-\log_NB>0,\qquad
 \Theta_Q=4\li(N)\sum_B
       \frac{\sigma(B)C_{BN}}{\varphi(B)\log(Q/B)},\qquad
 \alpha=\frac{1-\log_NB}{L}\ge2.
\]
指数为 $Lu$ 的单素数子节点，或指数为 $Lu,Lv$ 的素数对子节点，其剩余筛级和参数分别为
\begin{equation}\label{sv:child-state}
 L'=L(1-u),\quad
 \alpha'=1+\frac{\alpha-1}{1-u};\qquad
 L'=L(1-u-v),\quad
 \alpha'=1+\frac{\alpha-1}{1-u-v}.
\end{equation}
总筛级 $Q$、选定位置及素数上限均保留。特别地，$\alpha$ 增大，子节点不会重新设为 $\alpha=2$。

这里，可用带标记族指带标签素数盒示性函数的有限非负组合，全部带标记素数指数及 $L$ 都有固定正下界。在经典情形，完整模数筛级比 $N^{1/2}$ 低一个固定幂次；在增强情形，$1/2<\lambda<153/256$ 带固定余量，并保留全部完整及选定补集平方条件，以及附录~\ref{app:coefficients} 的一个分解分支。还保留一个 $\chi>0$，使接入指数低于 $\chi$ 的素数后，该分支及其条件不变。$\chi$ 不大于已指定的最小指数。若某分支使用尚未指定的第三位置界 $c$，还要求 $\chi\le c$。自变量为 $s$ 的每个增强普通节点，其实际截点 $N^{L/s}$ 都低于 $N^\chi$。这些限制都在取端点极限之前，以每个盒上的固定余量施加。估计针对带标记平均，而不是某个单独固定的素数前缀。

子节点相对父节点的测度分别为 $du/[u(1-u)]$、$du\,dv/[uv(1-u-v)]$，均为非负。为指定单步换位代价，定义
\begin{equation}\label{sv:one-step-mass}
 I(s,s_0,\alpha)=
 \int_{1/s_0<u_1<u_2<u_3<1/s}
 \omega\!\left(\frac{\alpha-u_1-u_2-u_3}{u_2}\right)
                    \frac{du_1\,du_2\,du_3}{u_1u_2^2u_3}.
\end{equation}
单步不等式给出候选
\begin{align}\label{sv:one-step-operator}
 \mathcal W_1[H,h](s;s_0,\alpha)
 ={}&A_+(s)-A_+(s_0)+H(s_0)-\frac K8 I(s,s_0,\alpha)\notag\\
 &+\frac12\int_{1/s_0}^{1/s}
 \frac{A_-(s_0(1-u))+h(s_0(1-u))}{u(1-u)}\,du,
 \qquad 2\le s<s_0\le5.
\end{align}
其中 $A_+(s)=sF(s)/(2e^\gamma)$、$A_-(s)=sf(s)/(2e^\gamma)$，并在 $2$ 以下延拓为零。实际单素数子节点的截点为 $N^{L/s_0}$，自变量为 $s_0(1-u)$，而不是另一个总筛级下的自变量。非负换位族为 $Bq_1q_2q_3n\le N$，其中 $N^{L/s_0}\le q_1<q_2<q_3<N^{L/s}$、$P^-(n)\ge q_2$。在父节点归一化下，其未筛质量为 $I/4$；逐点加权不等式中的因子 $1/2$ 给出 \eqref{sv:one-step-operator} 中的代价 $KI/8$。所有素因子都具有固定正指数。由于 $\alpha-3/s\ge1/2$，其单位商情形具有固定幂次节省。

令 $c_\omega=70893/125000$，并取
\begin{equation}\label{sv:omega-envelope}
 \omega_{\rm env}(u)=
 \begin{cases}
 1,&u<1,\\
 \max\{1/u,c_\omega\},&u\ge1.
 \end{cases}
\end{equation}
可得 \eqref{sv:one-step-mass} 中适用的一致界。对 $v\ge u\ge1$ 有 $\omega(v)\le\omega_{\rm env}(u)$。在 $[1,2]$ 上由 $\omega(v)=1/v$ 推出；在 $[2,3]$ 上对 $1+\log(v-1)-c_\omega v$ 求最大值。最大值为 $-\log c_\omega-c_\omega<0$，因为精确有理不等式 $c_\omega\sum_{j=0}^8c_\omega^j/j!>1$ 蕴含 $c_\omega e^{c_\omega}>1$。延迟方程将该界传播到 $v\ge3$。因此，把被积函数替换为 $\omega_{\rm env}((2-u_1-u_2-u_3)/u_2)$，即可对每个 $\alpha\ge2$ 控制代价。$1$ 以下的分支对使用它的换位估计给出保守界；若单位商所在区域质量不可忽略，则单独计入单位商项。

对双步作用，记 $a_0=1/s_0$、$a_i=1/\phi_i$、$b=1/s$，参数条件为
\begin{gather}\label{sv:two-step-gates}
 2<s<\phi_3<\phi_2<\phi_1<s_0\le5,\qquad
 a_3+2a_1<1,\quad3a_1<1,\quad a_2+2a_1<1,\notag\\
 a_0<(1-a_2)/2<1/3<a_2<1/2,\qquad
 s_0(1-b)>2,\quad s_0-1\le4,\quad s_0(2-2a_2)>2.
\end{gather}
三个单素数区域为
\begin{equation}\label{sv:singleton-domains}
 (a_0,b),\qquad(a_0,a_2),\qquad(a_0,a_3),
 \quad \rho(u)=s_0(1-u),
 \quad\text{实际截点 }N^{La_0}.
\end{equation}
普通素数对估计及其反馈适用的区域为
\[
\begin{array}{c|c|c}
 &\text{区域}&\text{子节点自变量}\\ \hline
 D_5^+&a_0<u<v<a_2,\ v+2u<1&s_0(1-u-v)\\
 D_6&(a_0,a_1)\times(a_2,a_3)&s_0(1-u-v)\\
 D_7&a_0<u<v<a_1&(1-u-v)/u\\
 D_8&(a_0,a_1)\times(a_1,a_2)&(1-u-v)/u,
\end{array}
\]
在逐点加权不等式中，第五项具有完整母区域
\begin{equation}\label{sv:complete-fifth-region}
 D_5=\{a_0<u<v<a_2\}
     =D_5^+\sqcup D_5^-,\qquad
 D_5^-=D_5\cap\{v+2u\ge1\}.
\end{equation}
因此 $\Gamma_5=\Gamma_5^++\Gamma_5^-$；组合不等式中不能将其限制为 $D_5^+$。普通估计及子节点改进量仅在 $D_5^+$ 使用；在 $D_5^-$ 保留完整计数，并用非负乘积估计控制。$D_5$ 上处处满足检验 $2v<1$；由 \eqref{sv:two-step-gates}，其余三个素数对区域同时满足 $2v<1$、$v+2u<1$。若 $J_2^0$ 表示所选可用参数组的普通及非负换位基础项，则加入的反馈为
\begin{equation}\label{sv:two-step-feedback}
 \frac15\left\{4H(s_0)+H(\phi_1)
 +\sum_{c\in\{b,a_2,a_3\}}\int_{a_0}^{c}
       \frac{h(s_0(1-u))}{u(1-u)}\,du
 +\sum_{i=5^+,6,7,8}\iint_{D_i}
       \frac{H(\rho_i(u,v))}{uv(1-u-v)}\,du\,dv\right\}.
\end{equation}
下面定义这一基础项，并证明各贡献的符号。

\paragraph{二十一项不等式。}
对带标记状态，令 $R=Q/B$，本段另记 $\mathcal A_{N,B}=\{(N-p)/B:p<N,\ p\text{ 为素数},\ B\mid N-p\}$。用 $\Phi_s$ 表示在 $R^{1/s}$ 处作筛后的非负带标记平均。用 $[l,r]^k_\uparrow$ 表示 $[l,r]$ 中 $k$ 个递增排列的指数。对正整数 $n$，令
\[
 Q_d(y;n)=\one_{d\mid n}\,
                  \one_{P^-(n/d)\ge y}.
\]
下式全部新增素数互不相同；与带标记前缀的碰撞及平方因子均以固定幂次节省误差去掉。随后对每个表达式都采用相同的非负带标记权作平均。前四项为
\[
 \Gamma_1=4\Phi_{s_0}+\Phi_{\phi_1},\qquad
 \Gamma_2,\Gamma_3,\Gamma_4
   =\sum_{R^{a_0}\le q<R^c}Q_q(R^{a_0};n),
 \quad c=b,a_2,a_3,
\]
其中最后各式包含对 $n\in\mathcal A_{N,B}$ 的平均。在 \eqref{sv:complete-fifth-region} 中，对 $\Gamma_5$ 使用完整区域 $D_5$；对 $\Gamma_6,\Gamma_7,\Gamma_8$ 使用上述其余三个素数对区域。被加项在 $i=5,6$ 时为 $Q_{q_1q_2}(R^{a_0};n)$，在 $i=7,8$ 时为 $Q_{q_1q_2}(q_1;n)$，其中 $q_1=R^u<q_2=R^v$。对 $9\le i\le21$，有序指数区域为
\begin{equation}\label{sv:switch-domains}
\begin{array}{c|l@{\qquad}c|l}
9 &[a_1,a_3]^3_\uparrow
 &10 &[a_1,a_2]^2_\uparrow[a_2,b]\\
11 &[a_1,a_2][a_2,a_3]^2_\uparrow
 &12 &[a_0,a_1]^2_\uparrow[a_3,b]\\
13 &[a_0,a_1][a_1,a_2][a_2,b]
 &14 &[a_0,a_1][a_2,b]^2_\uparrow\\
15 &[a_1,a_2][a_2,a_3][a_3,b]
 &16 &[a_2,a_3]^4_\uparrow\\
17 &[a_2,a_3]^3_\uparrow[a_3,b]
 &18 &[a_2,a_3]^2_\uparrow[a_3,b]^2_\uparrow\\
19 &[a_1,a_2][a_3,b]^3_\uparrow
 &20 &[a_2,a_3][a_3,b]^4_\uparrow\\
21 &[a_3,b]^6_\uparrow.&&
\end{array}
\end{equation}
一个含 $r$ 个素数的区域，其被加项为 $Q_{q_1\cdots q_r}(q_{r-1};n)$。\eqref{sv:switch-domains} 中的并置表示相互独立的连续区带，仅在每个指定的有序块内部排序。所用不等式是
\begin{equation}\label{sv:twenty-one-weight}
 5\Phi_s\le\Gamma_1-\Gamma_2-\Gamma_3-\Gamma_4
                       +\sum_{i=5}^{21}\Gamma_i+o(\Theta_Q).
\end{equation}
这是 \cite[引理~4.2]{Wu2004} 的第二个加权不等式；以下验证明确其完整范围与重数。

对无平方因子且 $R^{a_0}$-粗糙的商，用 $A,B',C,D$ 分别计数其在四个区带
\[
[a_0,a_1),[a_1,a_2),[a_2,a_3),[a_3,b)
\]
中的素因子。（本验证中的符号 $B'$ 是区带计数，不是带标记前缀 $B$。）逐点看，前四项为
\[
 \Gamma_1=4+\one_{A=0},\qquad
 (\Gamma_2,\Gamma_3,\Gamma_4)
       =(A+B'+C+D,A+B',A+B'+C).
\]
其余项的值如下，其中 $x_+=\max(x,0)$：
\begin{equation}\label{sv:band-counts}
\begin{array}{c|l@{\quad}c|l}
5&\binom{A+B'}2&6&AC\\
7&(A-1)_+&8&\one_{A\ge1}B'\\
9&\one_{A=0}(B'+C-2)_+
 &10&\one_{A=0,B'\ge2}(C+D)\\
11&\one_{A=0,B'=1}(C-1)_+
 &12&\one_{A\ge2}D\\
13&\one_{A=1,B'\ge1}(C+D)
 &14&\one_{A=1,B'=0}(C+D-1)_+\\
15&\one_{A=0,B'=1,C\ge1}D
 &16&\one_{A=B'=0}(C-3)_+\\
17&\one_{A=B'=0,C\ge3}D
 &18&\one_{A=B'=0,C=2}(D-1)_+\\
19&\one_{A=0,B'=1,C=0}(D-2)_+
 &20&\one_{A=B'=0,C=1}(D-3)_+\\
21&\one_{A=B'=C=0}(D-5)_+.&&
\end{array}
\end{equation}
事实上，在倒数第二个选定素数处对商作筛，迫使此前各选定素数为 $n$ 的最小素因子，由此得到 \eqref{sv:band-counts} 的每一项。\eqref{sv:twenty-one-weight} 的带符号右端分别为
\[
\begin{array}{c|l}
A\ge2&(A+B'-2)(A+B'-3)/2+(A-2)C\\
A=1,\ B'\ge1&(B'-1)(B'-2)/2\\
A=1,\ B'=0&(1-C-D)_+\\
A=0,\ B'\ge2&(B'-2)(B'-3)/2\\
A=0,\ B'=1,\ C\ge1&0\\
A=0,\ B'=1,\ C=0&(2-D)_+\\
A=B'=0,\ C\ge3&0\\
A=B'=0,\ C=2,1,0&(1-D)_+,\ (3-D)_+,\ (5-D)_+.
\end{array}
\]
这些值均非负，且在 $A=B'=C=D=0$ 时为五；左端就是最后这一情形的示性函数乘以五。在 $R^{a_0}$-粗糙支撑之外，每一被加项均为零，故逐点不等式成立。在固定正指数盒上，被略去的平方和前缀交叉项总质量为 $N^{1-\epsilon+o(1)}$，其中 $\epsilon>0$ 固定；经过有限带标记平均，得到所述误差。

\paragraph{逐项估计与普通基础项。}
$\Gamma_1$ 的两个部分具有上界乘子 $4A_+(s_0)+A_+(\phi_1)$。三个负项具有下界乘子
\begin{equation}\label{sv:negative-singleton-base}
 \mathcal B(c)=\int_{a_0}^{c}
   \frac{A_-(s_0(1-u))}{u(1-u)}\,du,
 \qquad c=b,a_2,a_3.
\end{equation}
可用条件使这些区间上的 $s_0(1-u)$ 位于 $(2,4]$，故 $A_-(s_0(1-u))=\log(s_0(1-u)-1)$。合并的单素数与素数对系数满足父状态的完整平方及选定补集检验。特别地，$D_5^+$、$D_6$、$D_7$、$D_8$ 相对于 $du\,dv/(uv)$ 具有普通上界核 $(1-u-v)^{-1}$。子节点自变量低于 $2$ 时，与更小筛截点比较仍给出同一核，因为直到 $3$ 均有 $A_+=1$。在 $D_5^-=\{a_0<u<v<a_2:v+2u\ge1\}$ 上不使用普通素数对估计。非负乘积估计给出有效核 $Kc_\omega s_0/4$，因为其粗糙自变量至少为 $s_0(2-2a_2)>2$。定义
\begin{align}\label{sv:pair-base}
 \mathcal U_5&=
  \iint_{D_5^+}\frac{du\,dv}{uv(1-u-v)}
   +\frac{Kc_\omega s_0}{4}\iint_{D_5^-}\frac{du\,dv}{uv},\notag\\
 \mathcal U_i&=\iint_{D_i}\frac{du\,dv}{uv(1-u-v)}
                  \qquad(i=6,7,8).
\end{align}
这样，以非负上界估计替换不具普通估计支撑的角部，不再同时保留其普通质量或反馈。

对具有 $r$ 个素数的换位区域 $\mathcal D_i$，令 $w=u_{r-1}$、$U=u_1+\cdots+u_r$。在实际父节点比值 $\alpha$ 处，其归一化未筛质量为
\begin{equation}\label{sv:full-switched-mass}
 \mathcal M_i(\alpha)=
 \int_{\substack{\boldsymbol u\in\mathcal D_i\\U\le\alpha-w}}
 \frac{\omega((\alpha-U)/w)}{w\prod_j u_j}\,d\boldsymbol u
 +\int_{\substack{\boldsymbol u\in\mathcal D_i\\U=\alpha}}
                     \frac{du_1\cdots du_{r-1}}{\prod_j u_j}.
\end{equation}
曲面积分中以导数一消去一个坐标。该曲面计数商为 $1$ 的情形，与普通 Buchstab 项分开。固定正盒上的素数求和及粗糙数公式给出 \eqref{sv:full-switched-mass}，按父节点单位归一化的上界计数为 $K\mathcal M_i(\alpha)/4$。

以下显式界对 $\alpha\ge2$ 一致。令 $\ell_{ij}=\log(\phi_i/\phi_j)$、$\ell_{01}=\log(s_0/\phi_1)$、$\ell_{3b}=\log(\phi_3/s)$。对 $i=9,\ldots,15$，取 $C_i=\max\{c_\omega,1/\max(1,m_i)\}$，其中
\[
\begin{array}{c|l@{\quad}c|l}
i&m_i&i&m_i\\ \hline
9&2\phi_3-3&10&2\phi_2-2-\phi_2/s\\
11&2\phi_3-\phi_3/\phi_2-2
 &12&2\phi_1-2-\phi_1/s\\
13&2\phi_2-\phi_2/\phi_1-1-\phi_2/s
 &14&2s-s/\phi_1-2\\
15&2\phi_3-\phi_3/\phi_2-1-\phi_3/s.&&
\end{array}
\]
这些量分别是各区域粗糙自变量的下界。由 \eqref{sv:omega-envelope}，得
\begin{align}\label{sv:switch-mass-bounds}
 \widehat M_9&=C_9\{(\phi_1+\phi_3)\ell_{13}-2(\phi_1-\phi_3)\},\notag\\
 \widehat M_{10}&=C_{10}\{\phi_1-\phi_2(1+\ell_{12})\}
                                          \log(\phi_2/s),\notag\\
 \widehat M_{11}&=C_{11}\ell_{12}\{\phi_2(\ell_{23}-1)+\phi_3\},\notag\\
 \widehat M_{12}&=C_{12}\{s_0-\phi_1(1+\ell_{01})\}\ell_{3b},\notag\\
 \widehat M_{13}&=C_{13}\ell_{01}(\phi_1-\phi_2)\log(\phi_2/s),\notag\\
 \widehat M_{14}&=C_{14}\ell_{01}
                       \{\phi_2(\log(\phi_2/s)-1)+s\},\notag\\
 \widehat M_{15}&=C_{15}\ell_{12}(\phi_2-\phi_3)\ell_{3b},\notag\\
 \widehat M_{16}&=\phi_2\ell_{23}^4/24,
 &\widehat M_{17}&=\phi_2\ell_{23}^3\ell_{3b}/6,\notag\\
 \widehat M_{18}&=\phi_3\ell_{23}^2\ell_{3b}^2/4,
 &\widehat M_{19}&=\phi_3\ell_{12}\ell_{3b}^3/6,\notag\\
 \widehat M_{20}&=\phi_3\ell_{23}\ell_{3b}^4/24
                         +\phi_2\ell_{3b}^4/24,\notag\\
 \widehat M_{21}&=\phi_3\ell_{3b}^6/720+\phi_3\ell_{3b}^5/120.
 \end{align}
对 $i\le19$，指定指数之和至多为 $4/s<2\le\alpha$，故没有单位商曲面。对 $i=20$，在曲面上消去第一个变量并使用 $1/u_1\le\phi_2$；其余四变量在 $[a_3,b]$ 中有序排列，得到 $\widehat M_{20}$ 的第二个被加项。对 $i=21$，消去末变量，使用 $1/u_6\le\phi_3$，并对其余五个有序变量积分。\eqref{sv:switch-mass-bounds} 的其余被加项由在所示区带上积分 $1/(w\prod u_j)$ 得到，其中 $9\le i\le15$ 时取 $\omega\le C_i$，其余情形取 $\omega\le1$。因此 $\mathcal M_i(\alpha)\le\widehat M_i$。

普通基础项现在成为显式泛函
\begin{equation}\label{sv:two-step-base}
 J_2^0=\frac15\left\{
 5A_+(s)-4A_+(s_0)-A_+(\phi_1)
 +\mathcal B(b)+\mathcal B(a_2)+\mathcal B(a_3)
 -\sum_{i=5}^8\mathcal U_i
 -\frac K4\sum_{i=9}^{21}\widehat M_i\right\}.
\end{equation}
在 \eqref{sv:twenty-one-weight} 中，对正项代入上界，对负项代入下界。$\Gamma_1$ 两个部分中的改进量，从其上界估计中减去 $4H(s_0)+H(\phi_1)$。三个单素数下界改进量，从各负项中减去相应积分；四个良好素数对的上界改进量，从各自非负上界估计中减去素数对积分。除以 $5$ 后，这些恰为 \eqref{sv:two-step-feedback}。所以
\begin{equation}\label{sv:two-step-admissibility}
 \Phi_s\le\{A_+(s)-J_2^0-\mathcal P_2[H,h]+o(1)\}\Theta_Q,
\end{equation}
其中 $\mathcal P_2[H,h]$ 为 \eqref{sv:two-step-feedback} 的表达式。由此，从数学上明确的基础项和非负加权的前一代改进量，证明一个有限双步作用的可用性。

可以用更精确的上界替换 \eqref{sv:pair-base} 与 \eqref{sv:full-switched-mass} 中的正代价。该上界必须对每个允许的 $\alpha\ge2$ 覆盖同一带标签族，并包含商为 $1$ 的曲面。在 \eqref{sv:two-step-base} 中仅替换该代价一次，再按所需方向估计所得显式泛函，即得相应的转移常数。这样保留了两步不等式的符号论证。

\subsection{可用带标记族上的有限初值}

令
\begin{equation}\label{sv:seed-parameters}
 S=\frac{11}{5},\qquad s_0=\frac{227}{50},\qquad
 (\phi_1,\phi_2,\phi_3)
 =\left(\frac{353}{100},\frac{29}{10},\frac{61}{25}\right),
 \qquad K_{\rm seed}=\frac{681}{100}.
\end{equation}
由于 $K<K_{\rm seed}$，命题~\ref{prop:current-positive-product} 也给出常数为 $K_{\rm seed}$ 的非负上界估计。令 $\mathcal J_{\rm seed}$ 为 \eqref{sv:two-step-base} 在参数 \eqref{sv:seed-parameters} 处的右端，并将 $s=S$、$K$ 替换为 $K_{\rm seed}$。特别地，其第五项代价采用完整的分拆代价 \eqref{sv:pair-base}，并计入 \eqref{sv:switch-mass-bounds} 的两个单位商曲面项。

\begin{lemma}[有限初值]\label{sv:finite-seed}
对满足 $\alpha\ge2$ 的可用带标记族，在增强情形另假定 $L/S<\chi$，则
\[
 \Phi_S\le\left(1-\frac{2327}{100000}+o(1)\right)\Theta_Q.
\]
只要状态允许比较截点 $N^{L/S}$，同一上界改进量在每个 $1\le w\le S$ 都可使用。结论对每个固定的此类有限族一致，并由初始改进量 $1/50$ 之后的四个加权步骤推出。
\end{lemma}

\begin{proof}
二十一项不等式中所有新增素数的指数至多为 $L/S$，因而低于保留的上限，使最大的已指定位置、选定子集及任何第三位置界均仍有效。单素数剩余平方条件为 $2u<1$，素数对平方条件为 $2v<1$、$v+2u<1$，并在 $D_5^+,D_6,D_7,D_8$ 上成立。将相应剩余平方乘以 $B$，即得每个新增完整平方。原有选定位置出现后，新的选定补集平方等同于完整平方。普通剩余原子在其实际筛级 $Q/(Bq)$ 或 $Q/(Bq_1q_2)$ 上计算。其余部分 $D_5^-$ 保留并承担非负代价，不对该处使用带符号分布结论。经典情形下完整普通模数低于 $Q$，低于二分之一水平的绝对余项给出同一保守泛函。换位项只需要 $\alpha\ge2$ 及固定的正素数截点。

显式泛函的定向包围为
\begin{equation}\label{sv:seed-base-certificate}
 \mathcal J_{\rm seed}\ge
 \frac{11672890052052913}{576460752303423488}
 >\frac{20249}{10^6}>\frac1{50}.
\end{equation}
该不等式的有限证书由附属数据中的 \texttt{seed\_certificate} 指定。其普通积分为 \eqref{sv:negative-singleton-base}、\eqref{sv:pair-base}，换位质量采用 \eqref{sv:switch-mass-bounds} 的显式上界。普通二重积分先对一个坐标积分，所得一维对数积分用有理 Darboux 区间及 \eqref{num:log-enclosure} 的对数余项包围。在 $\mathcal J_{\rm seed}$ 的定义中，正代价取上端点，负代价取下端点。因此，\eqref{sv:seed-base-certificate} 在使用后续数组递推之前即证明初始改进量 $J_0=1/50$。

假设已在这些可用族上证明 $[1,S]$ 上的上界改进量 $J>0$。非负性和截点比较给出
\[
 H_J(w)=J\one_{[1,S]}(w).
\]
同一总筛级上的下界 Buchstab 恒等式给出先前已证明的备选
\begin{equation}\label{sv:seed-lower-child}
 h_J(w)=J\log\frac{S}{w-1}\,
                    \one_{[2,1+S]}(w).
\end{equation}
归一化不依赖于 $w$，且 $1\le w\le S$ 时有 $A_+(w)=1$。为核验 \eqref{sv:seed-lower-child} 内的比较，由 $v\in[w-1,S]$ 标记的 Buchstab 子节点，其接入指数为 $L/(v+1)$，剩余指数为 $Lv/(v+1)$，初值比较指数满足
\[
 \frac{Lv}{S(v+1)}\le\frac{L}{v+1}.
\]
因此它低于新增素数及先前稳定上限。由于 $v\ge1$，新增平方条件成立。下界子节点自身无需允许自变量 $S$；初值实际应用于它的上界 Buchstab 子节点。

将 $H_J,h_J$ 代入 \eqref{sv:two-step-feedback}。由于 $s_0,\phi_1>S$，未改变前缀的改进量为零。记 $u_*=67/227$、$b=1/S$，并令
\[
 E_i=\begin{cases}
 D_i\cap\{u+v\ge117/227\},&i=5^+,6,\\
 D_i\cap\{v\ge1-16u/5\},&i=7,8.
 \end{cases}
\]
反馈恰为 $\lambda_{\rm seed} J$，其中
\begin{align}\label{sv:seed-feedback-integral}
 \lambda_{\rm seed}=\frac15\bigg\{&
 \sum_{c\in\{b,a_2,a_3\}}
 \int_{u_*}^{c}
 \frac{\log\{S/[s_0(1-u)-1]\}}{u(1-u)}\,du\notag\\
 &+\sum_{i=5^+,6,7,8}\iint_{E_i}
                     \frac{du\,dv}{uv(1-u-v)}\bigg\}.
\end{align}
$i=5^+,6$ 的支撑条件为 $s_0(1-u-v)\le S$；$i=7,8$ 的支撑条件为 $(1-u-v)/u\le S$。全部保留的素数对子节点自变量都超过 $1$，单素数下界自变量都超过 $2$。$D_5^-$ 不提供反馈。

\eqref{sv:seed-feedback-integral} 中全部边界均为有理仿射边界。对竖直切片 $l(u)<v<r(u)$，精确素数对积分为
\[
 \frac1{u(1-u)}
 \log\frac{r(u)(1-u-l(u))}{l(u)(1-u-r(u))}.
\]
对 $E_{5^+}$，连续 $u$ 区间的端点为 $50/227,117/454,19/58,1/3$。区域 $E_6$ 是其完整矩形。$E_7$ 的非空支撑从 $1265/5648$ 开始，并在 $5/21$ 改变下边界；$E_8$ 的下边界在 $1265/5648$ 改变。对这些区域的有理 Darboux 包围给出
\begin{equation}\label{sv:seed-feedback-certificate}
 \lambda_{\rm seed}\ge
 \frac{75057365825327405}{576460752303423488}
 >\frac{13}{100}.
\end{equation}
精确端点、细分和来源对应，与 \eqref{sv:seed-base-certificate} 一同由 \texttt{seed\_certificate} 数据指定。

由 \eqref{sv:two-step-admissibility}、\eqref{sv:seed-base-certificate}、\eqref{sv:seed-feedback-certificate} 得到
\[
 J_{n+1}=\frac{20249}{10^6}+\frac{13}{100}J_n.
\]
每一步均先由 $J_n$ 推出 \eqref{sv:seed-lower-child}，再应用加权不等式。从 $J_0=1/50$ 出发，四个精确值为
\[
 \frac{22849}{10^6},\quad
 \frac{2321937}{10^8},\quad
 \frac{232675181}{10^{10}},\quad
 \frac{23273777353}{10^{12}}>\frac{2327}{100000}.
\]
接入素数的操作只有有限次。所有紧区域选择固定后，其指数始终为正，因此每一步都适用非负乘积估计及分布余项。有限多个误差为 $o(\Theta_Q)$，按所述极限次序可由最后的严格剩余量吸收。
\end{proof}

\paragraph{在根节点与后继中的应用。}
对 $U_3,U_4$，已指定素数的指数为 $t>{a_{\rm s}}$，剩余指数为 $L=v(t)-t$，上界自变量为 $\max(2,L/{a_{\rm s}})$。在增强范围 $t<57/224$，最大素数分支满足 $2t<v(t)\le(1+t)/2$，故稳定上限可取在 $t$ 以下并任意接近它。实际截点指数 $\min({a_{\rm s}},L/2)\le {a_{\rm s}}<t$ 低于该上限。每次使用初值时，比较指数 $L/S$ 不大于实际截点。经典尾部采用低于二分之一水平的接口。对 $U_{10},U_{11}$，同一论证使用最小已指定指数 $w$ 和实际截点 $\min(w,L/2)$；原有选定子集及前三个位置已固定。对下界素数对项 $L_5,L_6$，初值只出现在 \eqref{sv:seed-lower-child} 内。首次接入素数至多为 $N^{a_{\rm s}}$，而由前一显示不等式，全部初值比较素数都低于它。这些比较也保留等于 ${a_{\rm s}}$ 的第三位置界。

更一般地，\eqref{sv:one-step-operator} 或 \eqref{sv:two-step-feedback} 中每个有限作用都保留总筛级 $Q$、所有原有平方条件及稳定上限。新单素数满足 $2u<1$；普通部分的新素数对满足 $2v<1$、$v+2u<1$。其实际截点不大于最小新增素数。式 \eqref{sv:child-state} 表明其比值仍至少为 $2$。由归纳，初值及其 Buchstab 推论始终处于可用族内。没有把初值改进量直接代入空前缀下界起点。

有限归纳的每个右端均使用前一代。在每个状态上取备选估计的最大值；当子自变量跨越节点时，在所有相邻闭单元上取最小值。因此，每一步都给出同一非负带标记计数的界，且只使用有限多步。

下面给出显式的共同逼近界。非负算子质量满足
\begin{align*}
 \|\mathcal B\|&\le\log5<2,\\
 \|\mathcal W_1\|&\le1+\tfrac12\log4<2,\\
 \|\mathcal W_2\|&\le1+\tfrac35\log4
                 +\tfrac45\,125(3/10)^2<12.
\end{align*}
最后一个估计用到 $1/5\le u,v\le1/2$、$1-u-v\ge u\ge1/5$，故四个区域上各自的素数对核至多为 $125$。未变前缀的系数和为一。换位非负代价属于作用的基础项，不是额外的反馈边。若每个局部定向逼近误差至多为 $\delta$，则
\begin{equation}\label{sv:feedback-error}
 e_{j+1}\le\delta+12e_j,\qquad
 e_j\le\delta\frac{12^{j+1}-1}{11}.
\end{equation}
十一项泛函的系数绝对值在乘因子 $1/4$ 之前之和为 $14$，故根部范数为 $7/2$。

在深度至多为 $440$ 时，每个剩余筛级至少为 $\lambda_*101^{-440}>0$，新增素数指数至少为此数除以 $6$。对每个固定非负父节点递归细分新增指数，使误差不超过其归一化质量乘以 $\delta$。真正的上限与节点边界为零测集；结构性截断及继承恒等区域保留。分母一致远离零。由 \eqref{sv:feedback-error} 可得一组有限查询，总误差任意小。仅在该组查询固定之后，才令 $N$ 增大。正质量顶点具有正的极限 $\Theta_Q/H_N$，故绝对分布误差化为 $o(\Theta_Q)$；零质量部分用绝对误差处理。

\section{端点估计与极限次序}\label{app:cross-review}

尖锐素数引理（引理~\ref{lem:current-sharp-prime}）适用于上面构造的同一系数。下面记录使其可用于带符号下界权的绝对界。若 $|\rho_N^{\rm low}(q)|\le C$，其无平方因子支撑为 $q\le N^\theta$，且 $\theta<1$、$\delta=(\log N)^{-B}$，则端点贡献至多为
\begin{equation}\label{sv:sharp-boundary}
 \sum_{\substack{p\text{ 为素数}\\0<|p-N|\le\delta N}}
          \sum_{q\mid N-p}|\rho_N^{\rm low}(q)|
 \le2C\sum_{1\le n\le\lceil\delta N\rceil}\tau(n)
 \ll\delta N\log N+\log N.
\end{equation}
当 $N>2$ 为偶数时，差值 $p-N=0$ 不可能出现。因 $\sum_{q\le N^\theta}1/\varphi(q)\ll\log N$，精确互素均值的边界也有同一界。$2N^{1-\kappa_0}$ 以下的光滑下截尾，在等差数列和中的代价为 $N^{1-\kappa_0+o(1)}$，在均值中的代价为 $O(N^{1-\kappa_0}\log N)$，均具有固定幂次节省。选择 $B>A+3$，即可使全部边界项小于规定的对数误差。

对精确互素素数均值，设 $\omega_{\mathrm{pf}}(q)$ 为 $q$ 的不同素因子个数。每个素因子低于 $N$，无平方因子性给出
\begin{equation}\label{sv:coprime-prime-error}
 \sum_q\frac{|\rho_N^{\rm low}(q)|\omega_{\mathrm{pf}}(q)}{\varphi(q)}
 \ll\sum_{\ell\le N^\theta}\frac1{\ell-1}
                  \sum_{e\le N^\theta/\ell}\frac1{\varphi(e)}
 \ll\log^2N.
\end{equation}
$\ell$ 上的和只对素数求取；即使扩大到全部整数，也足以得到所述粗界。由除数界，从原始序列中移去 $N$ 的 $O(\log N)$ 个素因子的贡献为 $N^{o(1)}$。这些估计说明可以在 \eqref{sv:physical-lower} 中使用尖锐的原始素数计数，而无需从非负乘积上界常数 $K$ 推出素数下界。

上、下界计数遵循第~\ref{sec:lower-limits} 节的参数次序。先固定矩阶数，再按~\eqref{coef:parameter-order} 选择 $\kappa$；边界指数仍满足上述 $B>A+3$。
\section{有限数值估计}\label{app:numeric}

本附录说明决定引理~\ref{lem:numerical-endpoints} 中两个常数的有限量。构造分别处理积分包围、有限改进量递推及最终有理相减，使每项贡献都有明确的数学定义和估计方向。

\subsection{补充数据}
有限数据收录于 \emph{Theorem (1+1.75) on the Goldbach Conjecture --- Numerical Supplement}，版本 20261106~\cite{ChanSupplement}，包括 $106$ 个仿射转移、深度 $248$、$312$、$430$ 的 $200$ 单元改进量数组，以及 $17572$ 个配对矩形。\filepath{finite-construction-data.json} 中的有理量表示为整数之商。字段 \texttt{upper\_saving\_bounds}、\texttt{lower\_root\_bounds}、\filepath{lower_integral_bounds} 给出定向有理端点。\filepath{certificate-map.json} 标识其来源证书及相应的区域、细分和区间记录。数学估计与证书字段的对应关系见表~\ref{num:certificate-table}。

\hypertarget{num:verification-scope}{}
\paragraph{核验命令与输入。}
下列命令从数值包根目录运行，使用 Python 3.11 或更高版本，并保留断言检查。六阶计算还使用 NumPy。表中区分从已有证书重建数据、对有限构造作精确核验，以及重算指定记录的数值界。通过 \texttt{-{}-output} 或 \texttt{-{}-output-dir} 指定的输出文件均写到数值包之外。

\begingroup
\renewcommand{\LTcaptype}{}
\stepcounter{LT@tables}
\small
\setlength{\tabcolsep}{4pt}
\begin{longtable}{@{}>{\raggedright\arraybackslash}p{.40\textwidth}>{\raggedright\arraybackslash}p{.29\textwidth}>{\raggedright\arraybackslash}p{.25\textwidth}@{}}
\toprule 入口 & 核验或重算的量 & 保留的数值输入 \\
\midrule
\endhead
\bottomrule
\endfoot
\filepath{verify-package.py} & 文件集合、大小和已记录的 SHA-256 标识。 & 版本对应的 \filepath{MANIFEST.json}。\\[3pt]
\filepath{finite-data/reproduce-data.py} \texttt{-{}-check} & 从源证书精确提取有限数据及字段对应。 & 已提供的源证书值。\\[3pt]
\filepath{finite-data/verify-finite-data.py} & 有限改进量递推；配对单元的几何、质量与损失控制式；$\Gamma_5$ 分解；密度和与严格标量余量。 & \filepath{certificate-map.json} 标明的动作系数及积分端点。\\[3pt]
\filepath{replay-records.py} & 有序覆盖的索引、已存整数下界之和及代价恒等式。 & 有序记录中已存的整数下界。\\[3pt]
\filepath{replay-ordered-cover.py} & 网格 48 基础覆盖每行选定方法的下界；加 \texttt{-{}-grid 32} 时处理其前一覆盖。 & 已记录的方法选择及第 248 代改进量表。\\[3pt]
\filepath{verify-sixth-domains.py} & 已存多面体单元上的仿射筛函数自变量范围。 & 已记录的顶点、坐标映射及上限。\\[3pt]
\filepath{replay-sixth-rank.py} & 网格 400 六阶立方体的贡献及其和，以及上一项的多面体范围检验。 & 已提供的立方体记录和多面体记录。\\
\end{longtable}
\endgroup

递推与配对单元核验使用表~\ref{num:certificate-table} 所列动作及积分输入。其原始包围要求仍由本附录及附录~\ref{app:upper}--\ref{app:lower} 中对应的积分和规则规定。\filepath{certificate-map.json} 的 \texttt{enclosure\_obligations} 条目逐族列出当前断言、完整记录和已提供的算法源文件。数值包 README 说明记录的生成方法；其中 \filepath{witnesses/generator-0041.py} 生成有序基础覆盖，\filepath{witnesses/generator-0045.py} 生成六阶立方体记录。重算时保留的第 248 代数组及其它数值输入通过各自条目核验，使每项核验均具有明确的输入集合。命令选项见数值包 README，有限数据的定点精度在下文规定。

\subsection{精确标量端点}
上界密度估计使用的精确值为
\begin{align}
 E_{\rm up}&=\frac{9979702576710911234850619807816847864907359959}
 {1475913550962201227946780619099353215139840000},
 \label{eq:current-E}\\
 u&=\frac{145007961227048312210593427530128000784681}
 {11807308407697609823574244952794825721118720000}.
 \label{eq:current-u}
\end{align}
对下界构造，在计入配对纤维增量之前，令
\begin{align}
 R_{\rm base}&=
 \frac{1345575129530003739571681432301188857312754737}
 {685738127494661151198444592354557952000000000},
 \label{eq:base-R}\\
 B_{\rm base}&=
 \frac{2504748318091942808229694219379799307550987}
 {5485905019957289209587556738836463616000000000}.
 \label{eq:base-B}
\end{align}
$u$ 和 $B_{\rm base}$ 各出现两次，分别支付有限逼近和系数实现的不同误差。求积误差已经包含在定向积分端点中。附属字段 \texttt{scalar\_bounds} 记录最终精确和及严格剩余量。

\subsection{定向算术}
证明上界时采用向外包围，证明下界时采用向内估计。更具体地，在分母 $\mathsf Q=2^{b_{\rm fp}}$ 下，有理区间 $[x,y]$ 表示为 $[\lfloor \mathsf Q x\rfloor,\lceil \mathsf Q y\rceil]$。加减在整数层面精确进行。乘法先取四个端点乘积的最小值和最大值，再除以 $\mathsf Q$，按规定方向取整。除法只在验证分母区间为正后进行。

对数计算中，令 $z=(x-1)/(x+1)$。若 $1\le x\le2$，则
\begin{equation}
 2\sum_{j=0}^{n-1}\frac{z^{2j+1}}{2j+1}
 \le\log x\le
 2\sum_{j=0}^{n-1}\frac{z^{2j+1}}{2j+1}
 +\frac{2z^{2n+1}}{(2n+1)(1-z^2)}.
 \label{num:log-enclosure}
\end{equation}
这由 $2\operatorname{arctanh}z$ 的正项幂级数求和，并将其余分母从下方以 $2n+1$ 控制得到。对 $n=96$，余项至多为 $9/(4\cdot193\cdot3^{193})$。范围缩减配合 $\log2$ 的单独包围，将公式推广至这里用到的每个正有理自变量。配对单元计算的级数分母为 $2^{224}$，返回区间的分母为 $2^{160}$。因此，不等式方向不由任何浮点对数值决定。

\subsection{延迟函数与积分的包围}
递推分子较为方便，记
\[
 A_+(s)=\frac{sF(s)}{2e^\gamma},\qquad
 A_-(s)=\frac{sf(s)}{2e^\gamma}.
\]
初始分支分别为 $1\le s\le3$ 时的 $A_+(s)=1$、$s\le2$ 时的 $A_-(s)=0$，后续分支满足
\[
 A_+'(s)=\frac{A_-(s-1)}{s-1},\qquad
 A_-'(s)=\frac{A_+(s-1)}{s-1}.
\]
每个分支的方程只涉及前一个单位区间。因此，可在有理网格上依次生成有理区间包围，并在每个整数分点处分开。对被积函数使用相应区间界，即得积分上下界。弦线界和曲率界可以在不改变方向的前提下收紧包围。任意查询都覆盖其完整支撑；端点恰为网格节点时，包括相邻两单元。$s\omega(s)$ 的类似构造使用 $(s\omega(s))'=\omega(s-1)$，其在 $[1,2]$ 上的初始分支为 $\omega(s)=1/s$。

弦线积分可以不借助数值求积实现。在 $[j/n,(j+1)/n]$ 上，分子 $V$ 的仿射端点插值除以自变量后的积分为
\[
 \bigl((j+1)\log(1+1/j)-1\bigr)V(j/n)
 +\bigl(1-j\log(1+1/j)\bigr)V((j+1)/n).
\]
两个权都是非负仿射基函数相对于 $du/u$ 的积分，因此为正。$F$、$f$ 的单调性说明，$A_+$ 在非恒定分支上为凸，$A_-$ 为凹。在一个延迟单元上，其分子曲率的绝对值至多为 $1/(j/n-1)$，故积分后的弦线误差至多为
\[
 \frac{1}{12nj(j-n)}\qquad(j>n).
\]
对弦线下界减去这一量，对弦线上界加上这一量。权本身具有交错有理级数
\[
 (j+1)\log(1+1/j)-1
  =\sum_{k\ge1}\frac{(-1)^{k+1}}{k(k+1)j^k},\qquad
 1-j\log(1+1/j)
  =\sum_{k\ge1}\frac{(-1)^{k+1}}{(k+1)j^k}.
\]
第四和第五部分和包围精确权。因此，即使直接相减两个几乎相等的对数在数值上不便，也仍有有理包围。

有限积分采用有理盒或有理多面体。对仿射不等式定义的多面体，精确求解每个可能的活动基，并对所得顶点检验每条原始不等式。先验证有界性及面之间的关联，再构造拉取三角剖分。随后多项式矩可以精确计算：在重心坐标为 $\lambda_0,\ldots,\lambda_d$ 的 $d$ 维单形 $T$ 上，
\begin{equation}
 \int_T\prod_{j=0}^d\lambda_j^{\alpha_j}\,dx
 =d!\operatorname{vol}(T)
   \frac{\prod_{j=0}^d\alpha_j!}{(d+\sum_j\alpha_j)!}.
 \label{num:simplex-moment}
\end{equation}
其余核因子的有理上下界，结合这些矩及 \eqref{num:log-enclosure}，给出所需积分包围。对负项代价的上界积分，边界盒必须保留并估计其全部贡献。对非负节省量，证明略去贡献非负后，可以省去该盒。两条规则的符号不同，须分别使用。

\subsection{上界密度}
写作
\[
 U_0=\frac4\Theta+I_1+I_{3,c}+I_{3,t},\qquad
 E_{\rm up}=U_0-\sum_{j=1}^{22}G_j.
\]
其中 $I_{3,c}$ 是末坐标至少为 $1/100$ 的区域积分，$I_{3,t}$ 控制完整补尾。令 $D=(F-f)/(2e^\gamma)$，由其延迟方程，利用 $D$ 的非负性和单调性，得
\[
 D'(s)=-\frac{D(s)+D(s-1)}s\le-\frac{2D(s)}s
 \quad(s\ge3),
\]
故 $s\ge10$ 时有 $D(s)\le D(10)(10/s)^2$。若 $\delta=1/100$、$\rho_0=19/96$，上界树的上限条件保证剩余指数至少为 $\rho_0$。当 $0<z<\delta$ 时，其延迟自变量至少为 $(\rho_0-z)/z\ge10$。将其余两个有序坐标扩大到 $z\le y\le x\le\rho_0$，得到完整尾界
\begin{equation}
 I_{3,t}\le
 \frac{200D(10)\delta}{(\rho_0-\delta)^2}
 \left\{\log^2\frac{\rho_0}{\delta}
        +2\log\frac{\rho_0}{\delta}+2\right\}.
 \label{num:complete-tail}
\end{equation}
这里有序 $x,y$ 积分为 $\tfrac12\log^2(\rho_0/z)$，余下关于 $0<z<\delta$ 的积分是初等积分。特别地，$z=0$ 处的对数奇性被完整积分，而非略去。第一个积分用正项级数及其几何余项包围。紧的第三积分采用不交分割覆盖，并对边界盒保留保守界。

表~\ref{num:upper-table} 给出非负节省量 $G_j$ 的便于阅读的下界，其整数条目满足 $G_j\ge m_j10^{-14}$，完整有理界见补充数据。同一区域上的连续变更均相对于各自父项计算。特别地，最后两行只给出相对于该积分前一估计的曲率及混合增量。

\begin{longtable}{@{}r p{.62\textwidth} r@{}}
\caption{上界密度节省量的有理下界 $10^{-14}m_j$。}
\label{num:upper-table}\\
\toprule
$j$ & 节省量 & $m_j$\\
\midrule
\endfirsthead
\toprule $j$ & 节省量 & $m_j$\\\midrule
\endhead
1 & 四阶上限替换 & 3507783546602\\
2 & 五阶父单元细分 & 1005163763573\\
3 & 不交五阶区域 & 136378695993\\
4 & 六阶上限替换 & 96281257455\\
5 & $\mu_0$ 以上的二阶备选 & 4209783419\\
6 & $\mu_0$ 以上的三阶备选 & 3908979560\\
7 & 六阶对数矩 & 1525398138\\
8 & 六阶递减区域 & 187855687\\
9 & 六阶递增上限差 & 387157212\\
10 & 六阶递减上限差 & 236377744\\
11 & 六阶裁剪下窗口 & 846201484\\
12 & 六阶递减下窗口 & 500857555\\
13 & 六阶局部上限修正 & 50721206\\
14 & 六阶非单点区域 & 11743733\\
15 & 八阶径向界 & 94105689\\
16 & 八阶递增区域 & 1543640141\\
17 & 八阶递减区域 & 590010094\\
18 & 八阶不等父上限矩 & 1304343681\\
19 & 八阶第三坐标矩 & 357686854\\
20 & 八阶前缀受限矩 & 37235787\\
21 & 八阶曲率增量 & 97705\\
22 & 八阶混合增量 & 607514\\
\bottomrule
\end{longtable}

上界包围使用 $3699127$ 个四阶盒、细分为 $18728512$ 个子盒的 $292633$ 个五阶父盒、独立五阶区域中的 $1763420$ 个盒，以及 $632104$ 个六阶立方体。八阶区域分为大小为 $30$、$15$、$13$、$29$ 的互不相交族。不同阶之间的筛级差作望远镜求和。在每个父区域上，将其本身的界与完整子项和比较，取有效节省下界的最大值。

在同一分母 $10^{14}$ 下，四个基础项的上界分子分别为
\[
 668896321070235,\quad76020940667,\quad
 11960259849166,\quad12313775.
\]
令 $E_{\rm up,14}$ 为其和减去表~\ref{num:upper-table} 之和，再除以 $10^{14}$。精确整数运算给出
\[
 E_{\rm up}\le E_{\rm up,14},\qquad E_{\rm up,14}+2u<K,
\]
其中 $u$ 见 \eqref{eq:current-u}；最后一个差事实上大于 $9.65\cdot10^{-8}$。未舍入的有理和给出 \eqref{eq:current-E}。两项 $u$ 分别支付有限逼近和系数实现，求积误差已包含在单侧积分端点内。

\subsection{有限改进量递推}
令 $\mathsf Q=2^{80}$，并记
\[
 C_i=[1+i/50,1+(i+1)/50],\qquad 0\le i<200.
\]
本小节全部表值都是分母 $\mathsf Q$ 下的整数分子。四步初值从 $J_0=1/50$ 开始，使用
\[
 J_{n+1}=\frac{20249}{10^6}+\frac{13}{100}J_n.
\]
初始改进量及该递推表示的蕴含由引理~\ref{sv:finite-seed} 给出，并保留其带标记状态及实际截点条件。附属 \texttt{seed\_certificate} 含有其中使用的独立普通项及反馈下端点。因此，算术起点来自对所需族已证明的解析节省量。随后 $J_4=0.023273777353>2327/100000$。初始化
\begin{equation}
 H_i^{(0)}=
 \begin{cases}
 \lfloor \mathsf Q\,2327/100000\rfloor,&1+(i+1)/50\le11/5,\\
 0,&\text{其余情形},
 \end{cases}
 \qquad h_i^{(0)}=0.
 \label{num:table-seed}
\end{equation}

为完整起见，下面指定精确整数转移。一个作用 $\alpha$ 由整数 $b_\alpha$、除数 $d_\alpha\in\{1,5\}$、点对 $(j,m)$、非负整数权 $w_{\alpha j},v_{\alpha j}$ 及目标集合 $T_\alpha$ 构成，其候选值为
\begin{equation}
 b_\alpha+\left\lfloor\frac1{d_\alpha}
 \left\{\sum_{(j,m)}mH_j+
 \sum_j\left\lfloor\frac{w_{\alpha j}h_j}{\mathsf Q}\right\rfloor+
 \sum_j\left\lfloor\frac{v_{\alpha j}H_j}{\mathsf Q}\right\rfloor
 \right\}\right\rfloor.
 \label{num:action-map}
\end{equation}
$70$ 个单步作用满足 $d_\alpha=1$，$36$ 个双步作用满足 $d_\alpha=5$，均在上界常数 $K_*=6.761767>K$ 处计算。各作用的参数、权、常数及目标集合逐一列于 \texttt{actions} 数组。作用参数 $s_\alpha$ 的目标条件为 $1+(i+1)/50\le s_\alpha$。

设 $\ell_j$ 为数据中 $\mathsf Q\log((1+(j+1)/50)/(1+j/50))$ 的整数下界，并令
\[
 S_H(j)=\sum_{k=j}^{199}\left\lfloor\frac{H_k\ell_k}{\mathsf Q}\right\rfloor,
 \qquad
 S_h(j)=\sum_{k=j}^{199}\left\lfloor\frac{h_k\ell_k}{\mathsf Q}\right\rfloor.
\]
对 $50\le i<200$，Buchstab 候选在 $H_i$ 时为 $S_h(\max(50,i-49))$，在 $h_i$ 时为 $S_H(i-49)$。每个新表值取其前值及所有适用候选的最大值。每个右端只使用前一数组。因此这是同时进行的有限递推，不假设收敛到不动点。

解析作用不等式及其后继封闭性通过归纳说明，有理数组在整个闭单元上给出改进量函数的下界。由于反馈权非负，\eqref{num:action-map} 中向下取整保留该方向。对 \eqref{num:table-seed} 施加映射 $430$ 次，即得 \texttt{defect\_tables.430} 中的 $400$ 个值。深度 $248$ 的数组由对应有限迭代从下方控制。深度 $430$ 的数组只用于 $U_3,U_4$ 两处，其余指定位置使用深度 $248$ 的数组。因此，表的变更只在下界泛函实际使用它的位置计入。

\subsection{十一项计数对应的积分}
\label{num:eleven-integrals}
算术项 $Y_j$ 的对应关系固定如下。记 $A_+(s)=sF(s)/(2e^\gamma)$、$A_-(s)=sf(s)/(2e^\gamma)$，并令 $\nu=249/448$、$w_*=25/224$。参考的单前缀项为
\begin{align*}
 L_1&=\frac4{a_{\rm s}}\overline f(\nu/{a_{\rm s}}),\\
 L_2&=4\left\{(1/{h_{\rm s}}+1)\overline f\left(\frac{1+{h_{\rm s}}}{2{h_{\rm s}}}\right)
                 -\overline f(249/50)\right\}+H_{\rm rem},
\end{align*}
其中 $H_{\rm rem}$ 见 \eqref{num:removed-defect}。$L_1$ 的未带标记根节点不代入改进量表，因为其比值 $1/\nu$ 小于 $2$。

对双前缀下界项，令
\[
 \mathcal P_5=\{{a_{\rm s}}<u<t<{h_{\rm s}}\},\qquad
 \mathcal P_6=\{{a_{\rm s}}<u<{h_{\rm s}}<t<{d_{\rm s}}\}.
\]
在增强部分 $t<\mu_0$，令 $V(t,u)$ 为可用候选 $v(t),v(t+u),\Lambda_{\rm cap}-3{a_{\rm s}}/8$ 的最大值；在 $t>\mu_0$ 上取 $V=v(t)$。这里的可用性按附录~\ref{app:coefficients} 保留完整及选定平方条件、选定最大位置及实际素数上限。令 $L=V-t-u$，下界积分为
\begin{equation}
 L_j=4\iint_{\mathcal P_j}
        \frac{A_-(L/{a_{\rm s}})+h_{248,\rm adm}(t,u)}{tuL}\,du\,dt,
 \qquad j=5,6.
 \label{num:pair-lower-integrals}
\end{equation}
仅对满足 $2\le L/{a_{\rm s}}\le5$ 的完全可用状态，改进量项才等于 $h_{248}(L/{a_{\rm s}})$；其余为零。当 $L/{a_{\rm s}}<2$ 时，普通下界项也为零。因此，无法确定的改进量查询可以放弃，不会改变普通下界估计。

对 $m=248,312,430$，单前缀上界积分为
\begin{equation}
 U_j^{(m)}=4\int_{I_j}
       \frac{A_+(s_j(t))-H_m(s_j(t))}{tL_j(t)}\,dt,
 \quad I_3=({a_{\rm s}},1/3),\quad I_4=({a_{\rm s}},{d_{\rm s}}),
 \label{num:one-prefix-upper-integrals}
\end{equation}
其中 $L_j,s_j$ 由下文 \eqref{num:table-change-integral} 给出。十一项和的参考值使用 $m=248$，后续两个表改进仅通过 \eqref{num:successive-table-gains} 插入。

最后两个代价在区域
\[
 \mathcal D_{10}=\{{a_{\rm s}}<w<z<y<x<{h_{\rm s}}\},\qquad
 \mathcal D_{11}=\{{a_{\rm s}}<w<z<y<{h_{\rm s}}<x<{d_{\rm s}}\}.
\]
上采用递减坐标 $x>y>z>w$，积分测度为
\[
 {d_{\rm s}}\sigma_4=\frac{dx\,dy\,dz\,dw}{xyzw}.
\]
令 $S=x+y+z+w$。换位备选的核为 $K_0\omega((1-S)/z)/z$；剩余指数为 $L=\lambda-S$ 的每个可用普通分支提供核 $4\{1-H_{248}(\max(2,L/w))\}/L$，其全部完整和选定条件必须在整个单元上成立。所选单元控制函数从上方控制这些备选的最小值，其积分定义 $U_{10}+U_{11}$ 的上包围。四前缀证书记录单元选择、剩余区间、闭单元改进量最小值以及完整有序覆盖。若普通分支未通过条件，仍可采用换位备选。由于 ${d_{\rm s}}+3{h_{\rm s}}<1$，这两个区域没有单位商贡献。其余换位代价 $U_7,U_8,U_9$ 定义于第~\ref{num:cost-slope} 节。

\subsection{下界修正项的定义}
\label{num:lower-definitions}
先指定修正项，再给出其数值组合。令 $\nu=249/448$、$w_*=25/224$，记 $\overline F=F/(2e^\gamma)$、$\overline f=f/(2e^\gamma)$。${h_{\rm s}}$ 处的普通起点与最小素因子分解中普通部分的差为
\begin{equation}
 B_{\rm ord}=\frac4{h_{\rm s}}\overline f(\nu/{h_{\rm s}})
 -4\left\{(1/{h_{\rm s}}+1)\overline f\left(\frac{1+{h_{\rm s}}}{2{h_{\rm s}}}\right)
                   -\overline f(249/50)\right\}.
 \label{num:ordinary-start-difference}
\end{equation}
从该分解中移去的改进量部分是完整积分
\begin{equation}
 H_{\rm rem}=4\int_{w_*}^{{h_{\rm s}}}
   \frac{H_{248}((1-t)/(2t))}{t(1-t)/2}\,dt.
 \label{num:removed-defect}
\end{equation}
只有去掉这一整个改进量项之后，才代入 ${h_{\rm s}}$ 处的有限普通估计。在改进量单元 $[s_l,s_r]$ 上，\eqref{num:removed-defect} 的贡献区间为 $(w_*,{h_{\rm s}})\cap[1/(2s_r+1),1/(2s_l+1)]$；在乘以因子 $4$ 之前，其原函数为 $2\log(t/(1-t))$。

对 $w={a_{\rm s}},{h_{\rm s}}$，设 $C_3,C_4,C_5$ 为附录~\ref{app:lower} 中同一下界规则的连续上限。上限或者继承，或者从有限指定的子乘积分支中选取，并满足全部完整及选定补集条件。特别地，$\nu\le C_3\le C_4\le C_5\le\Lambda_{\rm cap}$。对递减的 $t_1,\ldots,t_r$，记 $S_r=\sum_{j\le r}t_j$，并令
\[
 {d_{\rm s}}\mu_r=\frac{dt_1\cdots dt_r}{t_1\cdots t_{r-1}t_r^2}.
\]
三个保留区域为
\begin{align*}
 \mathcal E_3(w)&=\{0<t_3<t_2<t_1<w:t_1+3t_2<\nu\},\\
 \mathcal E_4(w)&=\{(t_1,\ldots,t_4):(t_1,t_2,t_3)\in\mathcal E_3(w),
          \ 0<t_4<t_3,\ S_3+3t_4<C_3\},\\
 \mathcal E_5(w)&=\{(t_1,\ldots,t_5):(t_1,\ldots,t_4)\in\mathcal E_4(w),
                                      \ 0<t_5<t_4\}.
\end{align*}
因此，三个非负密度增量为
\begin{align}
 I_3(w)&=4\int_{\mathcal E_3(w)}
  \left\{\overline F\left(\frac{\nu-S_3}{t_3}\right)
        -\overline F\left(\frac{C_3-S_3}{t_3}\right)\right\}\,{d_{\rm s}}\mu_3,
                                      \label{num:third-lower-gain}\\
 I_4(w)&=4\int_{\mathcal E_4(w)}
  \left\{\overline f\left(\frac{C_4-S_4}{t_4}\right)
        -\overline f\left(\frac{C_3-S_4}{t_4}\right)\right\}\,{d_{\rm s}}\mu_4,
                                      \label{num:fourth-lower-gain}\\
 I_5(w)&=4\int_{\mathcal E_5(w)}
  \left\{\overline F\left(\frac{C_4-S_5}{t_5}\right)
        -\overline F\left(\frac{C_5-S_5}{t_5}\right)\right\}\,{d_{\rm s}}\mu_5.
                                      \label{num:fifth-lower-gain}
\end{align}
第五平方条件由 $S_4+2t_5<S_3+3t_4<C_3\le C_4$ 推出。上限选择、分支标签及父立方条件共同固定；采用另一第四上限计算的积分不能追加到 \eqref{num:fifth-lower-gain}。有理证书可以只积分这些非负增量的子集，因此其端点是下界，并不声称选定盒穷尽了理想区域。

对第四节点恢复，设 $\mathcal R_{h_{\rm s}}$ 为 \eqref{sv:restore-gates} 指定的有理斜薄层的不交并。每层上在第四立方检验之前选择上限 $\lambda_4$，满足 $C_3\le S_3+3t_4<\lambda_4$，并保留完整及选定支撑条件。额外增量为
\begin{equation}
 T_{h_{\rm s}}=4\int_{\mathcal R_{h_{\rm s}}}
       \overline f\left(\frac{\lambda_4-S_4}{t_4}\right)\,{d_{\rm s}}\mu_4.
 \label{num:restoration-integral}
\end{equation}
精确薄层体积为 \eqref{sv:slab-volume}。该项以系数 $1/4$ 进入陈氏泛函，且其区域与已经计入的其余增量互不相交。

\subsection{固定上界代价斜率与表的变更}
\label{num:cost-slope}
下文的系数 $A_{\rm cost}$ 是五项代价的一个指定仿射上包围的安全斜率，不必是核最小值的导数。为列出其组成，令 $K_0=6.7619$、$A_+(s)=s\overline F(s)$。三个一维换位代价为
\begin{align*}
 \frac{U_7}{K_0}&=\int_{d_{\rm s}}^{1/3}
  \frac{\log((1-2t)/t)+A_+(1/t-1)-1}{t(1-t)}\,dt,\\
 \frac{U_8}{K_0}&=\int_{a_{\rm s}}^{1/3}\frac{\log(2-3t)}{t(1-t)}\,dt,\\
 \frac{U_9}{K_0}&=\int_{h_{\rm s}}^{d_{\rm s}}\frac{\log((1-t-{d_{\rm s}})/{d_{\rm s}})}{t(1-t)}\,dt.
\end{align*}
首个积分中额外的普通函数贡献无需插值即可计算：
\[
 \int_{d_{\rm s}}^{1/4}\frac{A_+(1/t-1)-1}{t(1-t)}\,dt
 =\int_2^{1/{d_{\rm s}}-2}\frac{\log(u-1)}u
                  \log\frac{1-{d_{\rm s}}}{{d_{\rm s}}(u+1)}\,du.
\]
用 $\underline U_j$ 表示这些关于 $K_0$ 线性代价的有理下端点，它们来自 \texttt{five\_cost\_slope\_certificate} 指定的四阶导数 Simpson 包围。根部代价上端点可采用独立的中点包围；这里上下端点承担不同作用。它们对安全斜率的贡献为 $A_{{\rm cost},789}=(2\underline U_7+\underline U_8+\underline U_9)/(4K_0)$。对 $40000$ 优先父单元之外的保留四前缀单元 $\mathscr B$，用 $A_{\mathscr B}$ 表示在 $K_0$ 处分配的上界质量，用 $s_{\mathscr B}$ 表示所附非负仿射斜率。所需逐单元条件为
\begin{equation}
 0\le s_{\mathscr B}\le A_{\mathscr B}/K_0,\qquad A_{\mathscr B}\ge K_0\mathcal T_{\mathscr B},
 \label{num:affine-majorant-slope}
\end{equation}
其中 $\mathcal T_{\mathscr B}$ 是精确的非负换位质量。由此得到
\[
 A_{\mathscr B}-(K_0-K)s_{\mathscr B}\ge(K/K_0)A_{\mathscr B}\ge K\mathcal T_{\mathscr B}.
\]
这里不需要不等式 $s_{\mathscr B}\le\mathcal T_{\mathscr B}$。在分母 $\mathsf Q=2^{80}$ 下，保留行的斜率为 $s_{\mathscr B}=10000d_{\mathscr B}/\mathsf Q$；对 $U_{10},U_{11}$，其和分别为 $1277985388931956822$、$7976183963523684871$。因此
\begin{equation}
 A_{\rm cost}=A_{{\rm cost},789}+
 \frac{10000(1277985388931956822+7976183963523684871)}{4\mathsf Q}.
 \label{num:five-cost-slope}
\end{equation}
完整保留单元不等式由证书对应表中的 \texttt{feedback\_original\_slope} 条目索引。

对两个有效改进量包络 $H_m\le H_n$，定义
\begin{equation}
 \mathcal T(H_m,H_n)=
 \int_{a_{\rm s}}^{1/3}\frac{(H_n-H_m)(s_3(t))}{tL_3(t)}\,dt
 +\int_{a_{\rm s}}^{d_{\rm s}}\frac{(H_n-H_m)(s_4(t))}{tL_4(t)}\,dt,
 \label{num:table-change-integral}
\end{equation}
其中 $s_j(t)=\max(2,L_j(t)/{a_{\rm s}})$，在 $[1,5]$ 外贡献为零；对 $e=57/224$，定义
\[
 L_3(t)=\begin{cases}
 (1-t)/2,&{a_{\rm s}}<t<\mu_0,\\
 57/64-5t/2,&\mu_0<t<e,\\
 1/2-t,&e<t<1/3.
 \end{cases}
\]
函数 $L_4$ 使用限制在 $({a_{\rm s}},{d_{\rm s}})$ 的前两个分支。两个上界计数中的因子 $4$ 抵消单个陈氏因子 $1/4$，故
\begin{equation}
 T_{312}=\mathcal T(H_{248},H_{312}),\qquad
 T_{430}=\mathcal T(H_{312},H_{430}).
 \label{num:successive-table-gains}
\end{equation}
节点处查询取相邻较小值；对正测度分支 $s_j=2$，这意味着取单元 $49$、$50$ 的最小值。在仿射段 $L=A_0-B_0t$ 上，精确的逆单元原函数为
\[
 \int_l^r\frac{dt}{t(A_0-B_0t)}
 =\frac1{A_0}\log\frac{r(A_0-B_0l)}{l(A_0-B_0r)}.
\]
由此确定每个表变更区间及其有理对数包围，无需跨节点作数值插值。

\subsection{数学对象与证书字段}
\label{num:certificate-chain}
表~\ref{num:certificate-table} 关联计数估计、积分界及附属证书。最后一列的路径为 \filepath{certificate-map.json} 中的目标；每个目标包含来源文件名、精确 JSON 指针、方向以及所需有理提取公式。对应表还包含每类原始对象的包围与覆盖断言，以及完整见证文件名。\texttt{domain\_constructions} 条目标识实际的 $C_3,C_4,C_5$ 区域生成器及其选定记录，包括两个三阶起点、四阶弦线及斜向区域、五阶子单元增量和斜层恢复。对 $U_{10},U_{11}$，它们标识生成器 \filepath{witnesses/generator-0041.py} 及完整有序见证 \filepath{witnesses/record-0638.json}。见证记录指定这些区域的父规则、方法选择、第 $248$ 代轮廓及完整面描述。

\begingroup
\small
\setlength{\tabcolsep}{4pt}
\begin{longtable}{@{}>{\raggedright\arraybackslash}p{.15\textwidth}>{\raggedright\arraybackslash}p{.43\textwidth}>{\raggedright\arraybackslash}p{.35\textwidth}@{}}
\caption{定义、方向与证书目标。}
\label{num:certificate-table}\\
\toprule 对象 & 数学估计 & 证书目标\\
\midrule
\endfirsthead
\toprule 对象 & 数学估计 & 证书目标\\
\midrule
\endhead
$G_j$, $1\le j\le22$ &
上界树中指定的子节点与父节点密度差；非负节省量可只采用子区域。 &
\filepath{/upper_saving_bounds/}\,$(j-1)$；下端点。\\[3pt]
有限作用 &
\eqref{sv:one-step-operator}, \eqref{sv:two-step-base},
\eqref{sv:two-step-feedback}；完整 $\Gamma_5$，并对 $D_5^-$ 单独支付代价。 &
\filepath{/actions/}\,$k$；基础项下界与反馈权；
\filepath{/gamma5_cost_decompositions}.\\[3pt]
初值与表 &
引理~\ref{sv:finite-seed} 及同时映射
\eqref{num:action-map}；闭单元估计。 &
\filepath{/seed_certificate}, \filepath{/defect_tables/}\,$m$;
下端点，$m=248,312,430$。\\[3pt]
$L_i,U_i$ &
第~\ref{num:eleven-integrals} 节的积分；每个负项代价覆盖完整保留区域，每个正项采用下界。 &
\filepath{/lower_root_bounds}；$L_i$ 取下界，$U_i$ 取上界。\\[3pt]
$B_{\rm ord},H_{\rm rem}$ &
\eqref{num:ordinary-start-difference}, \eqref{num:removed-defect};
去掉完整的第 $248$ 代改进量项。 &
\filepath{/lower_auxiliary_bounds}；$B_{\rm ord}$ 取下界，$H_{\rm rem}$ 取上界。\\[3pt]
$I_r(w),T_{h_{\rm s}}$ &
\eqref{num:third-lower-gain}--\eqref{num:restoration-integral}；共同的下界规则及不交恢复薄层。 &
\filepath{/lower_integral_bounds} 及
\filepath{/lower_auxiliary_bounds/restoration_L2_lower}；下界。\\[3pt]
$A_{\rm cost}$, $T_{312},T_{430}$ &
\eqref{num:affine-majorant-slope}--\eqref{num:successive-table-gains};
上界代价斜率及连续表变更。 &
\filepath{/five_cost_slope_certificate},
\filepath{/lower_auxiliary_bounds}；安全的下界系数。\\[3pt]
配对删除 &
$\mathscr E_0\cup \mathscr E_1$ 上分配的非负上界及被删下界估计量；\eqref{num:paired-extra}、\eqref{num:saved-loss}。 &
\filepath{/paired_rectangles}；精确几何、质量区间及下界估计量整数。\\[3pt]
$E_{\rm up},R_{\rm low}$ 及预留量 &
\eqref{num:lower-sum} 及上界密度表达式中的唯一组合。 &
\filepath{/scalar_bounds}；精确和与严格剩余量。\\
\bottomrule
\end{longtable}
\endgroup

对每个负代价，映射中的记录指定其保留区域的完整覆盖；对正节省量，则指定计算其下界所用的非负子族。

\subsection{下界泛函}
对带符号十一项表达式，每个正项使用下端点，每个负项使用上端点。下列整数除以 $10^{14}$ 后给出这样的端点：
\begin{center}
\begin{tabular}{@{}l r r@{}}
\toprule 项 & 除以 $4$ 前的系数 & 端点分子\\
\midrule
$L_1$ & $3$ & $1346406373358053$\\
$L_2$ & $1$ & $710262601191162$\\
$L_5$ & $1$ & $237450056438739$\\
$L_6$ & $1$ & $301984961684469$\\
$U_3$ & $-1$ & $2031242249325404$\\
$U_4$ & $-1$ & $1405868301028074$\\
$U_7$ & $-2$ & $138343028785273$\\
$U_8$ & $-1$ & $405707594761181$\\
$U_9$ & $-1$ & $325703365221150$\\
$U_{10}+U_{11}$ & $-1$ & $67568703438104$\\
\bottomrule
\end{tabular}
\end{center}
定义 $C_{11}$ 时采用精确有理端点，而非这些缩短值。四前缀最小值覆盖全部
\[
 \binom{51}{4}+48\binom{50}{3}=1190700
\]
有序行，包括坐标相等的单形。优先细分以其完整 $516779$ 个子单元之和与原父估计的最大值，替换 $40000$ 个父单元。额外系数节省量为
\[
 P=\frac{106520738437439293219}{604462909807314587353088}.
\]

项 $B_{\rm ord},H_{\rm rem}$, $I_r(a_{\rm s}),I_r(h_{\rm s})$, $T_{h_{\rm s}}$ 定义于 \eqref{num:ordinary-start-difference}--\eqref{num:restoration-integral}。其精确界由 \filepath{lower_auxiliary_bounds} 与 \filepath{lower_integral_bounds} 指定。在 $K_0=6.7619$ 处，令
\[
 C_{K_0}=C_{11}+P+\frac14\left(
 B_{\rm ord}-H_{\rm rem}
 +\sum_{r=3}^5\{3I_r({a_{\rm s}})+I_r({h_{\rm s}})\}+T_{h_{\rm s}}\right).
\]
三阶积分覆盖两个截点区域；四阶积分保留全部指定掩膜及不交附加区带；五阶积分在 ${h_{\rm s}}$ 一侧使用 $158251$ 个下界估计量和 $23662$ 个附加薄层或相等坐标单元，在 ${a_{\rm s}}$ 一侧使用 $679$ 个积分。斜层恢复使用 $98668$ 个盒和 $789344$ 层。只有非负增量可以略去边界条带。非对称核在其有序单形上估计，不对它假设阶乘对称性。

取 \eqref{num:five-cost-slope} 中的 $A_{\rm cost}$ 及 \eqref{num:successive-table-gains} 中的 $T_{312},T_{430}$，初始下界端点至少为
\begin{equation}
 C_{K_0}+A_{\rm cost}(K_0-K)+T_{312}+T_{430}
 +0.002095161+0.003799514-0.002372094.
 \label{num:lower-sum}
\end{equation}
直接有理代入所得值不小于 \eqref{eq:base-R} 中的 $R_{\rm base}$。下文配对单元细化恰增加 \eqref{eq:current-refinement-delta} 中的 $\Delta$。扣除 \eqref{eq:current-B} 的两项预留量后，所得下界端点超过 $c$ 的量为
\[
 \frac{13457260926863633771775678584711245753}
 {2742952509978644604793778369418231808000000000}>0.
\]
结合上界计算，这就从所述有限包围证明了引理~\ref{lem:numerical-endpoints} 的标量不等式。

\section{配对单元的细化}\label{app:refinement}

配对矩形按以下方式精确指定，无需几何舍入约定。令
\[
 t_i=\mu_0+i\frac{{d_{\rm s}}-\mu_0}{512},\qquad
 u_j={a_{\rm s}}+j\frac{{h_{\rm s}}-{a_{\rm s}}}{512}\qquad(0\le i,j\le512).
\]
\texttt{paired\_rectangles} 的每行指定族 $\mathscr E_0$ 或 $\mathscr E_1$，以及矩形 $[t_i,t_{i+1}]\times[u_j,u_{j+1}]$。两族分别有 $6091$、$11481$ 个单元，没有公共单元，边界积分测度为零。对矩形 $\mathscr B=[t_l,t_r]\times[u_l,u_r]$，定义
\[
 Z_{\mathscr B}=\min\left(u_l,\frac{3(1-t_r-u_r)}{13}\right).
\]
每行有理数据均满足 $Z_{\mathscr B}>{a_{\rm s}}$、$1-t_r-u_r-2Z_{\mathscr B}\ge(7/3)Z_{\mathscr B}$，故
\[
 \frac{1-t-u-z-w}{z}\ge\frac73
 \qquad({a_{\rm s}}<w<z<Z_{\mathscr B},\ (t,u)\in \mathscr B).
\]

在每个矩形上，引理~\ref{lem:current-omega-refinement} 给出整个纤维上的 $\omega\ge1/2$，以及指定内区域上的 $\omega\ge11/20$。第二个估计改善第一个贡献的一部分，并以增量计入。

因此，非负父上界至少包含整条纤维上的 $1/2$ 贡献，以及内纤维上的 $1/20$ 贡献。这是同时成立的两个非负贡献。积分并计入单个陈氏因子 $1/4$ 后，额外贡献为
\begin{equation}
 G_{\mathscr B}=\frac K{80}\log\frac{t_r}{t_l}\log\frac{u_r}{u_l}
 \left\{\frac1{a_{\rm s}}-\frac{1+\log(Z_{\mathscr B}/{a_{\rm s}})}{Z_{\mathscr B}}\right\}.
 \label{num:paired-extra}
\end{equation}
事实上，$w$ 积分为 $\log(z/{a_{\rm s}})$，且
\[
 \int_{a_{\rm s}}^{Z_{\mathscr B}}\frac{\log(z/{a_{\rm s}})}{z^2}\,dz
 =\frac1{a_{\rm s}}-\frac{1+\log(Z_{\mathscr B}/{a_{\rm s}})}{Z_{\mathscr B}}.
\]
此处的归一化与~\eqref{eq:current-extra-fibre} 一致。

数据每行都包含 \eqref{num:paired-extra} 的有理上下端点，以分母 $2^{160}$ 下的整数分子表示。对其中三个对数应用 \eqref{num:log-enclosure}，再将下端点相加，得
\begin{equation}
 \sum_{\mathscr B\in \mathscr E_0\cup \mathscr E_1}G_{\mathscr B}>\frac{558583}{10^9}.
 \label{num:gross-bound}
\end{equation}
两族之和分别约为 $1.9318091250\cdot10^{-4}$、$3.6540308174\cdot10^{-4}$；实际只使用有理下界 \eqref{num:gross-bound}。

\section{被删除的下界估计量与最终预留量}
\label{app:refinement-review}

配对删除在移去正项的同时移去其已分配的下界估计量。为控制由此产生的损失，下面同时给出该估计量的精确有限和及直接有理上包围。

在一个单元上，记
\[
 R(t,u)=\frac{57}{64}-\frac52t-u,\qquad
 r=R(t_r,u_r),\quad R_0=R(t_l,u_l),\quad
 M=\log(t_r/t_l)\log(u_r/u_l).
\]
$\mathscr E_0$ 的单元满足 $R_0/{a_{\rm s}}\le2$，已分配下界估计量为零。$\mathscr E_1$ 中有 $227$ 个单元满足 $r/{a_{\rm s}}<2$，其估计量也为零。其余 $11254$ 个单元满足 $2\le r/{a_{\rm s}}\le R_0/{a_{\rm s}}\le9/4$。令
\[
 f_l=\frac{\log(r/{a_{\rm s}}-1)}{r/{a_{\rm s}}},\qquad
 f_h=\frac{\log(R_0/{a_{\rm s}}-1)}{R_0/{a_{\rm s}}},
\]
及
\[
 M_R=\frac{57}{64}M
 -\frac52(t_r-t_l)\log(u_r/u_l)
 -(u_r-u_l)\log(t_r/t_l).
\]
由于 $r\le R(t,u)\le R_0$，按陈氏权四分之一归一化，两个已分配普通估计量至多为
\begin{equation}
 \frac{f_lM}{{a_{\rm s}}},\qquad
 \frac{f_l(R_0M-M_R)+f_h(M_R-rM)}{a_{\rm s}(R_0-r)}.
 \label{num:ordinary-proxy}
\end{equation}
$f_l,f_h$ 的系数均非负，因此将这些函数的下端点替换为精确值只会增大估计量。\eqref{num:ordinary-proxy} 中两个表达式的最大值，仍控制两个已分配备选的最大值。

设 $h_{\mathscr B}^{\min}$ 为与 $[r/{a_{\rm s}},R_0/{a_{\rm s}}]$ 相交的全部闭单元上第 $248$ 代下界表值的最小值，节点处包括两侧单元。精确支撑包含于已分配估计量向外取整的查询支撑中，所以 $h_{\mathscr B}^{\min}$ 不小于后者采用的最小值。因此，附加项 $h_{\mathscr B}^{\min}M/R_0$ 从上方控制已分配改进量下界。该比较涉及有限下界估计量的最小值，不是对数学函数 $h$ 给上界。将它加到 \eqref{num:ordinary-proxy} 的最大值上，即得所需逐单元控制量。

数据以分母 $\mathsf Q=2^{80}$ 下的整数记录各单元的下界估计量。

这些整数之和为 $11290806845175652020270$。除以 $4\mathsf Q$ 后，
\begin{equation}
 L_{\mathscr E}
 =\frac{5645403422587826010135}
 {2417851639229258349412352}
 <\lambda_{\rm loss}:=\frac{466977}{200000000}.
 \label{num:saved-loss}
\end{equation}
对上述控制量作向外计算，也得到低于同一 $\lambda_{\rm loss}$ 的上端点。计算时可将对数自变量缩减到 $[1,3/2]$，采用 $\log(1+x)$ 的 $320$ 项交错级数，并以下一项控制余项。全部除式的分母区间下端点均为正。所得控制量总和被包围在 $0.0023348841363598131014744273487904$ 附近。

删除之前，完整 $L_6$ 的保留下端点为 $3650774173156900914101685/\mathsf Q$；其非增强部分采用刚才的单元分解。因此，从该端点减去 $4\lambda_{\rm loss}$ 即可从下方控制保留的 $L_6$，增强部分不受影响。相对于 \eqref{num:lower-sum}，支付代价的减少恰为
\[
 0.002372094-\lambda_{\rm loss}=\frac{37209}{10^9}.
\]
结合 \eqref{num:gross-bound}，得到原始增量
\[
 \Delta=\frac{558583+37209}{10^9}
       =\frac{37237}{62500000}.
\]
给有限逼近和系数实现各预留 $\Delta/8$，则保留的净增量为 $3\Delta/4=0.000446844$。全部估计均针对同一组固定矩形，每个密度贡献只出现一次。它们给出 $R_{\rm low}=R_{\rm base}+\Delta$、$B_{\rm low}=B_{\rm base}+\Delta/8$，与 \eqref{eq:current-R}、\eqref{eq:current-B} 一致。

\section*{AI 辅助声明}
OpenAI ChatGPT/Codex 协助了数学探索、证明发展与评估、参数搜索、计算核验及文稿编辑。
作者对数学内容及最终稿件负责。

\begin{thebibliography}{99}
\bibitem{ABL}
E. Assing, V. Blomer and J. Li,
\emph{Uniform Titchmarsh divisor problems},
\href{https://arxiv.org/abs/2005.13915v1}{arXiv:2005.13915v1} (2020).

\bibitem{ChanSupplement}
S. Chan,
\emph{Theorem $(1+1.75)$ on the Goldbach Conjecture --- Numerical Supplement},
version 20261106, Zenodo.
\href{https://doi.org/10.5281/zenodo.23755453}{doi:10.5281/zenodo.23755453}.

\bibitem{Drappeau}
S. Drappeau,
\emph{Sums of Kloosterman sums in arithmetic progressions, and the error
term in the dispersion method}, Proc. Lond. Math. Soc. (3)
\textbf{114} (2017), no.~4, 684--732.
\href{https://doi.org/10.1112/plms.12022}{doi:10.1112/plms.12022}.

\bibitem{Kedlaya}
K. S. Kedlaya,
\emph{Analytic Number Theory}, Chapter~18, Theorem~18.2,
\url{https://kskedlaya.org/ant/chapter-18.html}.

\bibitem{KimSarnak}
H. Kim and P. Sarnak,
\emph{Refined estimates towards the Ramanujan and Selberg conjectures},
J. Amer. Math. Soc. \textbf{16} (2003), 175--181,
Appendix~2 to H. Kim, \emph{Functoriality for the exterior square of
$\mathrm{GL}_4$ and the symmetric fourth of $\mathrm{GL}_2$}.
\href{https://doi.org/10.1090/S0894-0347-02-00410-1}{doi:10.1090/S0894-0347-02-00410-1}.

\bibitem{LiLiu}
J. Li and J. Liu,
\emph{Theorem $(1+1.9)$ on the Goldbach Conjecture},
\href{https://arxiv.org/abs/2606.05224v2}{arXiv:2606.05224v2} (2026).

\bibitem{RunboLi}
R. Li,
\emph{On Chen's theorem, Goldbach's conjecture and almost prime twins II},
Math. Rep. (Bucur.) \textbf{28(78)} (2026), nos.~1--2, 39--61.
\href{https://doi.org/10.59277/mrar.2026.28.78.1.2.39}{doi:10.59277/mrar.2026.28.78.1.2.39};
\href{https://arxiv.org/abs/2405.05727v4}{arXiv:2405.05727v4}.

\bibitem{Lichtman}
J. D. Lichtman,
\emph{Primes in arithmetic progressions to large moduli, and Goldbach
beyond the square-root barrier},
\href{https://arxiv.org/abs/2309.08522v1}{arXiv:2309.08522v1} (2023).

\bibitem{MaynardI}
J. Maynard,
\emph{Primes in arithmetic progressions to large moduli I: Fixed residue
classes},
\href{https://arxiv.org/abs/2006.06572v2}{arXiv:2006.06572v2} (2021).

\bibitem{MaynardII}
J. Maynard,
\emph{Primes in arithmetic progressions to large moduli II: Well-factorable
estimates},
\href{https://arxiv.org/abs/2006.07088v1}{arXiv:2006.07088v1} (2020).

\bibitem{Pascadi}
A. Pascadi,
\emph{Large sieve inequalities for exceptional Maass forms and the
greatest prime factor of $n^2+1$}, Forum Math. Pi \textbf{14} (2026), e8.
\href{https://doi.org/10.1017/fmp.2026.10025}{doi:10.1017/fmp.2026.10025};
\href{https://arxiv.org/abs/2404.04239v3}{arXiv:2404.04239v3}.

\bibitem{Shiu}
P. Shiu,
\emph{A Brun--Titchmarsh theorem for multiplicative functions},
J. Reine Angew. Math. \textbf{313} (1980), 161--170.
\href{https://doi.org/10.1515/crll.1980.313.161}{doi:10.1515/crll.1980.313.161}.

\bibitem{Suzuki}
Y. Suzuki,
\emph{The Rosser--Iwaniec sieve}, unpublished notes, 14 March 2023,
102 pp., Proposition~11.8 and Lemma~14.7.
\url{https://www.suzuyu1729.com/notes/NagoyaSeminar2022/RosserIwaniec.pdf}.

\bibitem{TaoSW}
T. Tao,
\emph{254A, Notes 2: Complex-analytic multiplicative number theory},
9 December 2014, Exercise~64,
\url{https://terrytao.wordpress.com/2014/12/09/254a-notes-2-complex-analytic-multiplicative-number-theory/}.

\bibitem{Wu2004}
J. Wu,
\emph{Chen's double sieve, Goldbach's conjecture and the twin prime
problem}, Acta Arith. \textbf{114} (2004), 215--273.
\href{https://doi.org/10.4064/aa114-3-2}{doi:10.4064/aa114-3-2}.

\bibitem{Wu}
J. Wu,
\emph{Chen's double sieve, Goldbach's conjecture and the twin prime
problem, 2}, Acta Arith. \textbf{131} (2008), 367--387.
\href{https://arxiv.org/abs/0709.3764v1}{arXiv:0709.3764v1}.

\bibitem{Young}
M. P. Young,
\emph{Explicit calculations with Eisenstein series},
\href{https://arxiv.org/abs/1710.03624v2}{arXiv:1710.03624v2} (2017).
\end{thebibliography}

\end{document}
